\chapter{Flat base change for the pushforward of a quasi-coherent sheaf ($i=0$)}
\label{chap:Cohomology_FlatBaseChange}
% archon:covers AlgebraicJacobian/Cohomology/FlatBaseChange.lean AlgebraicJacobian/Cohomology/FlatBaseChangeGlobal.lean
% SOURCE: [Stacks Project], Cohomology of Schemes, Section "Cohomology and
% base change, I", and Tag 02KH (read from references/stacks-coherent.tex and
% references/stacks-coherent.md).

\section{Setup and motivation}

This chapter establishes the $i=0$ (i.e.\ \(H^0\) / direct-image) case of flat
base change: the formation of the pushforward \(f_*\mathcal{F}\) of a
quasi-coherent sheaf commutes with flat base change on the target. It is the
foundational, fully self-contained slice of the cohomology-and-base-change
package that ultimately underlies relative representability statements; the
higher derived cases \(R^i f_*\) for \(i \geq 1\) are treated separately and are
out of scope here.

Throughout, \(f : X \to S\) is a morphism of schemes that is quasi-compact and
quasi-separated, and \(\mathcal{F}\) is a quasi-coherent \(\mathcal{O}_X\)-module.
A base change of \(f\) along an arbitrary morphism \(g : S' \to S\) is recorded
by the cartesian square obtained from the fibre product
\(X' = X \times_S S'\):
\[
  \begin{array}{ccc}
    X' & \xrightarrow{\;g'\;} & X \\[2pt]
    {\scriptstyle f'}\big\downarrow & & \big\downarrow{\scriptstyle f} \\[2pt]
    S' & \xrightarrow{\;g\;} & S
  \end{array}
\]
Here \(f' : X' \to S'\) is the base change of \(f\), and \(g' : X' \to X\) is the
first projection. We write \(\mathcal{F}' = (g')^*\mathcal{F}\) for the pullback
of \(\mathcal{F}\) to \(X'\).

\section{The canonical base-change map}

\begin{definition}\leanok
[Base-change map for the pushforward]
  \label{def:pushforward_base_change_map}
  \lean{AlgebraicGeometry.pushforwardBaseChangeMap}
  % SOURCE: [Stacks Project], Cohomology of Schemes, Section "Cohomology and
  % base change, I", base-change diagram (\ref{equation-base-change-diagram})
  % (read from references/stacks-coherent.tex, L877--904).
  % SOURCE QUOTE: "Let $f : X \to S$ be a morphism of schemes.
  % Let $\mathcal{F}$ be a quasi-coherent sheaf on $X$.
  % Suppose further that $g : S' \to S$ is any morphism of schemes. Denote
  % $X' = X_{S'} = S' \times_S X$ the base change of $X$ and denote
  % $f' : X' \to S'$ the base change of $f$.
  % Also write $g' : X' \to X$ the projection,
  % and set $\mathcal{F}' = (g')^*\mathcal{F}$."
  \textit{Source: Stacks Project, Cohomology of Schemes, \S\,Cohomology and base
  change, I.}
  Let \(f : X \to S\) be a morphism of schemes, \(\mathcal{F}\) a quasi-coherent
  \(\mathcal{O}_X\)-module, and \(g : S' \to S\) an arbitrary morphism, giving the
  cartesian square above with projection \(g' : X' \to X\) and
  \(\mathcal{F}' = (g')^*\mathcal{F}\). The \emph{base-change map for the
  pushforward} is the canonical morphism of \(\mathcal{O}_{S'}\)-modules
  \[
    g^*\bigl(f_*\mathcal{F}\bigr) \longrightarrow f'_*\bigl((g')^*\mathcal{F}\bigr)
    = f'_*\mathcal{F}' .
  \]
  It is adjoint to the morphism
  \(f_*\mathcal{F} \to g_*\, f'_*\mathcal{F}' = f_*\,(g')_*\mathcal{F}'\) obtained
  by applying \(f_*\) to the unit
  \(\mathcal{F} \to (g')_*(g')^*\mathcal{F}\) of the
  \(((g')^*, (g')_*)\)-adjunction, using the commutativity
  \(g \circ f' = f \circ g'\) of the square.
\end{definition}

\section{The affine case}

\subsection{Locality of isomorphisms for quasi-coherent sheaf morphisms}

Before treating the affine computation we record three supplementary locality
lemmas not available in Mathlib. Mathlib provides the per-open criterion
(\(\mathrm{IsIso}\,\varphi \iff \forall U,\ \mathrm{IsIso}(\varphi_U)\)) and a
stalkwise isomorphism criterion for \(\mathrm{TopCat.Sheaf}\)-valued morphisms, but
it does not package the stalk-local or basis-local criterion at the level of
morphisms of \(\mathcal{O}_X\)-modules. The following lemmas bridge that gap; they
are exactly the locality tools used in the first reduction of the affine
base-change lemma below, where one checks the base-change map after restricting to
the affine opens of the base. They carry no external citation.

\begin{lemma}\leanok
[Stalk-local criterion for isomorphisms of \(\mathcal{O}_X\)-modules]
  \label{lem:modules_isIso_iff_stalk}
  \lean{AlgebraicGeometry.Modules.isIso_iff_isIso_stalkFunctor_map}
  Let \(X\) be a scheme and \(\varphi : M \to N\) a morphism of sheaves of
  \(\mathcal{O}_X\)-modules. Then \(\varphi\) is an isomorphism if and only if the
  underlying morphism of abelian-group presheaves induces an isomorphism on the
  stalk at every point \(x \in X\).
\end{lemma}

\begin{proof}



  If \(\varphi\) is an isomorphism, then so is its image under the forgetful
  functor to abelian-group presheaves, and the stalk functor at each \(x\),
  being a functor, carries isomorphisms to isomorphisms. Conversely, assume the
  stalk map is an isomorphism at every \(x\). Package the underlying abelian-group
  presheaves of \(M\) and \(N\) as sheaves of abelian groups, and \(\varphi\) as a
  morphism of these sheaves; the stalkwise-isomorphism hypothesis lets us apply the
  \(\mathrm{TopCat.Sheaf}\)-level stalk-local isomorphism criterion to conclude that
  this morphism of sheaves of abelian groups is an isomorphism. Pushing it through
  the forgetful functor to presheaves yields an isomorphism of the underlying
  abelian-group presheaf morphisms, and since the forgetful functor from
  \(\mathcal{O}_X\)-modules to abelian-group presheaves reflects isomorphisms,
  \(\varphi\) itself is an isomorphism.
\end{proof}

\begin{lemma}\leanok
[Basis-local criterion for isomorphisms of \(\mathcal{O}_X\)-modules]
  \label{lem:modules_isIso_of_isBasis}
  \lean{AlgebraicGeometry.Modules.isIso_of_isIso_app_of_isBasis}
  \uses{lem:modules_isIso_iff_stalk}
  Let \(X\) be a scheme, let \(\{B_i\}_{i \in \iota}\) be a family of opens whose
  range is a basis of the topology of \(X\), and let \(\varphi : M \to N\) be a
  morphism of sheaves of \(\mathcal{O}_X\)-modules. If \(\varphi\) restricts to an
  isomorphism on the sections over every basic open \(B_i\), then \(\varphi\) is an
  isomorphism.
\end{lemma}

\begin{proof}
  \uses{lem:modules_isIso_iff_stalk}
  By Lemma~\ref{lem:modules_isIso_iff_stalk} it suffices to prove that the induced
  stalk map is bijective at each point \(x \in X\). For injectivity, a germ that
  vanishes at \(x\) already vanishes on some open neighbourhood, which may be
  shrunk to a basic open \(B_i\); on \(B_i\) the section map \(\varphi_{B_i}\) is
  injective, so the original section is zero, giving injectivity of the stalk map
  from injectivity over the basis. For surjectivity, an arbitrary germ at \(x\) is
  represented by a section of \(N\) over some basic open \(B_i\) containing \(x\);
  since \(\varphi_{B_i}\) is surjective, this section lifts to a section of \(M\)
  over \(B_i\), whose germ at \(x\) maps to the given germ. Thus the stalk map is
  bijective, and \(\varphi\) is an isomorphism.
\end{proof}

\begin{lemma}\leanok
[Affine-open locality criterion for isomorphisms of \(\mathcal{O}_X\)-modules]
  \label{lem:modules_isIso_iff_affineOpens}
  \lean{AlgebraicGeometry.Modules.isIso_iff_isIso_app_affineOpens}
  \uses{lem:modules_isIso_of_isBasis}
  Let \(X\) be a scheme and \(\varphi : M \to N\) a morphism of sheaves of
  \(\mathcal{O}_X\)-modules. Then \(\varphi\) is an isomorphism if and only if it
  restricts to an isomorphism on the sections over every affine open of \(X\).
\end{lemma}

\begin{proof}
  \uses{lem:modules_isIso_of_isBasis}
  The forward implication is immediate, since the section functor over a fixed open
  preserves isomorphisms. For the converse, the affine opens of a scheme form a
  basis of its topology, so the claim is the special case of
  Lemma~\ref{lem:modules_isIso_of_isBasis} in which the basis is the family of
  affine opens.
\end{proof}

\subsection{The affine computation}

The whole result rests on the following elementary computation in the affine
setting, where the base-change map is visibly an isomorphism because tensoring is
associative. Before stating it we isolate the one ingredient not yet available as
a ready-made result on which the affine reduction turns: the description of the
pushforward of a tilde-module along a morphism of affine schemes.

We first record three supplementary lemmas computing the global sections of an
affine pushforward; they carry no external citation.

\begin{lemma}



\leanok
[Naturality of the global-sections comparison square]
  \label{lem:globalSectionsIso_hom_comp_specMap_appTop}
  \lean{AlgebraicGeometry.globalSectionsIso_hom_comp_specMap_appTop}
  Let \(\varphi : R \to R'\) be a homomorphism of commutative rings, and write
  \(\mathrm{gs}_R : \Gamma(\operatorname{Spec} R, \top) \cong R\) and
  \(\mathrm{gs}_{R'} : \Gamma(\operatorname{Spec} R', \top) \cong R'\) for the
  global-sections isomorphisms of the structure sheaves. Then the ring square
  \[
    \mathrm{gs}_R \;\circ\; (\operatorname{Spec}\varphi)^{\sharp}_{\top}
    \;=\; \varphi \;\circ\; \mathrm{gs}_{R'}
  \]
  commutes, where \((\operatorname{Spec}\varphi)^{\sharp}_{\top}\) is the
  top-open component of the structure-sheaf comparison map underlying
  \(\operatorname{Spec}\varphi\).
\end{lemma}

\begin{proof}



  This is the naturality of the unit/counit of the
  \(\Gamma \dashv \operatorname{Spec}\) adjunction read off at the top open: the
  global-sections map of \(\operatorname{Spec}\varphi\) is, under the
  global-sections isomorphisms on each side, conjugate to \(\varphi\) itself.
  Transposing the naturality square of the \(\operatorname{Spec}\) unit gives
  exactly the asserted ring equation.
\end{proof}

\begin{lemma}



\leanok
[Global sections of an affine pushforward]
  \label{lem:gammaPushforwardIso}
  \lean{AlgebraicGeometry.gammaPushforwardIso}
  \uses{lem:globalSectionsIso_hom_comp_specMap_appTop}
  Let \(\varphi : R \to R'\) be a homomorphism of commutative rings and let
  \(N\) be an \(\operatorname{Spec}(R')\)-module. Then there is a canonical
  isomorphism of \(R\)-modules
  \[
    \Gamma\bigl((\operatorname{Spec}\varphi)_* N\bigr)
    \;\cong\;
    \operatorname{restr}_\varphi\bigl(\Gamma N\bigr),
  \]
  identifying the \(R\)-module of global sections of the pushforward with the
  restriction of scalars along \(\varphi\) of the \(R'\)-module of global sections
  of \(N\).
\end{lemma}

\begin{proof}
  \uses{lem:globalSectionsIso_hom_comp_specMap_appTop}
  Because \((\operatorname{Spec}\varphi)^{-1}(\top) = \top\), the global sections of
  the pushforward \((\operatorname{Spec}\varphi)_* N\) are, on underlying abelian
  groups, the global sections of \(N\) with no transport along the open-set
  identification. The two sides therefore carry the same underlying abelian group;
  both unwind to nested restriction-of-scalars towers along the structure-sheaf
  global-sections maps. These towers are reconciled by applying the
  restriction-of-scalars composition identity twice and substituting the ring
  equation of
  Lemma~\ref{lem:globalSectionsIso_hom_comp_specMap_appTop}, which shows the
  resulting \(R\)-action is exactly restriction of scalars along \(\varphi\).
\end{proof}

\begin{lemma}



\leanok
[Global sections of an affine pushforward of a tilde-module]
  \label{lem:gammaPushforwardTildeIso}
  \lean{AlgebraicGeometry.gammaPushforwardTildeIso}
  \uses{lem:gammaPushforwardIso}
  Let \(\varphi : R \to R'\) be a homomorphism of commutative rings and let \(M\)
  be an \(R'\)-module. Then there is a canonical isomorphism of \(R\)-modules
  \[
    \Gamma\bigl((\operatorname{Spec}\varphi)_* \widetilde{M}\bigr)
    \;\cong\;
    \operatorname{restr}_\varphi M .
  \]
\end{lemma}

\begin{proof}
  \uses{lem:gammaPushforwardIso}
  Specialise Lemma~\ref{lem:gammaPushforwardIso} to \(N = \widetilde{M}\) and
  compose with the tilde-unit isomorphism
  \(\Gamma(\widetilde{M}) \cong M\) of \(R'\)-modules (which restricts to an
  isomorphism of \(R\)-modules after applying \(\operatorname{restr}_\varphi\)).
\end{proof}

\begin{lemma}



\leanok
[Sections of an affine pushforward over an arbitrary open]
  \label{lem:gammaPushforwardIsoAt}
  \lean{AlgebraicGeometry.gammaPushforwardIsoAt}
  \uses{lem:globalSectionsIso_hom_comp_specMap_appTop}
  Let \(\varphi : R \to R'\) be a homomorphism of commutative rings, let \(N\) be
  an \(\operatorname{Spec}(R')\)-module, and let \(U\) be an open of
  \(\operatorname{Spec}(R)\), with preimage
  \(V = (\operatorname{Spec}\varphi)^{-1}(U)\) in \(\operatorname{Spec}(R')\). Then
  there is a canonical isomorphism of \(R\)-modules
  \[
    \Gamma\bigl((\operatorname{Spec}\varphi)_* N,\, U\bigr)
    \;\cong\;
    \operatorname{restr}_\varphi\bigl(\Gamma(N,\, V)\bigr) .
  \]
  This is the open-indexed generalization of
  Lemma~\ref{lem:gammaPushforwardIso}, which is the case \(U = \top\) (so that
  \(V = \top\) as well). The family \(\{e_U\}_U\) of these isomorphisms is moreover
  natural in the open argument: the structure-sheaf restriction maps on the two
  sides commute with the isomorphisms, so the family is a natural isomorphism of the
  two presheaves of \(R\)-modules in the open variable (this naturality is
  established in the proof below). For \(U = D(a)\) it is exactly the linear
  equivalence \(e_{D(a)}\) of movement~(1) in the proof of
  Lemma~\ref{lem:pushforward_spec_tilde_iso}.
\end{lemma}

\begin{proof}
  \uses{lem:globalSectionsIso_hom_comp_specMap_appTop}
  The construction copies that of Lemma~\ref{lem:gammaPushforwardIso} verbatim,
  with the top open \(\top\) replaced throughout by the arbitrary open \(U\) (and
  its preimage \(V\)). The point that makes this replacement free is that the
  \(R\)-module structure on the sections is supplied \emph{uniformly} by
  restriction of scalars along the global-sections ring map at the top open,
  independently of the evaluation open; consequently both sides unwind to the same
  nested restriction-of-scalars towers, reconciled by applying the
  restriction-of-scalars composition identity twice and substituting the
  \(\top\)-level ring equation of
  Lemma~\ref{lem:globalSectionsIso_hom_comp_specMap_appTop}. Naturality in the open
  follows because every constituent of the construction is either a structure-sheaf
  restriction map or conjugation by a fixed ring map, and all of these commute with
  further restriction along an inclusion of opens.

  \medskip
  \emph{Naturality of the family \(\{e_U\}_U\) in the open (remark).} The family is
  in fact a natural isomorphism between two presheaves of \(R\)-modules on
  \(\operatorname{Spec}(R)\): the source presheaf
  \(U \mapsto \Gamma\bigl((\operatorname{Spec}\varphi)_* N,\, U\bigr)\), with the
  structure-sheaf restriction maps of the pushforward, and the target presheaf
  \(U \mapsto
    \operatorname{restr}_\varphi\bigl(\Gamma(N,\, (\operatorname{Spec}\varphi)^{-1}U)\bigr)\),
  with the structure-sheaf restriction maps of \(N\) transported through the preimage
  operation \((\operatorname{Spec}\varphi)^{-1}\) and read as \(R\)-linear maps by
  restriction of scalars along \(\varphi\). Concretely, for any inclusion of opens
  \(U' \subseteq U\), with \(V = (\operatorname{Spec}\varphi)^{-1}U\) and
  \(V' = (\operatorname{Spec}\varphi)^{-1}U'\), the square with horizontal maps
  \(e_U, e_{U'}\) and vertical maps the respective structure-sheaf restrictions
  commutes. The reason is structural: each \(e_U\) is a composite whose constituents
  are of exactly two kinds — structure-sheaf restriction maps (of \(N\),
  respectively of \((\operatorname{Spec}\varphi)_* N\)) and identity-on-carrier
  \(\operatorname{restrictScalars}\) repackagings (conjugation by the \emph{fixed},
  open-independent ring map supplied by the restriction-of-scalars reconciliation and
  by the \(\top\)-level ring equation of
  Lemma~\ref{lem:globalSectionsIso_hom_comp_specMap_appTop}). Restriction maps commute
  with further restriction by the presheaf functoriality axiom (using
  \(V' = (\operatorname{Spec}\varphi)^{-1}U' \subseteq
  (\operatorname{Spec}\varphi)^{-1}U = V\), as \((\operatorname{Spec}\varphi)^{-1}\) is
  monotone), and conjugation by a fixed ring map commutes with restriction because the
  ring map does not depend on \(U\); the naturality square is the paste of these
  component squares, and each \(e_U\) being an isomorphism makes the family a natural
  isomorphism. This open-naturality is exactly the compatibility invoked, at the
  inclusion \(D(a) \subseteq \top\), in the proof of
  Lemma~\ref{lem:pushforward_spec_tilde_iso}.
\end{proof}

\begin{lemma}



\leanok
[The tilde restriction from the top open to a basic open is a localization]
  \label{lem:tildeRestriction_isLocalizedModule}
  \lean{AlgebraicGeometry.tildeRestriction_isLocalizedModule}
  Let \(M\) be an \(R'\)-module and \(b \in R'\). Then the structure-sheaf
  restriction map
  \[
    \Gamma(\widetilde M,\, \top) \longrightarrow \Gamma(\widetilde M,\, D(b)),
  \]
  read as a map of \(R'\)-modules, exhibits \(\Gamma(\widetilde M, D(b))\) as the
  localization of \(\Gamma(\widetilde M, \top)\) at \(\mathrm{powers}(b)\). This is
  the \(R'\)-side localization fact transported in movement~(3) of the proof of
  Lemma~\ref{lem:pushforward_spec_tilde_iso}.
\end{lemma}

\begin{proof}


  By the defining property of \(\widetilde{(-)}\), the tilde-localization map
  \(M \to \Gamma(\widetilde M, D(b))\) exhibits its target as the localization of
  \(M\) at \(\mathrm{powers}(b)\). The global-sections map
  \(M \to \Gamma(\widetilde M, \top)\) is the localization of \(M\) at the trivial
  submonoid \(\mathrm{powers}(1)\) — since \(D(1) = \top\) — and is therefore
  bijective. The triangle identity relating these two localization maps to the
  structure-sheaf restriction \(\Gamma(\widetilde M, \top) \to
  \Gamma(\widetilde M, D(b))\) then exhibits the restriction itself, post-composed
  with the inverse of the bijection \(M \cong \Gamma(\widetilde M, \top)\), as the
  localization of \(\Gamma(\widetilde M, \top)\) at \(\mathrm{powers}(b)\).
\end{proof}

We record three further supplementary lemmas driving the basic-open verification
of the affine-pushforward isomorphism below. They carry no external citation.

\begin{lemma}



\leanok
[Restriction of scalars of a localized module]
  \label{lem:powers_restrictScalars}
  \lean{AlgebraicGeometry.IsLocalizedModule.powers_restrictScalars}
  Let \(A\) be a commutative \(R\)-algebra, \(S \subseteq R\) a submonoid, and
  \(f : M \to N\) an \(A\)-linear map that exhibits \(N\) as the localization of
  \(M\) at the image submonoid
  \(\operatorname{algMap}_A(S) \subseteq A\). Then the underlying \(R\)-linear map
  \(\operatorname{restr}_R f\) exhibits \(N\) as the localization of \(M\) at \(S\)
  itself.
\end{lemma}

\begin{proof}



  This is the exact converse of the standard fact that a localization of \(M\) at a
  submonoid \(S\) of \(R\) remains a localization after enlarging the base to an
  \(R\)-algebra \(A\) along the image submonoid \(\operatorname{algMap}_A(S)\). One
  checks the three defining conditions of a localized module directly. Each
  \(s \in S\) acts invertibly on \(N\) because its image
  \(\operatorname{algMap}_A(s)\) does and the two actions coincide by the scalar
  tower \(R \to A \to N\). Surjectivity of \(f\) up to denominators and the
  equality-witness condition transport along the map
  \(s \mapsto \operatorname{algMap}_A(s)\) from \(S\) onto
  \(\operatorname{algMap}_A(S)\), again using
  \(\operatorname{algMap}_A(s)\cdot{} = s\cdot{}\). In the application
  \(A = R'\), \(S = \mathrm{powers}(a)\), one has
  \(\operatorname{algMap}_{R'}(\mathrm{powers}(a)) = \mathrm{powers}(\varphi a)\), so
  this identifies the localization of \(\operatorname{restr}_\varphi M\) at \(a\)
  with the localization of \(M\) at \(\varphi a\); it is the ring-change core of the
  basic-open computation below.
\end{proof}

\begin{lemma}



\leanok
[Section-level counit isomorphism from a localization hypothesis]
  \label{lem:fromTildeGamma_app_isIso_of_localized}
  \lean{AlgebraicGeometry.fromTildeΓ_app_isIso_of_isLocalizedModule}
  Let \(N\) be an \(\operatorname{Spec}(R)\)-module and \(a \in R\). Suppose the
  structure-sheaf restriction map \(\Gamma(N, \top) \to \Gamma(N, D(a))\), read as
  an \(R\)-linear map of the global-section modules, exhibits \(\Gamma(N, D(a))\) as
  the localization of \(\Gamma(N, \top)\) at \(\mathrm{powers}(a)\). Then the
  tilde--\(\Gamma\) counit
  \(\mathrm{fromTilde}\Gamma\,N : \widetilde{\Gamma(N,\top)} \to N\) restricts to an
  isomorphism on the sections over the basic open \(D(a)\).
\end{lemma}

\begin{proof}



  Over \(D(a)\) the section map \(L\) of the counit sits in the triangle identity
  \(L \circ j = \rho\), where
  \(j : \Gamma(N,\top) \to \Gamma\bigl(\widetilde{\Gamma(N,\top)}, D(a)\bigr)\) is
  the tilde-localization map and \(\rho : \Gamma(N,\top) \to \Gamma(N, D(a))\) is
  the structure-sheaf restriction. Both \(j\) and \(\rho\) exhibit their targets as
  the localization of \(\Gamma(N,\top)\) at \(\mathrm{powers}(a)\) — \(j\) by the
  defining property of \(\widetilde{(-)}\), and \(\rho\) by hypothesis. By the
  uniqueness of localized modules there is a unique \(R\)-linear isomorphism \(e\)
  with \(e \circ j = \rho\); the triangle identity then forces \(L = e\), so \(L\) is
  an isomorphism. (This is the per-basic-open input feeding the basis-local
  criterion Lemma~\ref{lem:modules_isIso_of_isBasis}.)
\end{proof}

\begin{lemma}



\leanok
[Affine pushforward of a tilde-module, conditional form]
  \label{lem:pushforward_spec_tilde_iso_conditional}
  \lean{AlgebraicGeometry.pushforward_spec_tilde_iso_of_isLocalizedModule}
  \uses{lem:modules_isIso_of_isBasis, lem:gammaPushforwardTildeIso,
  lem:fromTildeGamma_app_isIso_of_localized}
  Let \(\varphi : R \to R'\) be a homomorphism of commutative rings and \(M\) an
  \(R'\)-module. Suppose that for every \(a \in R\) the structure-sheaf restriction
  of the pushforward \(N := (\operatorname{Spec}\varphi)_*\widetilde M\) exhibits its
  sections over \(D(a)\) as the localization of its global sections at
  \(\mathrm{powers}(a)\). Then there is a canonical isomorphism of
  \(\operatorname{Spec}(R)\)-modules
  \[
    (\operatorname{Spec}\varphi)_*\widetilde M \;\cong\;
    \widetilde{\bigl(\operatorname{restr}_\varphi M\bigr)} .
  \]
\end{lemma}

\begin{proof}
  \uses{lem:modules_isIso_of_isBasis, lem:gammaPushforwardTildeIso, lem:fromTildeGamma_app_isIso_of_localized}
  Write \(N := (\operatorname{Spec}\varphi)_*\widetilde M\). Under the hypothesis,
  Lemma~\ref{lem:fromTildeGamma_app_isIso_of_localized} applies for each \(a \in R\),
  so the counit \(\mathrm{fromTilde}\Gamma\,N\) restricts to an isomorphism on the
  sections over every basic open \(D(a)\). Since the basic opens form a basis of
  \(\operatorname{Spec}(R)\), the basis-local criterion
  Lemma~\ref{lem:modules_isIso_of_isBasis} upgrades this to: the counit
  \(\mathrm{fromTilde}\Gamma\,N\) is an isomorphism, i.e.\ \(N\) lies in the
  essential image of \(\widetilde{(-)}\). Composing the inverse counit with the
  tilde of the global-sections comparison
  Lemma~\ref{lem:gammaPushforwardTildeIso} yields the asserted object isomorphism.
\end{proof}

\begin{lemma}\leanok
[Affine pushforward of a tilde-module]
  \label{lem:pushforward_spec_tilde_iso}
  \lean{AlgebraicGeometry.pushforward_spec_tilde_iso}
  \uses{lem:modules_isIso_of_isBasis, lem:gammaPushforwardTildeIso,
        lem:pushforward_spec_tilde_iso_conditional, lem:powers_restrictScalars,
        lem:gammaPushforwardIso, lem:gammaPushforwardIsoAt,
        lem:tildeRestriction_isLocalizedModule}
% NOTE: The proof reduces the localization criterion to two inputs: the scalar-tower structure
% from \cref{lem:powers_restrictScalars} and the localized-module identification of
% \cref{lem:tildeRestriction_isLocalizedModule}. Open-naturality of the section comparison over
% each basic open \(D(a) \subseteq X\) is established in the proof of
% \cref{lem:gammaPushforwardIsoAt}; each local criterion then follows from the top-level
% localization and that naturality square.
  Let \(\varphi : R \to R'\) be a homomorphism of commutative rings and let \(M\)
  be an \(R'\)-module. Then there is a canonical isomorphism of
  \(\operatorname{Spec}(R)\)-modules
  \[
    \bigl(\operatorname{Spec}\varphi\bigr)_*\,\widetilde{M}
    \;\cong\;
    \widetilde{\bigl(\operatorname{restr}_\varphi M\bigr)},
  \]
  where \(\operatorname{Spec}\varphi : \operatorname{Spec}(R') \to
  \operatorname{Spec}(R)\) is the induced morphism of affine schemes,
  \(\widetilde{(-)}\) denotes the quasi-coherent sheaf associated to a module, and
  \(\operatorname{restr}_\varphi M\) is the \(R\)-module obtained from \(M\) by
  restriction of scalars along \(\varphi\). That is, pushing the quasi-coherent
  sheaf \(\widetilde{M}\) forward along \(\operatorname{Spec}\varphi\) is, up to
  this canonical identification, the tilde of the restriction of scalars of \(M\).

  This is the classical fact that the direct image of a quasi-coherent sheaf along
  a morphism of affine schemes is computed on modules by restriction of scalars; it
  is the affine companion of the pullback formula (the affine case of the
  pushforward analogue of Stacks ``$\widetilde{M}$ and pullback''). It does not
  appear as a ready-made result in Mathlib's quasi-coherent sheaf theory.
\end{lemma}

\begin{proof}
  \uses{lem:modules_isIso_of_isBasis, lem:gammaPushforwardTildeIso, lem:pushforward_spec_tilde_iso_conditional, lem:powers_restrictScalars, lem:gammaPushforwardIso, lem:gammaPushforwardIsoAt, lem:tildeRestriction_isLocalizedModule}
  We build the isomorphism \emph{directly on a basis of basic opens}. This is the
  non-circular route: checking the comparison map on basic opens yields the object
  isomorphism, and quasi-coherence of the pushforward is then a corollary, not a
  prerequisite. (One must not reverse this order: the global-sections comparison
  \(\Gamma\bigl((\operatorname{Spec}\varphi)_*\widetilde{M}\bigr) \cong
  \operatorname{restr}_\varphi M\) of
  Lemma~\ref{lem:gammaPushforwardTildeIso} alone would only upgrade to an object
  isomorphism once the pushforward is already known to be quasi-coherent, so
  quasi-coherence cannot be deduced from an object isomorphism built that way.)

  \emph{The comparison morphism.} There is a canonical morphism of
  \(\operatorname{Spec}(R)\)-modules
  \[
    \alpha : \widetilde{\bigl(\operatorname{restr}_\varphi M\bigr)}
    \longrightarrow \bigl(\operatorname{Spec}\varphi\bigr)_*\widetilde{M},
  \]
  namely the tilde-adjoint of the global-sections identification of
  Lemma~\ref{lem:gammaPushforwardTildeIso}. We show \(\alpha\) is an isomorphism.

  \medskip
  \emph{Reduction to basic opens.} The basic opens \(D(a)\), for \(a \in R\), form a
  basis of the topology of \(\operatorname{Spec}(R)\). By the basis-local criterion
  Lemma~\ref{lem:modules_isIso_of_isBasis} it suffices to check that \(\alpha\)
  restricts to an isomorphism on the sections over each basic open \(D(a)\).

  \medskip
  \emph{The computation over \(D(a)\).} On the source, the sections of
  \(\widetilde{(\operatorname{restr}_\varphi M)}\) over \(D(a)\) are the
  localization \(\bigl(\operatorname{restr}_\varphi M\bigr)[1/a]\) as an
  \(R[1/a]\)-module. On the target, since \((\operatorname{Spec}\varphi)^{-1}D(a) =
  D(\varphi a)\), the sections of \(\bigl(\operatorname{Spec}\varphi\bigr)_*
  \widetilde{M}\) over \(D(a)\) are the sections of \(\widetilde{M}\) over
  \(D(\varphi a)\), i.e.\ the localization \(M[1/\varphi a]\), regarded as an
  \(R[1/a]\)-module by restriction of scalars along the induced map
  \(R[1/a] \to R'[1/\varphi a]\). These two localizations agree: the element \(a\)
  acts on \(\operatorname{restr}_\varphi M\) through \(\varphi a\), so localizing
  \(\operatorname{restr}_\varphi M\) at the multiplicative set generated by \(a\) is
  the same module as localizing \(M\) at the multiplicative set generated by
  \(\varphi a\). Both are the universal localization of
  \(M\) inverting \(\varphi a\), so the comparison map \(\alpha\) is an isomorphism
  on the sections over \(D(a)\). As this holds for every basic open, \(\alpha\) is
  an isomorphism.

  \medskip
  \emph{The explicit formalization route over \(D(a)\).} The agreement of the two
  localizations just described is the entire content, but it must be realized
  through the global ring \(R\), because the \(R\)-action on the pushforward
  sections over \(D(a)\) is built from the global-sections module of
  \(N := (\operatorname{Spec}\varphi)_*\widetilde M\) rather than directly from
  \(R[1/a]\). The reduction therefore proceeds exactly as the conditional
  assembly Lemma~\ref{lem:pushforward_spec_tilde_iso_conditional} prescribes: it
  suffices to supply, for each \(a \in R\), the hypothesis \(\mathrm{hloc}(a)\) that
  the structure-sheaf restriction \(\Gamma(N, \top) \to \Gamma(N, D(a))\) exhibits
  \(\Gamma(N, D(a))\) as the localization of \(\Gamma(N, \top)\) at
  \(\mathrm{powers}(a)\). The \(R\)-action on \(\Gamma(N, D(a))\) is the honest
  restriction-of-scalars action carried by the scalar tower \(R \to R' \to {-}\)
  determined by \(\varphi\); against that structure the ring-change converse
  Lemma~\ref{lem:powers_restrictScalars} is applicable, so the carrier of the
  \(R\)-module is exactly the one to which the localization claim refers, and no
  recursively-defined or ad-hoc section-level scalar action intervenes.
  We obtain \(\mathrm{hloc}(a)\) by transporting the already-known \(R'\)-side
  localization (Lemma~\ref{lem:tildeRestriction_isLocalizedModule}), in
  element-free fashion, across the open-indexed section comparison of
  Lemma~\ref{lem:gammaPushforwardIsoAt} \emph{using its naturality in the open}
  (the naturality remark in the proof of
  Lemma~\ref{lem:gammaPushforwardIsoAt}). This proceeds in three
  movements.

  \emph{(1) The \(D(a)\)-level linear equivalence.} This is the component at
  \(U = D(a)\) of the open-indexed family of
  Lemma~\ref{lem:gammaPushforwardIsoAt}: the \(R\)-linear isomorphism
  \[
    e_{D(a)} : \Gamma\bigl((\operatorname{Spec}\varphi)_*\widetilde M,\, D(a)\bigr)
      \;\cong\;
      \operatorname{restr}_\varphi\bigl(\Gamma(\widetilde M,\, D(\varphi a))\bigr) ,
  \]
  using \((\operatorname{Spec}\varphi)^{-1}D(a) = D(\varphi a)\), which is
  definitional. No fresh \(D(a)\)-level construction is needed: this is precisely the
  open-indexed isomorphism already supplied by
  Lemma~\ref{lem:gammaPushforwardIsoAt} evaluated at the basic open \(D(a)\).

  \emph{(2) The naturality square for \(D(a) \subseteq \top\).} The compatibility
  that relates \(e_{D(a)}\) to the \(\top\)-level isomorphism
  \(e_\top\) (the global-sections comparison of
  Lemma~\ref{lem:gammaPushforwardIso}) is precisely the open-naturality square of the
  family \(\{e_U\}_U\) (established in the proof of
  Lemma~\ref{lem:gammaPushforwardIsoAt}), instantiated at the inclusion
  \(D(a) \subseteq \top\): the square
  \[
    \begin{array}{ccc}
      \Gamma\bigl((\operatorname{Spec}\varphi)_*\widetilde M,\, \top\bigr)
        & \xrightarrow{\;e_\top\;}
        & \operatorname{restr}_\varphi\bigl(\Gamma(\widetilde M,\, \top)\bigr) \\[4pt]
      {\scriptstyle \mathrm{res}_{\top \to D(a)}}\big\downarrow & &
        \big\downarrow{\scriptstyle \mathrm{res}_{\top \to D(\varphi a)}} \\[4pt]
      \Gamma\bigl((\operatorname{Spec}\varphi)_*\widetilde M,\, D(a)\bigr)
        & \xrightarrow{\;e_{D(a)}\;}
        & \operatorname{restr}_\varphi\bigl(\Gamma(\widetilde M,\, D(\varphi a))\bigr)
    \end{array}
  \]
  commutes, using \((\operatorname{Spec}\varphi)^{-1}D(a) = D(\varphi a)\), which is
  definitional. Because the family \(\{e_U\}_U\) is a natural isomorphism, this single
  square is available uniformly in \(a\); there is no per-\(a\) section-level identity
  to re-prove by hand. It intertwines the source restriction
  \(\mathrm{res}_{\top \to D(a)}\) on the pushforward with the (restriction-of-scalars
  reading of the) target restriction \(\mathrm{res}_{\top \to D(\varphi a)}\) on
  \(\widetilde M\).

  \emph{(3) Discharging \(\mathrm{hloc}(a)\) by transport.} The target restriction
  \(\mathrm{res}_{\top \to D(\varphi a)} : \Gamma(\widetilde M, \top) \to
  \Gamma(\widetilde M, D(\varphi a))\) exhibits \(\Gamma(\widetilde M, D(\varphi a)) =
  M[1/\varphi a]\) as the \(\mathrm{powers}(\varphi a)\)-localization of
  \(M = \Gamma(\widetilde M, \top)\) over \(R'\); this is exactly
  Lemma~\ref{lem:tildeRestriction_isLocalizedModule}. Apply the ring-change converse
  Lemma~\ref{lem:powers_restrictScalars} (with \(R\)-algebra \(A = R'\) and submonoid
  \(\mathrm{powers}(a)\), so that
  \(\operatorname{algMap}_{R'}(\mathrm{powers}(a)) = \mathrm{powers}(\varphi a)\)) —
  its \(\operatorname{Algebra} R\, R'\) and scalar-tower hypotheses being supplied by
  the restriction-of-scalars structure along \(\varphi\) — to read this same map as
  the \(\mathrm{powers}(a)\)-localization over \(R\). The commuting naturality square
  of movement (2) intertwines this target restriction with the source restriction
  \(\mathrm{res}_{\top \to D(a)} : \Gamma(N, \top) \to \Gamma(N, D(a))\) by way of the
  isomorphisms \(e_\top\) and \(e_{D(a)}\); transporting the localization property
  across that square — via the transport-along-an-equivalence principle for localized
  modules — yields that \(\mathrm{res}_{\top \to D(a)}\) is the
  \(\mathrm{powers}(a)\)-localization on \(\Gamma(N, D(a))\) over \(R\), i.e.\
  \(\mathrm{hloc}(a)\). No section-level computation specific to the individual \(a\)
  is performed: the only \(a\)-dependent input is the instantiation of the single
  natural-isomorphism square at \(D(a)\). Feeding the family
  \(\{\mathrm{hloc}(a)\}_{a \in R}\) to the conditional assembly
  Lemma~\ref{lem:pushforward_spec_tilde_iso_conditional} produces the unconditional
  object isomorphism \(\alpha\).

  \medskip
  \emph{Naturality in the open drives the transport.} The comparison assembled
  here, \((\operatorname{Spec}\varphi)_* \circ \widetilde{(-)} \cong
  \widetilde{(-)} \circ \operatorname{restr}_\varphi\), is natural in \emph{both}
  arguments — in the module \(M\) and in the open \(U\). The naturality in the
  module is the usual functoriality; the naturality in the open is the naturality
  remark in the proof of Lemma~\ref{lem:gammaPushforwardIsoAt}: the family
  \(\{e_U\}_U\) of Lemma~\ref{lem:gammaPushforwardIsoAt} commutes with the
  structure-sheaf restriction maps on both sides. It is precisely this
  open-naturality that drives the localization transport \emph{uniformly}: the
  per-\(a\) localization fact \(\mathrm{hloc}(a)\) follows from the \(\top\)-level
  localization Lemma~\ref{lem:tildeRestriction_isLocalizedModule} together with that
  naturality square for the inclusion
  \(D(a) \hookrightarrow \top\), rather than from a section-level naturality square
  re-proved by hand for each individual \(a\). This is the single remaining
  obligation of the construction; the carrier of the \(R\)-action on the pushforward
  sections is the restriction-of-scalars structure along \(\varphi\), against which
  Lemma~\ref{lem:powers_restrictScalars} applies directly.

  The \(R\)-module structure on the \(R'\)-side sections to which the ring-change
  converse Lemma~\ref{lem:powers_restrictScalars} is applied is the honest
  restriction-of-scalars structure: \(R\) acts through \(\varphi\), i.e.\ via the
  \(R\)-algebra structure on \(R'\) determined by \(\varphi\) and the associated
  scalar tower \(R \to R' \to M[1/\varphi a]\). It is \emph{not} an ad-hoc,
  recursively-defined scalar action. Lemma~\ref{lem:powers_restrictScalars} is
  stated against exactly this restriction-of-scalars / scalar-tower structure, and
  it is under that structure that localizing \(\operatorname{restr}_\varphi M\) at
  \(a\) coincides with localizing \(M\) at \(\varphi a\).

  % NOTE: This lemma is the affine case of the general fact that the pushforward of
  % a quasi-coherent sheaf along an affine morphism is quasi-coherent; the proof
  % proceeds by the basic-open route established above.
  \medskip
  \emph{Quasi-coherence as a corollary.} The object isomorphism just constructed
  exhibits \(\bigl(\operatorname{Spec}\varphi\bigr)_*\widetilde{M} \cong
  \widetilde{(\operatorname{restr}_\varphi M)}\) as the tilde of a module. Tilde
  modules are quasi-coherent and quasi-coherence of \(\operatorname{Spec}(R)\)-modules
  is closed under isomorphism, so the pushforward
  \(\bigl(\operatorname{Spec}\varphi\bigr)_*\widetilde{M}\) is quasi-coherent. This
  conclusion is stated \emph{after} the object isomorphism, of which it is a
  consequence — no separate ``pushforward preserves quasi-coherent'' input is
  needed for the affine lane.
\end{proof}

\begin{remark}
  This single isomorphism discharges both inputs that the affine reduction below
  requires of the pushforward of a tilde-module.
  \begin{enumerate}
    \item \emph{The global-sections comparison.} Applying the global-sections
      (module-of-sections) functor to the isomorphism identifies the
      \(R\)-module of sections of
      \(\bigl(\operatorname{Spec}\varphi\bigr)_*\widetilde{M}\) with
      \(\operatorname{restr}_\varphi M\); this is exactly the \(\Gamma\)-fragment
      comparison used to recognise the section-level base-change map.
    \item \emph{Quasi-coherence of the pushforward.} The tilde of any module is
      quasi-coherent, and quasi-coherence of \(\operatorname{Spec}(R)\)-modules is
      closed under isomorphism; hence
      \(\bigl(\operatorname{Spec}\varphi\bigr)_*\widetilde{M}\) is quasi-coherent.
      In particular no separate ``pushforward preserves quasi-coherent'' theorem is
      needed for the affine lane.
  \end{enumerate}
  Because both inputs the affine reduction asks of the pushforward are corollaries
  of this single object isomorphism — the section-level module identification and
  the quasi-coherence of the pushforward — this lemma is the sole supplementary
  ingredient the affine lane requires.
\end{remark}

\subsection{The pullback companion}

The affine reduction below requires, alongside the pushforward dictionary, the
\emph{pullback} companion: the affine description of the inverse image of a
tilde-module. It is the affine case of the classical fact that the pullback of a
quasi-coherent sheaf is computed on modules by base change of scalars.

The pullback companion is built by uniqueness of left adjoints from the
per-object \(\Gamma\)-fragment comparison of
Lemma~\ref{lem:gammaPushforwardIso}, packaged as a natural isomorphism of
right-adjoint functors (Lemma~\ref{lem:gammaPushforwardNatIso}).

\begin{lemma}



\leanok
[Naturality packaging of the pushforward \(\Gamma\)-comparison]
  \label{lem:gammaPushforwardNatIso}
  \lean{AlgebraicGeometry.gammaPushforwardNatIso}
  \uses{lem:gammaPushforwardIso}
  The per-object comparison of Lemma~\ref{lem:gammaPushforwardIso} assembles into a
  natural isomorphism of the two right-adjoint functors
  \((\operatorname{Spec}\varphi)_* \circ \widetilde{(-)}\) and
  \(\widetilde{(-)} \circ \operatorname{restr}_\varphi\) on modules, exhibiting the
  pushforward dictionary as a natural transformation rather than merely a family of
  object isomorphisms. This is the input to the uniqueness-of-left-adjoints argument
  that produces the pullback companion Lemma~\ref{lem:pullback_spec_tilde_iso}.
\end{lemma}
\begin{proof}
  \uses{lem:gammaPushforwardIso}
  The components are the object isomorphisms of Lemma~\ref{lem:gammaPushforwardIso};
  naturality is the corresponding naturality square, which holds because every
  constituent of the comparison is identity-on-elements.
\end{proof}

\begin{lemma}



\leanok
[Three-fold naturality of the global-sections comparison square]
  \label{lem:globalSectionsIso_hom_comp3_specMap_appTop}
  \lean{AlgebraicGeometry.globalSectionsIso_hom_comp3_specMap_appTop}
  \uses{lem:globalSectionsIso_hom_comp_specMap_appTop}
  Let \(\varphi : R \to S\) and \(\rho : S \to T\) be homomorphisms of commutative
  rings, and write \(\mathrm{gs}_R, \mathrm{gs}_S, \mathrm{gs}_T\) for the
  global-sections isomorphisms of the three structure sheaves. By contravariance of
  \(\operatorname{Spec}\) one has
  \(\operatorname{Spec}(\rho\circ\varphi) = \operatorname{Spec}\varphi \circ
  \operatorname{Spec}\rho\); transporting the top-open structure-sheaf component along
  this identity expresses the composite component as the pasting of the two single-map
  components,
  \[
    (\operatorname{Spec}(\rho\circ\varphi))^{\sharp}_{\top}
      \;=\;
    (\operatorname{Spec}\varphi)^{\sharp}_{\top}
      \;\text{followed by}\;
    (\operatorname{Spec}\rho)^{\sharp}_{\top},
  \]
  the two presentations being identified by the cast carried by that contravariance
  identity. Equivalently, conjugating by the global-sections isomorphisms, the ring
  square
  \[
    \mathrm{gs}_R \;\circ\; (\operatorname{Spec}(\rho\circ\varphi))^{\sharp}_{\top}
    \;=\; (\rho\circ\varphi) \;\circ\; \mathrm{gs}_T
  \]
  commutes. This is the three-fold \((R \to S \to T)\) composite analogue of the
  single-map square \cref{lem:globalSectionsIso_hom_comp_specMap_appTop}.
\end{lemma}

\begin{proof}
  \uses{lem:globalSectionsIso_hom_comp_specMap_appTop}
  By the contravariance of \(\operatorname{Spec}\) the morphism
  \(\operatorname{Spec}(\rho\circ\varphi)\) equals the composite
  \(\operatorname{Spec}\varphi \circ \operatorname{Spec}\rho\). Since the structure-sheaf
  comparison is multiplicative under composition of morphisms, transporting the top-open
  component along this identity expresses
  \((\operatorname{Spec}(\rho\circ\varphi))^{\sharp}_{\top}\) as the pasting of
  \((\operatorname{Spec}\varphi)^{\sharp}_{\top}\) followed by
  \((\operatorname{Spec}\rho)^{\sharp}_{\top}\). Now apply the single-map square
  \cref{lem:globalSectionsIso_hom_comp_specMap_appTop} twice: for \(\varphi\) it gives
  \(\mathrm{gs}_R \circ (\operatorname{Spec}\varphi)^{\sharp}_{\top}
  = \varphi \circ \mathrm{gs}_S\), and for \(\rho\) it gives
  \(\mathrm{gs}_S \circ (\operatorname{Spec}\rho)^{\sharp}_{\top}
  = \rho \circ \mathrm{gs}_T\). Splicing the two squares along \(\mathrm{gs}_S\) yields
  \(\mathrm{gs}_R \circ (\operatorname{Spec}(\rho\circ\varphi))^{\sharp}_{\top}
  = (\rho\circ\varphi) \circ \mathrm{gs}_T\), as claimed.
\end{proof}

\begin{lemma}


\leanok
[Bridge: the pushforward-composition cast on \(\Gamma\) is an \texttt{eqToHom}]
  \label{lem:gammaPushforwardIso_comp_bridge}
  \lean{AlgebraicGeometry.gammaPushforwardIso_comp_bridge}
  For ring maps \(\varphi : R \to S\), \(\rho : S \to T\) and a
  \(\operatorname{Spec}(T)\)-module \(N\), the sheaf-of-modules morphism on
  \(\operatorname{Spec} R\) obtained as the cast carried by the pushforward-composition
  identity \((\operatorname{Spec}(\rho\circ\varphi))_* = (\operatorname{Spec}\varphi)_*\circ
  (\operatorname{Spec}\rho)_*\), followed by the inverse component of the pushforward
  composition isomorphism, has the identity section at every open. Consequently its image
  under the global-sections functor \(\Gamma_R\) is the \texttt{eqToHom} of the induced
  object equality. This isolates the single value-module-category / scheme-module junction
  crossing of \cref{lem:gammaPushforwardIso_comp} into one bounded morphism identity.
\end{lemma}

\begin{proof}
  The pushforward composition isomorphism is, underneath, the identity isomorphism, so its
  inverse component at \(N\) is the identity morphism, and the cast factor is an
  \texttt{eqToHom} by construction; their composite is therefore an \texttt{eqToHom}.
  Applying the global-sections functor and using that the functor sends a morphism to its
  top-open section, the claim reduces, by morphism extensionality, to the per-open identity
  of the pushforward-composition inverse, reconciled with the \texttt{eqToHom} on the
  underlying carrier by the \texttt{eqToHom}-on-elements identity. No element of a
  substantive composite is reduced — the lone junction crossing happens once, here.
\end{proof}

\begin{lemma}



\leanok
[Per-component composition coherence of the affine-pushforward \(\Gamma\)-comparison]
  \label{lem:gammaPushforwardIso_comp}
  \lean{AlgebraicGeometry.gammaPushforwardIso_comp}
  \uses{lem:gammaPushforwardIso, lem:gammaPushforwardIso_hom_apply, lem:gammaPushforwardIso_comp_bridge}
  Let \(\varphi : R \to S\) and \(\rho : S \to T\) be homomorphisms of commutative
  rings and let \(N\) be an \(\operatorname{Spec}(T)\)-module. Then the global-sections
  comparison \cref{lem:gammaPushforwardIso} for the composite ring map
  \(\rho\circ\varphi\),
  \[
    \Gamma\bigl((\operatorname{Spec}(\rho\circ\varphi))_* N\bigr)
      \;\cong\;
    \operatorname{restr}_{\rho\circ\varphi}\bigl(\Gamma N\bigr),
  \]
  agrees with the composite of the two single-map comparisons — that for \(\rho\)
  applied to \(N\), and that for \(\varphi\) applied to the pushforward
  \((\operatorname{Spec}\rho)_* N\) — once the codomain is reorganised by the
  restrict-scalars composition identity
  \(\operatorname{restr}_{\rho\circ\varphi} = \operatorname{restr}_\varphi \circ
  \operatorname{restr}_\rho\) and the domain by the pushforward composition identity
  \((\operatorname{Spec}(\rho\circ\varphi))_* = (\operatorname{Spec}\varphi)_* \circ
  (\operatorname{Spec}\rho)_*\). This is the component at \(N\) of the coherence
  packaged in \cref{lem:gammaPushforwardNatIso_comp}, and it is the component where the
  genuine global-sections ring content lives.
\end{lemma}

\begin{proof}
  \uses{lem:gammaPushforwardIso, lem:gammaPushforwardIso_hom_apply, lem:gammaPushforwardIso_comp_bridge}
  We argue at the level of module morphisms. The realized route crosses the value/scheme
  junction exactly once, inside the bridge \cref{lem:gammaPushforwardIso_comp_bridge}; after
  substituting it, the comparison collapses on the carrier by
  \cref{lem:gammaPushforwardIso_hom_apply} and the residual cast is discharged by
  proof-irrelevance. The three-fold ring coherence
  \cref{lem:globalSectionsIso_hom_comp3_specMap_appTop} that the idealized argument below
  invokes is, in the realized proof, absorbed definitionally inside the per-map comparisons,
  so it is not directly cited.
  Both the comparison \cref{lem:gammaPushforwardIso} for \(\rho\circ\varphi\) and
  the right-hand composite are morphisms in the module category over \(R\), and the
  identity to be proved is an equality of two such morphisms.

  The decisive structural observation is that the right-hand composite is, verbatim, the
  composition law
  \[
    \mathrm{map}(\rho\circ\varphi)
      = \mathrm{map}(\varphi)\;\circ\;
        \operatorname{restr}_\varphi\bigl(\mathrm{map}(\rho)\bigr)\;\circ\;
        \bigl(\operatorname{restr}_{\rho\circ\varphi} \cong
        \operatorname{restr}_\varphi\circ\operatorname{restr}_\rho\bigr)^{-1}
  \]
  of a presheaf-of-modules-shaped structure, with \(\mathrm{map}\) read as the
  \(\Gamma\)-comparison, once two bookkeeping identifications are in place: the domain is
  corrected by the pushforward composition identity
  \((\operatorname{Spec}(\rho\circ\varphi))_* = (\operatorname{Spec}\varphi)_* \circ
  (\operatorname{Spec}\rho)_*\), and the codomain by the restrict-scalars composition
  identity \(\operatorname{restr}_{\rho\circ\varphi} = \operatorname{restr}_\varphi \circ
  \operatorname{restr}_\rho\). The proof is the explicit verification of this single
  coherence square.

  First normalise the codomain glue: write the restrict-scalars composition isomorphism
  \(\operatorname{restr}_{\rho\circ\varphi} \cong \operatorname{restr}_\varphi \circ
  \operatorname{restr}_\rho\) in its explicit object form at the section module, so that
  its inverse appears as a single named coherence morphism over \(\Gamma(N,\top)\). Next
  unfold each of the three \(\Gamma\)-comparisons — for \(\rho\circ\varphi\), for
  \(\varphi\), and for \(\rho\) — to its defining restrict-scalars-composition composite,
  again purely as an equality of morphisms.

  The workhorse step re-aligns the two sides without touching any carrier: the middle
  factor, the \(\operatorname{restr}_\varphi\)-image of the comparison for \(\rho\), is
  slid leftward past the coherence isomorphism using the morphism-level naturality of the
  restrict-scalars composition isomorphism. This naturality square expresses that the
  composition coherence is natural in the module argument; applying it brings the two
  composites into a common shape.

  At this point the value/scheme junction is crossed exactly once, isolated in a single
  morphism-level bridge identity for the domain reindexing cast \(\Gamma(\text{cast})\)
  induced by the pushforward composition identity. That bridge is a standalone equality of
  morphisms whose proof is the lone junction definitional identity — kernel-cheap precisely
  because it is taken in isolation rather than distributed across the whole composite.
  After substituting it the goal carries no pushforward and no value-category functor: it
  is a pure restrict-scalars coherence identity over \(\Gamma(N,\top)\).

  Finally, the two sides differ only by a restrict-scalars congruence isomorphism evaluated
  at two endpoints. Those endpoints are equal — the equality is exactly the three-fold
  global-sections ring coherence \cref{lem:globalSectionsIso_hom_comp3_specMap_appTop} — and
  since the congruence isomorphism depends on its endpoints only through an equality of
  propositional type, the two congruences coincide by proof-irrelevance of that endpoint
  equality. This closes the residual endpoint mismatch with no element argument, and the
  two morphisms agree.

  The clean structural fix, of which the above is the hand verification, is to reconstruct
  the \(\Gamma\)-comparison so that the value/scheme junction is crossed once and centrally
  — as the composition (\(\mathrm{map\_comp}\)) field of a presheaf-of-modules /
  restrict-scalars pseudofunctor structure — after which this coherence, together with any
  future unit and pentagon coherence, is discharged once at construction by pure
  restrict-scalars pseudofunctor algebra rather than re-derived per use.
\end{proof}

\begin{lemma}



\leanok
[Carrier action of the affine-pushforward \(\Gamma\)-comparison]
  \label{lem:gammaPushforwardIso_hom_apply}
  \lean{AlgebraicGeometry.gammaPushforwardIso_hom_apply}
  \uses{lem:gammaPushforwardIso}
  For \(\varphi : R \to S\) and an \(\operatorname{Spec}(S)\)-module \(N\), the forward
  comparison map of \cref{lem:gammaPushforwardIso} acts as the identity on the underlying
  carrier: for any global section \(x\), \((\mathrm{gammaPushforwardIso}\,\varphi\,N).\mathrm{hom}\,x = x\).
\end{lemma}

\begin{proof}
  \uses{lem:gammaPushforwardIso}
  Every constituent of the comparison isomorphism
  (\(\mathrm{restrictScalarsComp'App}\), \(\mathrm{restrictScalarsCongr}\)) is the identity
  on the underlying abelian group of global sections; only the recorded module structure
  changes. Hence the underlying map of \(\mathrm{hom}\) is the identity.
\end{proof}

\begin{lemma}



\leanok
[Carrier action of the inverse affine-pushforward \(\Gamma\)-comparison]
  \label{lem:gammaPushforwardIso_inv_apply}
  \lean{AlgebraicGeometry.gammaPushforwardIso_inv_apply}
  \uses{lem:gammaPushforwardIso}
  Under the same hypotheses, the inverse comparison map of \cref{lem:gammaPushforwardIso}
  also acts as the identity on the carrier: \((\mathrm{gammaPushforwardIso}\,\varphi\,N).\mathrm{inv}\,x = x\).
\end{lemma}

\begin{proof}
  \uses{lem:gammaPushforwardIso}
  Identical to \cref{lem:gammaPushforwardIso_hom_apply}, applied to the inverse, which is
  built from the same identity-on-carrier constituents.
\end{proof}

\begin{lemma}



\leanok
[Composition coherence of the pushforward \(\Gamma\)-comparison]
  \label{lem:gammaPushforwardNatIso_comp}
  \lean{AlgebraicGeometry.gammaPushforwardNatIso_comp}
  \uses{lem:gammaPushforwardNatIso, lem:gammaPushforwardIso_comp, lem:globalSectionsIso_hom_comp3_specMap_appTop}
  The naturality packaging \cref{lem:gammaPushforwardNatIso} is coherent under composition of
  ring maps: for ring maps \(\varphi : R \to S\) and \(\rho : S \to T\), the comparison for the
  composite \(\rho \circ \varphi\) is the pasting of the two single-map comparisons
  \(\bigl(\text{\cref{lem:gammaPushforwardNatIso}}\bigr)_\varphi\) and
  \(\bigl(\text{\cref{lem:gammaPushforwardNatIso}}\bigr)_\rho\), glued on the domain side by the
  pushforward composition isomorphism \(\mathrm{pushforwardComp}\) and on the codomain side by the
  restrict-scalars composition isomorphism \(\mathrm{restrictScalarsComp}\). This is the foundation
  coherence whose conjugate through the composite-adjunction mate equivalence produces the
  pseudofunctoriality of the tilde-pullback dictionary used in
  \cref{lem:pullback_spec_tilde_iso_ring_square_mate_glue}.

  The domain-side gluing isomorphism \(\mathrm{pushforwardComp}\) carries the
  \(\operatorname{Spec}\)-contravariance cast (an \(\mathrm{eqToIso}\)) coming from the
  identity \(\operatorname{Spec}(\rho\circ\varphi) = \operatorname{Spec}\varphi \circ
  \operatorname{Spec}\rho\). Far from being inert, that cast is the locus of the genuine
  content of the lemma: it is what reconciles the composite structure-sheaf comparison
  with the pasting of the two single-map comparisons, and unwinding it is precisely the
  three-fold global-sections ring coherence
  \cref{lem:globalSectionsIso_hom_comp3_specMap_appTop}.
  \begin{proof}
    \uses{lem:gammaPushforwardNatIso, lem:gammaPushforwardIso_comp, lem:globalSectionsIso_hom_comp3_specMap_appTop}
    On the domain side the gluing is trivial: the pushforward composition isomorphism
    \(\mathrm{pushforwardComp}\) acts as the identity on every module, so it contributes no
    genuine content. The substance is therefore entirely in reconciling the two natural
    isomorphisms component by component.

    By extensionality of natural transformations between functors, it suffices to verify
    the asserted equality on a single, arbitrary component; there is no need to expand
    every constituent of the two composite isomorphisms simultaneously, and it is precisely
    this reduction to a per-component check that keeps the verification tractable. Fixing a
    module \(N\), the component equality is exactly the per-component composition coherence
    \cref{lem:gammaPushforwardIso_comp} of the global-sections comparison
    \(\mathrm{gammaPushforwardIso}\) for the composite ring map \(\rho\circ\varphi\).

    That component coherence is not a pointwise identity. The comparison for the composite
    is indexed by the top-open structure-sheaf component of
    \(\operatorname{Spec}(\rho\circ\varphi)\), whereas the pasting of the two single-map
    comparisons is indexed by the top-open components of \(\operatorname{Spec}\varphi\) and
    \(\operatorname{Spec}\rho\) joined through the \(\operatorname{Spec}\)-contravariance
    cast. Reconciling these two indexings is precisely the three-fold global-sections ring
    coherence \cref{lem:globalSectionsIso_hom_comp3_specMap_appTop}, which discharges the
    component equality and thereby completes the proof.
  \end{proof}
\end{lemma}

\begin{lemma}



\leanok
[Affine pullback of a tilde-module]
  \label{lem:pullback_spec_tilde_iso}
  \lean{AlgebraicGeometry.pullback_spec_tilde_iso}
  \uses{lem:pushforward_spec_tilde_iso, lem:gammaPushforwardNatIso}
  % SOURCE: [Stacks Project], Schemes, Lemma `lemma-widetilde-pullback', Tag 01I9
  % (\S\,Quasi-coherent sheaves on affines) (read from references/stacks-schemes.tex,
  % L1241--1256).
  % SOURCE QUOTE: "Let
  % $(X, \mathcal{O}_X) = (\Spec(S), \mathcal{O}_{\Spec(S)})$,
  % $(Y, \mathcal{O}_Y) = (\Spec(R), \mathcal{O}_{\Spec(R)})$
  % be affine schemes.
  % Let $\psi : (X, \mathcal{O}_X) \to (Y, \mathcal{O}_Y)$ be a
  % morphism of affine schemes, corresponding to the ring map
  % $\psi^\sharp : R \to S$ (see Lemma \ref{lemma-category-affine-schemes}).
  % \begin{enumerate}
  % \item We have $\psi^* \widetilde M = \widetilde{S \otimes_R M}$
  % functorially in the $R$-module $M$.
  % \item We have $\psi_* \widetilde N = \widetilde{N_R}$ functorially
  % in the $S$-module $N$.
  % \end{enumerate}"
  \textit{Source: Stacks Project, Schemes, Lemma ``\(\widetilde M\) and pullback''
  (Tag 01I9), part (1).}
  Let \(\varphi : R \to R'\) be a homomorphism of commutative rings and let \(M\)
  be an \(R\)-module. Then there is a canonical isomorphism of
  \(\operatorname{Spec}(R')\)-modules
  \[
    \bigl(\operatorname{Spec}\varphi\bigr)^{*}\,\widetilde{M}
    \;\cong\;
    \widetilde{\bigl(R' \otimes_R M\bigr)},
  \]
  where \(\operatorname{Spec}\varphi : \operatorname{Spec}(R') \to
  \operatorname{Spec}(R)\) is the induced morphism of affine schemes,
  \(\widetilde{(-)}\) denotes the quasi-coherent sheaf associated to a module, and
  \(R' \otimes_R M\) is the base change of \(M\) along \(\varphi\) — the
  \(R'\)-module obtained by extension of scalars. The isomorphism is natural in
  \(M\). That is, pulling the quasi-coherent sheaf \(\widetilde{M}\) back along
  \(\operatorname{Spec}\varphi\) is, up to this canonical identification, the tilde
  of the base change of \(M\). It is the pullback companion of the affine
  pushforward dictionary Lemma~\ref{lem:pushforward_spec_tilde_iso}, and is part~(1)
  of the same Stacks lemma of which the pushforward formula is part~(2).
\end{lemma}

% SOURCE QUOTE PROOF: "The first assertion follows from the identification in
% Lemma \ref{lemma-compare-constructions}
% and the result of Modules, Lemma \ref{modules-lemma-restrict-quasi-coherent}."
\begin{proof}
  \uses{lem:pushforward_spec_tilde_iso}
  Pullback along \(\operatorname{Spec}\varphi\) is, on quasi-coherent sheaves, the
  extension-of-scalars operation \(- \otimes_R R'\) on modules transported through
  the \(\widetilde{(-)}\) / \(\operatorname{moduleSpec}\Gamma\) equivalence between
  modules and quasi-coherent sheaves. Concretely, the global sections of
  \(\bigl(\operatorname{Spec}\varphi\bigr)^{*}\widetilde{M}\) are the base change
  \(R' \otimes_R M\): the pullback functor \((\operatorname{Spec}\varphi)^{*}\) is
  left adjoint to the pushforward \((\operatorname{Spec}\varphi)_{*}\), and under
  the global-sections identification the latter is restriction of scalars along
  \(\varphi\) (Lemma~\ref{lem:pushforward_spec_tilde_iso}); the left adjoint of
  restriction of scalars is extension of scalars \(M \mapsto R' \otimes_R M\). The
  comparison map is the unit of the base-change adjunction, and the universal
  property of base change identifies it with the canonical isomorphism
  \(\bigl(\operatorname{Spec}\varphi\bigr)^{*}\widetilde{M} \cong
  \widetilde{(R' \otimes_R M)}\). Naturality in \(M\) is the naturality of that
  unit. (One may equally read this off the standard-open description: on
  \(D(\varphi a)\) the localization \((R' \otimes_R M)[1/a]\) is
  \(R'[1/\varphi a] \otimes_{R[1/a]} M[1/a]\), which is the base change of the
  sections of \(\widetilde M\) over \(D(a)\).)
\end{proof}

\begin{lemma}


\leanok
[Functor-level affine pullback dictionary]
  \label{lem:pullbackSpecTildeNatIso}
  \lean{AlgebraicGeometry.pullbackSpecTildeNatIso}
  \uses{lem:gammaPushforwardNatIso, lem:conjugateIsoEquiv_mathlib}
  The per-module affine pullback dictionary \cref{lem:pullback_spec_tilde_iso} assembles into a
  natural isomorphism of functors
  \(\mathrm{ModuleCat}\,R \to (\operatorname{Spec} R').\mathrm{Modules}\),
  \[
    \widetilde{(-)} \ggg \mathrm{pullback}(\operatorname{Spec}\varphi)
      \;\cong\;
    \mathrm{ext}_\varphi \ggg \widetilde{(-)},
  \]
  whose component at \(M\) is the isomorphism \cref{lem:pullback_spec_tilde_iso} for \(M\). It is the
  functor-level packaging of that dictionary, obtained as the
  \(\mathrm{conjugateIsoEquiv}\)-inverse \cref{lem:conjugateIsoEquiv_mathlib} of the
  \(\Gamma\)-pushforward natural isomorphism \cref{lem:gammaPushforwardNatIso}, for the two composite
  adjunctions
  \[
    \mathrm{adj}_L = (\widetilde{(-)}\dashv\Gamma)\,.\mathrm{comp}\,
      (\mathrm{pullback}\dashv\mathrm{pushforward}),
    \quad
    \mathrm{adj}_R = (\mathrm{ext}_\varphi\dashv\mathrm{restr}_\varphi)\,.\mathrm{comp}\,
      (\widetilde{(-)}\dashv\Gamma).
  \]
  This packaging lets the ring-square glue be carried out at the natural-transformation level, without
  evaluating any component.
  \begin{proof}
    \uses{lem:gammaPushforwardNatIso, lem:conjugateIsoEquiv_mathlib}
    Apply the iso-conjugation equivalence \cref{lem:conjugateIsoEquiv_mathlib} for
    \(\mathrm{adj}_L,\mathrm{adj}_R\), in its inverse direction, to the pushforward
    \(\Gamma\)-comparison \cref{lem:gammaPushforwardNatIso}, and take the inverse isomorphism of the
    result. The source and target functors are as displayed, and the component at \(M\) is
    \cref{lem:pullback_spec_tilde_iso} for \(M\) by construction.
  \end{proof}
\end{lemma}

\begin{lemma}


\leanok
[Component of the functor-level pullback dictionary]
  \label{lem:pullbackSpecTildeNatIso_app}
  \lean{AlgebraicGeometry.pullbackSpecTildeNatIso_app}
  \uses{lem:pullbackSpecTildeNatIso, lem:pullback_spec_tilde_iso}
  For every \(R\)-module \(M\), the component of the functor-level dictionary
  \cref{lem:pullbackSpecTildeNatIso} at \(M\) is the per-module dictionary
  \cref{lem:pullback_spec_tilde_iso} for \(M\):
  \[
    \bigl(\mathrm{pullbackSpecTildeNatIso}\,\varphi\bigr).\mathrm{app}\,M
      = \mathrm{pullback\_spec\_tilde\_iso}\,\varphi\,M .
  \]
  \begin{proof}
    \uses{lem:pullbackSpecTildeNatIso}
    The per-module dictionary is, by the construction of \cref{lem:pullbackSpecTildeNatIso}, exactly
    this component, so the identity holds definitionally.
  \end{proof}
\end{lemma}

\begin{lemma}



\leanok
[Defining unit triangle of the affine pullback dictionary]
  \label{lem:pullback_spec_tilde_iso_inv_unit_triangle}
  \lean{AlgebraicGeometry.pullback_spec_tilde_iso_inv_unit_triangle}
  \uses{lem:pullback_spec_tilde_iso, lem:gammaPushforwardNatIso}
  The inverse of the affine pullback dictionary \cref{lem:pullback_spec_tilde_iso} is, by
  construction, the adjoint conjugate of the \(\Gamma\)-pushforward natural isomorphism
  \cref{lem:gammaPushforwardNatIso}; this is recorded by its \emph{defining unit triangle}.
  Write \(L = \widetilde{(-)} \ggg \mathrm{pullback}(\operatorname{Spec}\varphi)\) for the
  composite left adjoint (right adjoint \(\mathrm{pushforward}(\operatorname{Spec}\varphi)\ggg
  \Gamma_R\)) and \(L' = \mathrm{ext}_\varphi \ggg \widetilde{(-)}\) for the algebraic one
  (right adjoint \(\Gamma_{R'}\ggg\mathrm{restr}_\varphi\)). Then, for every \(R\)-module \(M\),
  \[
    \eta^{L}_M \mathbin{;} \bigl(\text{\cref{lem:gammaPushforwardNatIso}}\bigr)_{L.\mathrm{obj}\,M}
      \;=\;
    \eta^{L'}_M \mathbin{;} \bigl(\Gamma_{R'}\ggg\mathrm{restr}_\varphi\bigr)
              \!\bigl((\text{\cref{lem:pullback_spec_tilde_iso}}_M)^{-1}\bigr),
  \]
  where \(\eta^{L}, \eta^{L'}\) are the two units. This is the per-leg bridge
  ``\(\mathrm{pullback\_spec\_tilde\_iso}\rightsquigarrow\mathrm{gammaPushforwardNatIso}\)''
  consumed by the ring-square crux \cref{lem:pullback_spec_tilde_iso_ring_square_natural}.
  \begin{proof}
    \uses{lem:pullback_spec_tilde_iso, lem:gammaPushforwardNatIso}
    By construction \cref{lem:pullback_spec_tilde_iso} is the \(M\)-component of
    \((\mathrm{conjugateIsoEquiv}\,\mathrm{adj}_L\,\mathrm{adj}_{L'})^{-1}\) applied to
    \cref{lem:gammaPushforwardNatIso}, so its inverse is the conjugate of that
    \(\Gamma\)-comparison and the bridge
    \(\bigl((\mathrm{conjugateEquiv}\,\mathrm{adj}_L\,\mathrm{adj}_{L'})^{-1}
    (\text{\cref{lem:gammaPushforwardNatIso}}).\mathrm{hom}\bigr)_M
    = (\text{\cref{lem:pullback_spec_tilde_iso}}_M)^{-1}\) holds by \texttt{rfl}. The displayed
    triangle is then exactly Mathlib's \texttt{CategoryTheory.unit\_conjugateEquiv\_symm} for this
    conjugate. (One line; the Lean declaration is axiom-clean.)
  \end{proof}
\end{lemma}

\subsection{The affine base-change lemma and its remaining obligations}

The affine base-change lemma below proceeds \emph{directly on global sections}: it
reduces, via the locality criterion
Lemma~\ref{lem:modules_isIso_iff_affineOpens} and the two affine dictionaries
Lemmas~\ref{lem:pushforward_spec_tilde_iso} and~\ref{lem:pullback_spec_tilde_iso},
to two obligations that are isolated here as named blocks. Both are extracted from
the proof of Stacks ``Affine base change'', and both are absent from Mathlib. The
first is the \emph{locality reduction}: restricting to the affine opens of \(S'\)
reduces the isomorphism claim to the affine--affine model
\(S = \operatorname{Spec}(R)\), \(S' = \operatorname{Spec}(R')\),
\(X = \operatorname{Spec}(A)\), \(\mathcal{F} = \widetilde{M}\). The second is the
\emph{section-level identification}: in that affine--affine model, the global-sections
incarnation \(\Gamma(\alpha)\) of the base-change map, transported through the two
proved tilde dictionaries, is the canonical cancellation isomorphism
\(\operatorname{cancelBaseChange}\) — an isomorphism with no flatness hypothesis.
The proof proceeds by locality on \(S'\) followed by a pure computation on the module
of global sections. No abstract adjoint-mate identification at the level of sheaves is
needed; the mate is unwound \emph{on sections}, where it becomes
\(\operatorname{cancelBaseChange}\) directly.

\begin{lemma}



\leanok
[Affine-local compatibility of the base-change map]
  \label{lem:base_change_map_affine_local}
  \lean{AlgebraicGeometry.base_change_map_affine_local}
  \uses{def:pushforward_base_change_map, lem:modules_isIso_iff_affineOpens}
  % SOURCE: [Stacks Project], Cohomology of Schemes, Lemma
  % `lemma-affine-base-change' (\S\,Cohomology and base change, I), proof, the
  % "local on $S$ and $S'$" step (read from references/stacks-coherent.tex,
  % L920--926).
  % SOURCE QUOTE: "The statement is local on $S$ and $S'$. Hence we may
  % assume $X = \Spec(A)$, $S = \Spec(R)$,
  % $S' = \Spec(R')$ and $\mathcal{F} = \widetilde{M}$
  % for some $A$-module $M$."
  \textit{Source: Stacks Project, Cohomology of Schemes, Lemma ``Affine base
  change'' (proof, local-on-\(S\)-and-\(S'\) step).}
  Let \(f : X \to S\) be an affine morphism and \(\mathcal{F}\) a quasi-coherent
  \(\mathcal{O}_X\)-module, with a base-change square as above. The formation of the
  base-change map (Definition~\ref{def:pushforward_base_change_map}) is compatible
  with restriction to affine opens of \(S'\): for each affine open \(U \subseteq S'\), the
  restriction of the base-change map to sections over \(U\) agrees with the
  base-change map of the square obtained by restricting along \(U\). Consequently,
  by the affine-open locality criterion
  Lemma~\ref{lem:modules_isIso_iff_affineOpens}, the assertion that the base-change
  map is an isomorphism over an arbitrary base-change square \((S, S', X, X')\)
  reduces to the affine--affine case \(S = \operatorname{Spec}(R)\),
  \(S' = \operatorname{Spec}(R')\), \(X = \operatorname{Spec}(A)\),
  \(\mathcal{F} = \widetilde{M}\) for an \(A\)-module \(M\).

\end{lemma}

\begin{proof}
  \uses{def:pushforward_base_change_map, lem:modules_isIso_iff_affineOpens}
  By the affine-open locality criterion
  Lemma~\ref{lem:modules_isIso_iff_affineOpens}, the base-change map is an isomorphism
  of \(\mathcal{O}_{S'}\)-modules as soon as its restriction to the sections over each
  affine open \(U \subseteq S'\) is an isomorphism. The only content beyond the bare criterion is to identify that
  restriction with the base-change map of the affine--affine square cut out over
  \(U\). We obtain the identification by unfolding the definition of the map and
  pushing the restriction functor through each of its building blocks. The
  compatibility this needs is \emph{not} a new coherence theorem: it is the naturality
  of the adjunction transpose, together with the standard fact that pushforward
  commutes with restriction to an open of the base. We spell this out in three steps.

  \medskip
  \emph{Step 1: unfold the map.} By
  Definition~\ref{def:pushforward_base_change_map}, the base-change map is the image,
  under the \((g^*, g_*)\)-adjunction transpose — the natural bijection
  \(\operatorname{Hom}(g^*A, B) \cong \operatorname{Hom}(A, g_*B)\) — of the morphism
  \[
    f_*\mathcal{F} \xrightarrow{\,f_*(\eta_{\mathcal{F}})\,}
    f_*(g')_*(g')^*\mathcal{F} = g_*\,f'_*(g')^*\mathcal{F},
  \]
  where \(\eta_{\mathcal{F}} : \mathcal{F} \to (g')_*(g')^*\mathcal{F}\) is the unit of
  the \(((g')^*, (g')_*)\)-adjunction and the identification
  \(f_*(g')_* = g_*\,f'_*\) is the pseudofunctoriality of pushforward applied to the
  square's commutativity \(g \circ f' = f \circ g'\). Thus the map is, by
  construction, the transpose of the \(f_*\)-image of a unit — a composite of
  adjunction data and pushforward functors, with no ingredient beyond these.

  \medskip
  \emph{Step 2: restriction to \(U\) commutes with each building block.} Fix an affine
  open \(U \subseteq S'\) and write \(\rho_U\) for restriction to \(U\) —
  equivalently, the section functor \(\Gamma(U, -)\) on \(\mathcal{O}_{S'}\)-modules,
  or pullback along the open immersion \(U \hookrightarrow S'\). Restriction along
  \(U\) is a functor, and every ingredient of Step 1 is a natural transformation or a
  natural bijection; so applying \(\rho_U\) to the transpose-of-a-unit of Step 1
  returns the transpose-of-the-unit built from the restricted data. Concretely:
  \begin{itemize}
    \item[(a)] the unit \(\eta_{\mathcal{F}}\) of the \(((g')^*, (g')_*)\)-adjunction
      restricts to the unit of the corresponding adjunction on the restricted square,
      because units are natural in their object and pulling the cartesian square back
      along \(U \hookrightarrow S'\) yields again a cartesian square (with total space
      the open \((f')^{-1}(U)\) of \(X'\));
    \item[(b)] \(f_*\) of a restricted morphism is the restriction of \(f_*\) of the
      morphism, because pushforward commutes with restriction to an open of the base:
      for an open \(U\) of the base one has
      \((f_*\mathcal{G})|_U = (f|_{f^{-1}U})_*\bigl(\mathcal{G}|_{f^{-1}U}\bigr)\),
      naturally in \(\mathcal{G}\) — and likewise for \((g')_*\) and \(f'_*\);
    \item[(c)] the \((g^*, g_*)\)-adjunction transpose is a bijection natural in both
      variables, so it commutes with restricting its source and target along \(U\).
  \end{itemize}
  Chaining (a), (b), (c) — i.e.\ transporting the transpose-of-unit composite of
  Step 1 through the restriction functor \(\rho_U\) and replacing each block by its
  restricted counterpart — shows that the restriction of the base-change map to
  sections over \(U\) is itself the base-change map formed from the cartesian square
  restricted over \(U\). Choosing an affine open
  \(\operatorname{Spec}(R) \subseteq S\) containing the image \(g(U)\), affineness of
  \(f\) (local on the base and stable under base change) makes the total space
  restrict to an affine \(\operatorname{Spec}(A)\), and quasi-coherence of
  \(\mathcal{F}\) makes it restrict to a tilde-module \(\widetilde{M}\) there; so this
  restricted map is exactly the base-change map of the affine--affine square cut out
  over \(U = \operatorname{Spec}(R')\) and \(\operatorname{Spec}(R)\). The whole
  identification is thus definitional plus the naturalities (a)--(c); it introduces no
  obligation beyond them.

  \medskip
  \emph{Step 3: conclude the reduction.} By Step 2 the per-affine-open hypothesis
  required by the locality criterion Lemma~\ref{lem:modules_isIso_iff_affineOpens} —
  that the base-change map restricts to an isomorphism on sections over each affine
  open \(U \subseteq S'\) — is, under the identification of Step 2, precisely the
  affine--affine section assertion. That assertion is supplied, with \emph{no} flatness
  hypothesis, by the section-level computation
  Lemma~\ref{lem:pushforward_base_change_mate_cancelBaseChange}, which identifies the
  affine--affine base-change map on global sections with the cancellation isomorphism
  \(\operatorname{cancelBaseChange}\). Applying the locality criterion
  Lemma~\ref{lem:modules_isIso_iff_affineOpens} to this family of per-open
  isomorphisms shows that the base-change map is an isomorphism, as asserted.
\end{proof}

\subsection{The section-level base-change identity}

The section-level identification
Lemma~\ref{lem:pushforward_base_change_mate_cancelBaseChange} below is the single
highest-leverage obligation of the affine close. In the affine--affine model the
base-change map \(\theta\) is, after its domain and codomain are read through the two
proved tilde dictionaries, a concrete \(R'\)-linear map
\(R' \otimes_R M \to (R' \otimes_R A) \otimes_A M\), and the content of the lemma is
that this map is the inverse of the cancellation isomorphism
\(\operatorname{cancelBaseChange}\). Because the base-change map is \emph{defined} as the
transpose of the inner pushforward comparison \(\theta_{\mathrm{in}}\) under the
\((g^* \dashv g_*)\)-adjunction (\(g = \operatorname{Spec}\psi\)), its section-level value
is read off by the counit form of the hom-set adjunction: the transpose
factors \(\theta\) as \(g^*(\theta_{\mathrm{in}})\) followed by the geometric counit
\(\varepsilon_g\), and conjugating this composite by the two tilde dictionaries
\(\Theta_{\mathrm{src}}\) (Lemma~\ref{lem:base_change_mate_domain_read}) and
\(\Theta_{\mathrm{tgt}}\) (Lemma~\ref{lem:base_change_mate_codomain_read}) replaces the
geometric counit by the \emph{algebraic} extension-of-scalars counit along \(\psi\). Two Mathlib results are recorded
first as dependency anchors.

\begin{lemma}[Pullback of a pair of affine \(\operatorname{Spec}\)-maps]
  \label{lem:pullbackSpecIso_mathlib}
  \lean{AlgebraicGeometry.pullbackSpecIso}
  \mathlibok
  \textit{Provided by Mathlib.}
  Let \(R \to S\) and \(R \to T\) be \(R\)-algebras. Then the fibre product of the
  induced morphisms \(\operatorname{Spec}(S) \to \operatorname{Spec}(R)\) and
  \(\operatorname{Spec}(T) \to \operatorname{Spec}(R)\) is canonically the spectrum of
  the tensor product,
  \[
    \operatorname{Limits.pullback}\bigl(\operatorname{Spec}(R\to S),\,
      \operatorname{Spec}(R\to T)\bigr)
    \;\cong\; \operatorname{Spec}\bigl(S \otimes_R T\bigr)
    \qquad (\texttt{AlgebraicGeometry.pullbackSpecIso}),
  \]
  and under this isomorphism the two cone legs are the \(\operatorname{Spec}\)-maps of
  the canonical tensor inclusions: \texttt{pullbackSpecIso\_hom\_fst} identifies the
  first projection with \(\operatorname{Spec}\) of
  \(\texttt{Algebra.TensorProduct.includeLeftRingHom} : S \to S \otimes_R T\), and
  \texttt{pullbackSpecIso\_hom\_snd} identifies the second projection with
  \(\operatorname{Spec}\) of
  \(\texttt{Algebra.TensorProduct.includeRight} : T \to S \otimes_R T\) (with companion
  forms \texttt{pullbackSpecIso\_inv\_fst}, \texttt{\_inv\_snd}).
\end{lemma}

\begin{lemma}[Cancellation of base change for the tensor tower]
  \label{lem:cancelBaseChange_mathlib}
  \lean{TensorProduct.AlgebraTensorModule.cancelBaseChange}
  \mathlibok
  \textit{Provided by Mathlib.}
  Let \(R \to A \to B\) be a tower of algebras, \(M\) a \(B\)-module that is also an
  \(A\)- and \(R\)-module compatibly, and \(N\) an \(R\)-module. Then there is a
  \(B\)-linear equivalence, carrying \emph{no} flatness hypothesis,
  \[
    \operatorname{cancelBaseChange} :
    M \otimes_A (A \otimes_R N) \;\xrightarrow{\ \sim\ }\; M \otimes_R N,
  \]
  with \(\operatorname{cancelBaseChange}\bigl(m \otimes_A (a \otimes_R n)\bigr)
  = (a \cdot m) \otimes_R n\) (\texttt{cancelBaseChange\_tmul}) and inverse
  \(\operatorname{cancelBaseChange}^{-1}(m \otimes_R n) = m \otimes_A (1 \otimes_R n)\)
  (\texttt{cancelBaseChange\_symm\_tmul}). Instantiated at the affine--affine model
  (\(M = R'\), \(B = R'\), \(N = M\)) this is the cancellation isomorphism
  \((R' \otimes_R A) \otimes_A M \cong R' \otimes_R M\) used below, read in either
  orientation.
\end{lemma}


\begin{lemma}




\leanok
[The pullback cone legs are the \(\operatorname{Spec}\)-maps of the tensor inclusions]
  \label{lem:pullback_fst_snd_specMap_tensor}
  \lean{AlgebraicGeometry.pullback_fst_snd_specMap_tensor}
  \uses{lem:pullbackSpecIso_mathlib}
  Let \(\psi : R \to R'\) and \(\varphi : R \to A\) be homomorphisms of commutative
  rings, and regard \(R'\) and \(A\) as \(R\)-algebras via \(\psi\) and \(\varphi\).
  The canonical isomorphism of Lemma~\ref{lem:pullbackSpecIso_mathlib} identifies the
  fiber product \(\operatorname{Spec}(A) \times_{\operatorname{Spec}(R)}
  \operatorname{Spec}(R')\) with \(\operatorname{Spec}(R' \otimes_R A)\). Under this
  identification the two cone legs are:
  \begin{align*}
    g' &\;=\; \operatorname{Spec}\bigl(A \to R' \otimes_R A,\; a \mapsto 1 \otimes a\bigr),\\
    f' &\;=\; \operatorname{Spec}\bigl(R' \to R' \otimes_R A,\; r' \mapsto r' \otimes 1\bigr),
  \end{align*}
  composed with the inverse of the canonical isomorphism.
\end{lemma}

\begin{proof}
  \uses{lem:pullbackSpecIso_mathlib}
  The identification of each cone leg follows from the companion identities of
  Lemma~\ref{lem:pullbackSpecIso_mathlib}. The only content beyond citing those
  identities is to match the given morphisms \(\operatorname{Spec}\varphi\) and
  \(\operatorname{Spec}\psi\) with the \(\operatorname{Spec}\)-maps of the algebra
  structures that \(\varphi\) and \(\psi\) define on \(A\) and \(R'\) — a definitional
  identification — after which the Mathlib identities apply verbatim. This identification
  of the cone legs is the one Mathlib-absent step of the section-level computation.
\end{proof}

\begin{lemma}




\leanok
[Domain read of the section-level base-change map]
  \label{lem:base_change_mate_domain_read}
  \lean{AlgebraicGeometry.base_change_mate_domain_read}
  \uses{lem:pushforward_spec_tilde_iso, lem:pullback_spec_tilde_iso}
  % SOURCE: [Stacks Project], Cohomology of Schemes, Lemma
  % `lemma-affine-base-change' (\S\,Cohomology and base change, I), proof, the
  % "use Schemes, Lemma widetilde-pullback to describe pullbacks and pushforwards"
  % step (read from references/stacks-coherent.tex, L927--938).
  % SOURCE QUOTE: "We use Schemes, Lemma \ref{schemes-lemma-widetilde-pullback}
  % to describe pullbacks and pushforwards of $\mathcal{F}$."
  \textit{Source: Stacks Project, Cohomology of Schemes, Lemma ``Affine base
  change'' (proof, ``describe pullbacks and pushforwards'' step).}
  In the affine--affine model (\(g = \operatorname{Spec}\psi\),
  \(f = \operatorname{Spec}\varphi\), \(\mathcal{F} = \widetilde M\) for an
  \(A\)-module \(M\)), the global sections of the domain \(g^*(f_*\widetilde M)\) of the
  base-change map are canonically \(R' \otimes_R M\) as an \(R'\)-module. Explicitly,
  the assembled isomorphism
  \[
    \Theta_{\mathrm{src}} :
    \Gamma\bigl(g^*(f_*\widetilde M)\bigr) \xrightarrow{\ \sim\ } R' \otimes_R M
  \]
  is the composite of the pushforward dictionary
  Lemma~\ref{lem:pushforward_spec_tilde_iso} (reading \(f_*\widetilde M\) as
  restriction of scalars \(\operatorname{restr}_\varphi M\)) followed by the pullback
  dictionary Lemma~\ref{lem:pullback_spec_tilde_iso} (reading \(g^*\) of a tilde as
  extension of scalars \(R' \otimes_R (-)\) along \(\psi\)).
\end{lemma}

\begin{proof}
  \uses{lem:pushforward_spec_tilde_iso, lem:pullback_spec_tilde_iso}
  The domain \(g^*(f_*\widetilde M)\) involves only the cospan maps
  \(f = \operatorname{Spec}\varphi\) and \(g = \operatorname{Spec}\psi\), each a genuine
  \(\operatorname{Spec}\)-map, so the two affine dictionaries apply directly with no
  pullback-leg identification needed. The pushforward dictionary
  Lemma~\ref{lem:pushforward_spec_tilde_iso} gives
  \(f_*\widetilde M = \widetilde{\operatorname{restr}_\varphi M}\) on global sections,
  i.e.\ \(M\) viewed as an \(R\)-module; the pullback dictionary
  Lemma~\ref{lem:pullback_spec_tilde_iso} gives
  \(g^*\widetilde N = \widetilde{R' \otimes_R N}\), so
  \(\Gamma\bigl(g^*(f_*\widetilde M)\bigr) = R' \otimes_R M\) as an \(R'\)-module.
  Composing the two dictionary isomorphisms in this order is \(\Theta_{\mathrm{src}}\).
\end{proof}

\begin{lemma}




\leanok
[Codomain read of the section-level base-change map]
  \label{lem:base_change_mate_codomain_read}
  \lean{AlgebraicGeometry.base_change_mate_codomain_read}
  \uses{lem:pullback_fst_snd_specMap_tensor, lem:pullback_spec_tilde_iso,
        lem:pushforward_spec_tilde_iso, lem:pullback_isEquivalence_of_iso}
  % SOURCE: [Stacks Project], Cohomology of Schemes, Lemma
  % `lemma-affine-base-change' (\S\,Cohomology and base change, I), proof, the
  % "Namely X' = Spec(R' \otimes_R A) and F' is the qcoh sheaf associated to
  % (R' \otimes_R A) \otimes_A M" step (read from references/stacks-coherent.tex,
  % L929--932).
  % SOURCE QUOTE: "Namely, $X' = \Spec(R' \otimes_R A)$ and
  % $\mathcal{F}'$ is the quasi-coherent sheaf associated
  % to $(R' \otimes_R A) \otimes_A M$."
  \textit{Source: Stacks Project, Cohomology of Schemes, Lemma ``Affine base
  change'' (proof, ``\(X' = \operatorname{Spec}(R' \otimes_R A)\)'' step).}
  In the affine--affine model, with \(f'\) and \(g'\) the canonical projections of the
  pullback square (Lemma~\ref{lem:pullback_fst_snd_specMap_tensor}), the global
  sections of the codomain \(f'_*(g')^*\widetilde M\) are canonically
  \((R' \otimes_R A) \otimes_A M\) as an \(R'\)-module. Explicitly, the assembled
  isomorphism
  \[
    \Theta_{\mathrm{tgt}} :
    \Gamma\bigl(f'_*(g')^*\widetilde M\bigr)
    \xrightarrow{\ \sim\ } (R' \otimes_R A) \otimes_A M
  \]
  is the composite of the pullback dictionary Lemma~\ref{lem:pullback_spec_tilde_iso}
  (for \((g')^*\)) followed by the pushforward dictionary
  Lemma~\ref{lem:pushforward_spec_tilde_iso} (for \(f'_*\)).
\end{lemma}

\begin{proof}
  \uses{lem:pullback_fst_snd_specMap_tensor, lem:pullback_spec_tilde_iso,
        lem:pushforward_spec_tilde_iso}
  Here the two functors \((g')^*\) and \(f'_*\) are taken along the pullback cone legs,
  which are \emph{not} a priori \(\operatorname{Spec}\)-maps; this is where
  Lemma~\ref{lem:pullback_fst_snd_specMap_tensor} enters. By
  Lemma~\ref{lem:pullback_fst_snd_specMap_tensor}, the first projection \(g'\) is the
  \(\operatorname{Spec}\)-map of the inclusion \(A \to R' \otimes_R A\), and the second
  projection \(f'\) is the \(\operatorname{Spec}\)-map of \(R' \to R' \otimes_R A\).
  With the legs so identified the affine dictionaries apply:
  the pullback dictionary Lemma~\ref{lem:pullback_spec_tilde_iso} along
  \(A \to R' \otimes_R A\) reads \((g')^*\widetilde M\) as the extension of scalars
  \((R' \otimes_R A) \otimes_A M\), and the pushforward dictionary
  Lemma~\ref{lem:pushforward_spec_tilde_iso} along \(R' \to R' \otimes_R A\) reads
  \(f'_*\) of it as restriction of scalars back to \(R'\). The result is
  \((R' \otimes_R A) \otimes_A M\) as an \(R'\)-module, and composing the two dictionary
  isomorphisms is \(\Theta_{\mathrm{tgt}}\). This matches the Stacks description
  \(X' = \operatorname{Spec}(R' \otimes_R A)\), with \(\mathcal{F}'\) the tilde of
  \((R' \otimes_R A) \otimes_A M\).
\end{proof}


% --- Supporting lemmas for the codomain read ---
% These helpers establish that pullback along a scheme isomorphism is an equivalence
% of module categories; the adjunction unit of an equivalence is an isomorphism,
% used in the codomain-read computation.

\begin{lemma}\leanok
[Pullback along an isomorphism is an equivalence of module categories]
  \label{lem:pullbackIsoEquivalenceOfIso}
  \lean{AlgebraicGeometry.pullbackIsoEquivalenceOfIso}
  Let \(f : X \to Y\) be an isomorphism of schemes. Then the pullback functor on
  \(\mathcal{O}\)-modules \(f^* : Y\text{-}\mathbf{Mod} \to X\text{-}\mathbf{Mod}\) is
  one half of an equivalence of categories
  \(Y\text{-}\mathbf{Mod} \simeq X\text{-}\mathbf{Mod}\), with quasi-inverse
  \((f^{-1})^*\).
\end{lemma}
\begin{proof}

  Assemble the equivalence from the functor \(f^*\) and the candidate inverse
  \((f^{-1})^*\). The unit and counit isomorphisms are built from the pseudofunctor
  coherences of pullback: \((f^{-1})^* \circ f^* \cong (f \circ f^{-1})^* = \mathrm{id}^*
  \cong \mathrm{id}\) and symmetrically, using
  \(f \circ f^{-1} = \mathrm{id}_Y\), \(f^{-1} \circ f = \mathrm{id}_X\)
  (composition-of-pullbacks, congruence under equal maps, and pullback-of-identity).
  The triangle identities are forced by these coherences.
\end{proof}

\begin{lemma}\leanok
[The pullback functor along an isomorphism is an equivalence]
  \label{lem:pullback_isEquivalence_of_iso}
  \lean{AlgebraicGeometry.pullback_isEquivalence_of_iso}
  \uses{lem:pullbackIsoEquivalenceOfIso}
  Consequently, for an isomorphism \(f : X \to Y\) the functor
  \(f^* : Y\text{-}\mathbf{Mod} \to X\text{-}\mathbf{Mod}\) is an equivalence of
  categories. In particular its adjunction unit is a natural isomorphism, so
  \((\eta_{f})_{W}\) is invertible on every object \(W\).
\end{lemma}
\begin{proof}

  The functor underlying the equivalence of
  Lemma~\ref{lem:pullbackIsoEquivalenceOfIso} is exactly \(f^*\), so \(f^*\) is an
  equivalence; an equivalence's unit is a natural isomorphism.
\end{proof}

\subsection{Tilde-transport direct proof of the section-level value}
\label{sec:fbc_tilde_transport_direct}

% This subsection gives a direct affine computation of the section-level value of
% the base-change map. It routes through the single-step tilde dictionaries
% (Lemma~\ref{lem:pullback_spec_tilde_iso}, Lemma~\ref{lem:pushforward_spec_tilde_iso})
% and the tensor-product universal property.
\begin{lemma}
[Direct section-level value of the affine base-change map]
  \label{lem:pushforward_base_change_mate_sections_direct}
  \uses{def:pushforward_base_change_map,
        lem:base_change_mate_domain_read,
        lem:base_change_mate_codomain_read,
        lem:pullback_spec_tilde_iso,
        lem:pushforward_spec_tilde_iso,
        lem:cancelBaseChange_mathlib}
  In the affine--affine model — \(S = \operatorname{Spec}(R)\),
  \(S' = \operatorname{Spec}(R')\), \(X = \operatorname{Spec}(A)\),
  \(\mathcal{F} = \widetilde M\) for an \(A\)-module \(M\), \(g = \operatorname{Spec}\psi\)
  for \(\psi : R \to R'\), \(f = \operatorname{Spec}\varphi\) for \(\varphi : R \to A\) —
  the composition
  \[
    \Theta_{\mathrm{src}}^{-1} \;\ggg\; \Gamma(\alpha) \;\ggg\; \Theta_{\mathrm{tgt}}
    \;:\; R' \otimes_R M
    \;\longrightarrow\; (R' \otimes_R A) \otimes_A M
  \]
  equals the inverse cancellation isomorphism
  \(\operatorname{cancelBaseChange}^{-1} : r' \otimes m \mapsto (r' \otimes 1) \otimes m\)
  of Lemma~\ref{lem:cancelBaseChange_mathlib}, where \(\alpha\) is the base-change map
  (Definition~\ref{def:pushforward_base_change_map}) and \(\Theta_{\mathrm{src}}\),
  \(\Theta_{\mathrm{tgt}}\) are the domain and codomain identifications of
  Lemmas~\ref{lem:base_change_mate_domain_read}
  and~\ref{lem:base_change_mate_codomain_read}.
\end{lemma}

% SOURCE: [Stacks Project], Cohomology of Schemes, Lemma
% `lemma-affine-base-change' (\S\,Cohomology and base change, I), proof, the
% "boils down to the equality" step (read from references/stacks-coherent.tex,
% L927--938).
% SOURCE QUOTE PROOF: "We use Schemes, Lemma \ref{schemes-lemma-widetilde-pullback}
% to describe pullbacks and pushforwards of $\mathcal{F}$.
% Namely, $X' = \Spec(R' \otimes_R A)$ and
% $\mathcal{F}'$ is the quasi-coherent sheaf associated
% to $(R' \otimes_R A) \otimes_A M$.
% Thus we see that the lemma boils down to the
% equality
% $$
% (R' \otimes_R A) \otimes_A M = R' \otimes_R M
% $$
% as $R'$-modules."
\begin{proof}
  \uses{def:pushforward_base_change_map,
        lem:base_change_mate_domain_read,
        lem:base_change_mate_codomain_read,
        lem:pullback_spec_tilde_iso,
        lem:pushforward_spec_tilde_iso,
        lem:cancelBaseChange_mathlib}
  \textit{Source: Stacks Project, Cohomology of Schemes, Lemma ``Affine base
  change'' (proof, ``boils down to the equality'' step).}

  The proof is a direct module computation: we unfold the \emph{definition} of
  \(\alpha = \mathtt{pushforwardBaseChangeMap}\) and evaluate each adjunction
  unit/counit through the single-step tilde dictionaries, never invoking the
  composite-mate recognition. By
  Definition~\ref{def:pushforward_base_change_map}, \(\alpha\) is the
  \((g^* \dashv g_*)\)-adjunct of the map
  \[
    \beta : f_*\widetilde M \longrightarrow g_*\,f'_*(g')^*\widetilde M
          = f_*\,(g')_*(g')^*\widetilde M
  \]
  obtained by applying \(f_*\) to the unit
  \(\eta : \widetilde M \to (g')_*(g')^*\widetilde M\) of the
  \(((g')^*, (g')_*)\)-adjunction, using the square's commutativity
  \(g \circ f' = f \circ g'\). Explicitly, as adjunct,
  \(\alpha = \varepsilon_{f'_*(g')^*\widetilde M} \circ g^*(\beta)\), where
  \(\varepsilon\) is the counit of \(g^* \dashv g_*\). We trace each factor on global
  sections over the affine model, on a generator \(r' \otimes m \in R' \otimes_R M\).

  \emph{Step 1 — the inner unit \(\eta\) is the extension-of-scalars unit.}
  On \(X = \operatorname{Spec}(A)\), the projection \(g'\) is the
  \(\operatorname{Spec}\)-map of the algebra inclusion \(A \to R' \otimes_R A\)
  (Lemma~\ref{lem:base_change_mate_codomain_read}). By the pullback dictionary
  Lemma~\ref{lem:pullback_spec_tilde_iso}, \((g')^*\widetilde M\) is the tilde of the
  extension of scalars \((R' \otimes_R A) \otimes_A M\), and that lemma identifies the
  adjunction unit with the algebraic unit of extension of scalars; hence on sections
  \(\eta\) is the \(A\)-linear map
  \[
    M \longrightarrow (R' \otimes_R A) \otimes_A M, \qquad
    m \longmapsto (1 \otimes 1) \otimes m .
  \]
  This is forced by the universal property of the tensor product and the algebra map
  \(A \to R' \otimes_R A\); no mate identity is used.

  \emph{Step 2 — apply \(f_* = \operatorname{restr}_\varphi\).} The pushforward
  dictionary Lemma~\ref{lem:pushforward_spec_tilde_iso} reads \(f_*\) of a tilde as
  restriction of scalars along \(\varphi\). Restriction of scalars does not move
  elements, so on sections over \(S = \operatorname{Spec}(R)\) the map
  \(\beta = f_*(\eta)\) (with codomain rewritten by \(g \circ f' = f \circ g'\)) is the
  same assignment \(m \mapsto (1 \otimes 1) \otimes m\), now read \(R\)-linearly into
  \((R' \otimes_R A) \otimes_A M\).

  \emph{Step 3 — apply \(g^* = R' \otimes_R (-)\).} The pullback dictionary
  Lemma~\ref{lem:pullback_spec_tilde_iso} for \(g = \operatorname{Spec}\psi\) reads
  \(g^*\) of a tilde as extension of scalars \(R' \otimes_R (-)\) along \(\psi\), and is
  functorial: \(g^*(\beta)\) on sections over \(S' = \operatorname{Spec}(R')\) is
  \(\operatorname{id}_{R'} \otimes \beta\),
  \[
    R' \otimes_R M \longrightarrow R' \otimes_R \bigl((R' \otimes_R A) \otimes_A M\bigr),
    \qquad
    r' \otimes m \longmapsto r' \otimes \bigl((1 \otimes 1) \otimes m\bigr).
  \]
  The domain \(R' \otimes_R M\) here is exactly \(\Theta_{\mathrm{src}}\)'s reading of
  \(\Gamma(g^*(f_*\widetilde M))\) (Lemma~\ref{lem:base_change_mate_domain_read}).

  \emph{Step 4 — the counit \(\varepsilon\) is the extension-of-scalars counit.}
  Writing \(W = (R' \otimes_R A) \otimes_A M\) for the target \(R'\)-module, the counit
  \(\varepsilon : g^* g_* \widetilde W \to \widetilde W\) of \(g^* \dashv g_*\) is, on
  sections over \(\operatorname{Spec}(R')\) and again by
  Lemma~\ref{lem:pullback_spec_tilde_iso} (which identifies the counit with the
  algebraic action map of extension of scalars), the \(R'\)-linear map
  \[
    R' \otimes_R \bigl(\operatorname{restr}_\psi W\bigr) \longrightarrow W,
    \qquad r' \otimes w \longmapsto r' \cdot w ,
  \]
  where \(r' \cdot w\) is the \(R'\)-action on \(W\). The codomain \(W\) is exactly
  \(\Theta_{\mathrm{tgt}}\)'s reading of \(\Gamma(f'_*(g')^*\widetilde M)\)
  (Lemma~\ref{lem:base_change_mate_codomain_read}).

  \emph{Composition.} Composing Steps 3 and 4, the conjugated section map sends
  \[
    r' \otimes m
    \;\overset{g^*(\beta)}{\longmapsto}\;
    r' \otimes \bigl((1 \otimes 1) \otimes m\bigr)
    \;\overset{\varepsilon}{\longmapsto}\;
    r' \cdot \bigl((1 \otimes 1) \otimes m\bigr)
    = (r' \otimes 1) \otimes m ,
  \]
  the last equality because the \(R'\)-action on \((R' \otimes_R A) \otimes_A M\)
  multiplies into the \(R'\) tensor slot: \(r' \cdot ((1 \otimes 1) \otimes m)
  = ((r' \cdot 1) \otimes 1) \otimes m = (r' \otimes 1) \otimes m\). Thus
  \(\Theta_{\mathrm{src}}^{-1} \ggg \Gamma(\alpha) \ggg \Theta_{\mathrm{tgt}}\) is the
  \(R'\)-linear map \(r' \otimes m \mapsto (r' \otimes 1) \otimes m\), which is
  precisely \(\operatorname{cancelBaseChange}^{-1}\) of
  Lemma~\ref{lem:cancelBaseChange_mathlib}
  (\(\mathtt{cancelBaseChange\_symm\_tmul}\)). No flatness hypothesis is used, and the
  value falls out of the tensor-product universal property together with the algebra
  map \(\varphi : R \to A\) — not from any conjugate/mate coherence identity.


\end{proof}

\begin{lemma}



\leanok
[Section-level value of the base-change map equals the cancellation isomorphism]
  \label{lem:pushforward_base_change_mate_cancelBaseChange}
  \lean{AlgebraicGeometry.pushforward_base_change_mate_cancelBaseChange}
  % NOTE: The isomorphism conclusion is stated directly. Stating the literal equality
  % \(\Theta_{\mathrm{tgt}} \circ \Gamma(\alpha) \circ
  % \Theta_{\mathrm{src}}^{-1} = \mathrm{cancelBaseChange}^{-1}\) would first require
  % identifying the pullback projections over the generic square with the Spec-of-tensor
  % inclusions. The proof instead assembles the isomorphism via the tilde-transport route:
  % the domain and codomain reads
  % (\cref{lem:base_change_mate_domain_read}, \cref{lem:base_change_mate_codomain_read})
  % identify the endpoints, and \cref{lem:pushforward_base_change_mate_sections_direct}
  % supplies the explicit conjugated value (no flatness hypothesis needed), hence
  % \(\Gamma(\alpha)\) is an isomorphism.
  \uses{def:pushforward_base_change_map, lem:base_change_mate_domain_read,
  lem:base_change_mate_codomain_read,
  lem:cancelBaseChange_mathlib}
  % SOURCE: [Stacks Project], Cohomology of Schemes, Lemma
  % `lemma-affine-base-change' (\S\,Cohomology and base change, I), proof, the
  % "boils down to the equality" step (read from references/stacks-coherent.tex,
  % L927--938).
  % SOURCE QUOTE: "We use Schemes, Lemma \ref{schemes-lemma-widetilde-pullback}
  % to describe pullbacks and pushforwards of $\mathcal{F}$.
  % Namely, $X' = \Spec(R' \otimes_R A)$ and
  % $\mathcal{F}'$ is the quasi-coherent sheaf associated
  % to $(R' \otimes_R A) \otimes_A M$.
  % Thus we see that the lemma boils down to the
  % equality
  % $$
  % (R' \otimes_R A) \otimes_A M = R' \otimes_R M
  % $$
  % as $R'$-modules."
  \textit{Source: Stacks Project, Cohomology of Schemes, Lemma ``Affine base
  change'' (proof, ``boils down to the equality'' step).}
  In the affine--affine case \(S = \operatorname{Spec}(R)\),
  \(S' = \operatorname{Spec}(R')\), \(X = \operatorname{Spec}(A)\),
  \(\mathcal{F} = \widetilde{M}\) — with \(g = \operatorname{Spec}(\psi)\) for
  \(\psi : R \to R'\), \(f = \operatorname{Spec}(\varphi)\) for \(\varphi : R \to A\),
  and \(M\) an \(A\)-module — write \(\Gamma(\alpha)\) for the global-sections
  incarnation of the base-change map (Definition~\ref{def:pushforward_base_change_map}).
  Recall \(\alpha\) runs
  \(g^*(f_*\widetilde{M}) \to f'_*(g')^*\widetilde{M}\). Transport \(\Gamma(\alpha)\)
  through the two proved affine dictionaries that identify its domain and codomain: the
  pushforward dictionary Lemma~\ref{lem:pushforward_spec_tilde_iso}, followed by the
  pullback dictionary Lemma~\ref{lem:pullback_spec_tilde_iso}, identifies the domain
  \(\Gamma\bigl(g^*(f_*\widetilde{M})\bigr)\) with \(R' \otimes_R M\); the pullback
  dictionary Lemma~\ref{lem:pullback_spec_tilde_iso}, followed by the pushforward
  dictionary Lemma~\ref{lem:pushforward_spec_tilde_iso}, identifies the codomain
  \(\Gamma\bigl(f'_*(g')^*\widetilde{M}\bigr)\) with \((R' \otimes_R A) \otimes_A M\),
  both read as \(R'\)-modules. Under these identifications \(\Gamma(\alpha)\) is
  precisely the canonical cancellation isomorphism
  \[
    \operatorname{cancelBaseChange} :
    (R' \otimes_R A) \otimes_A M \;\xrightarrow{\ \sim\ }\; R' \otimes_R M,
    \qquad (r' \otimes a) \otimes m \;\mapsto\; r' \otimes (a \cdot m)
  \]
  (Lemma~\ref{lem:cancelBaseChange_mathlib}; here \(a \cdot m\) is the \(A\)-action on \(M\)). Equivalently, in the direction in which
  \(\Gamma(\alpha)\) is presented — from the source \(R' \otimes_R M\) to the target
  \((R' \otimes_R A) \otimes_A M\) — it is \(\operatorname{cancelBaseChange}^{-1}\),
  \(r' \otimes m \mapsto (r' \otimes 1) \otimes m\); the two readings agree because
  \(\operatorname{cancelBaseChange}\) is an isomorphism. Since
  \(\operatorname{cancelBaseChange}\) is an isomorphism with \emph{no} flatness
  hypothesis, this identification closes the affine case: \(\Gamma(\alpha)\) is an
  isomorphism, hence (full faithfulness / counit isomorphism of \(\widetilde{(-)}\) on
  quasi-coherent modules) so is \(\alpha\) on the affine chart.

\end{lemma}

\begin{proof}
  \uses{def:pushforward_base_change_map, lem:base_change_mate_domain_read,
        lem:base_change_mate_codomain_read,
        lem:cancelBaseChange_mathlib}
  By Lemma~\ref{lem:base_change_mate_domain_read} the
  source \(\Gamma\bigl(g^*(f_*\widetilde M)\bigr)\) is \(R' \otimes_R M\) via
  \(\Theta_{\mathrm{src}}\), and by Lemma~\ref{lem:base_change_mate_codomain_read} the
  target \(\Gamma\bigl(f'_*(g')^*\widetilde M\bigr)\) is \((R' \otimes_R A) \otimes_A M\)
  via \(\Theta_{\mathrm{tgt}}\). By the direct section-value computation
  Lemma~\ref{lem:pushforward_base_change_mate_sections_direct}, the conjugated map
  \[
    \theta = \Theta_{\mathrm{src}}^{-1} \ggg \Gamma(\alpha) \ggg \Theta_{\mathrm{tgt}}
  \]
  equals \(\operatorname{cancelBaseChange}^{-1}\), the inverse of the cancellation
  isomorphism of Lemma~\ref{lem:cancelBaseChange_mathlib}. Since
  \(\operatorname{cancelBaseChange}\) is an isomorphism with \emph{no} flatness
  hypothesis, so is its inverse \(\theta\). As \(\theta\) is the conjugate of
  \(\Gamma(\alpha)\) by the two dictionary isomorphisms
  \(\Theta_{\mathrm{src}}, \Theta_{\mathrm{tgt}}\), it follows that \(\Gamma(\alpha)\) is
  itself an isomorphism — assembled by writing
  \(\Gamma(\alpha) = \Theta_{\mathrm{src}} \ggg \theta \ggg \Theta_{\mathrm{tgt}}^{-1}\)
  and inferring the instance. Both the source and target of \(\alpha\) are
  quasi-coherent, so the tilde-equivalence \(\widetilde{(-)}\) is fully faithful on
  them and \(\mathrm{IsIso}\,\Gamma(\alpha)\) yields \(\mathrm{IsIso}\,\alpha\); the
  closing conclusion the affine close consumes is \(\mathrm{IsIso}\ \Gamma(\alpha)\).
  This routes entirely through the tilde dictionaries and aligns the
  section-level map with the Stacks ``boils down to the equality'' step.
\end{proof}

% NOTE: The flat base change theorem
% (\cref{thm:flat_base_change_pushforward}) is reached via the concrete module comparison
% route (\cref{lem:base_changed_equalizer_diagram}) together with
% \cref{lem:flat_base_change_reduce_global_sections}; the per-affine identification
% is \cref{lem:pullback_spec_tilde_iso}. This block and the mate apparatus are retained
% as proved supporting lemmas whose statements are referenced by downstream declarations.
\begin{lemma}\leanok
[Affine base change]
  \label{lem:affine_base_change_pushforward}
  \lean{AlgebraicGeometry.affineBaseChange_pushforward_iso}
  \uses{def:pushforward_base_change_map, lem:base_change_map_affine_local,
        lem:modules_isIso_iff_affineOpens, lem:pullback_spec_tilde_iso,
        lem:pushforward_base_change_mate_cancelBaseChange,
        lem:pushforward_spec_tilde_iso}
  % SOURCE: [Stacks Project], Cohomology of Schemes, Lemma
  % `lemma-affine-base-change' (\S\,Cohomology and base change, I)
  % (read from references/stacks-coherent.tex, L906--918).
  % SOURCE QUOTE: "Let $f : X \to S$ be a morphism of schemes.
  % Let $\mathcal{F}$ be a quasi-coherent $\mathcal{O}_X$-module.
  % Assume $f$ is affine.
  % In this case $f_*\mathcal{F} \cong Rf_*\mathcal{F}$ is
  % a quasi-coherent sheaf, and for every base change diagram
  % (\ref{equation-base-change-diagram})
  % we have
  % $$
  % g^*f_*\mathcal{F} = f'_*(g')^*\mathcal{F}.
  % $$"
  \textit{Source: Stacks Project, Cohomology of Schemes, Lemma ``Affine base
  change''.}
  Assume \(f : X \to S\) is an affine morphism and \(\mathcal{F}\) is a
  quasi-coherent \(\mathcal{O}_X\)-module. Then for every base-change square as
  above the canonical map of
  Definition~\ref{def:pushforward_base_change_map} is an isomorphism
  \[
    g^*\bigl(f_*\mathcal{F}\bigr) \;\cong\; f'_*\bigl((g')^*\mathcal{F}\bigr).
  \]
  The quasi-coherence hypothesis on \(\mathcal{F}\) is essential and is not a mere
  convenience: the proof works over an affine open by writing
  \(\mathcal{F} = \widetilde{M}\) for a module \(M\), and a general
  \(\mathcal{O}_X\)-module need not be of this form
  \(\widetilde{M}\); without it neither the affine description of the pushforward
  nor the tensor-associativity computation below is available, and the conclusion
  fails.
\end{lemma}

% SOURCE QUOTE PROOF: "The vanishing of higher direct images is
% Lemma \ref{lemma-relative-affine-vanishing}.
% The statement is local on $S$ and $S'$. Hence we may
% assume $X = \Spec(A)$, $S = \Spec(R)$,
% $S' = \Spec(R')$ and $\mathcal{F} = \widetilde{M}$
% for some $A$-module $M$.
% We use Schemes, Lemma \ref{schemes-lemma-widetilde-pullback}
% to describe pullbacks and pushforwards of $\mathcal{F}$.
% Namely, $X' = \Spec(R' \otimes_R A)$ and
% $\mathcal{F}'$ is the quasi-coherent sheaf associated
% to $(R' \otimes_R A) \otimes_A M$.
% Thus we see that the lemma boils down to the
% equality
% $$
% (R' \otimes_R A) \otimes_A M = R' \otimes_R M
% $$
% as $R'$-modules."
% NOTE: This affine lemma is closed directly on global sections, via the two named
% obligations: lem:base_change_map_affine_local (locality reduction on the affine
% opens of S') and lem:pushforward_base_change_mate_cancelBaseChange (the section-
% level value of the base-change map is cancelBaseChange). Both proved affine tilde
% dictionaries -- lem:pushforward_spec_tilde_iso and lem:pullback_spec_tilde_iso --
% feed the second obligation. cancelBaseChange is in Mathlib with no flatness; the
% affine case carries no flatness hypothesis. No Čech infrastructure is used here.
\begin{proof}
  \uses{def:pushforward_base_change_map, lem:modules_isIso_iff_affineOpens, lem:pushforward_spec_tilde_iso, lem:pullback_spec_tilde_iso, lem:base_change_map_affine_local, lem:pushforward_base_change_mate_cancelBaseChange}
  The proof is the composite of the two named obligations, taken \emph{directly on
  global sections}.

  \medskip
  \emph{First reduction (locality on \(S'\)).} By
  Lemma~\ref{lem:modules_isIso_iff_affineOpens} it suffices to prove that the
  base-change map of Definition~\ref{def:pushforward_base_change_map} is an
  isomorphism on sections over every affine open of \(S'\). This is exactly the
  locality reduction Lemma~\ref{lem:base_change_map_affine_local}: restricting the
  cartesian square to an affine open \(U = \operatorname{Spec}(R') \subseteq S'\) and
  a chosen affine \(\operatorname{Spec}(R) \subseteq S\) containing \(g(U)\), the
  hypothesis \([\mathrm{IsAffineHom}\ f]\) makes \(X\) restrict to
  \(\operatorname{Spec}(A)\) and quasi-coherence of \(\mathcal{F}\) makes it restrict
  to \(\widetilde{M}\) for an \(A\)-module \(M\). (It is here that quasi-coherence is
  used: it is what licenses writing \(\mathcal{F} = \widetilde{M}\).) Computing the
  fibre product gives \(X' = \operatorname{Spec}(R' \otimes_R A)\) and
  \(\mathcal{F}' = (g')^*\mathcal{F}\) the tilde of the \((R' \otimes_R A)\)-module
  \((R' \otimes_R A) \otimes_A M\). The matter is thus reduced to the affine--affine
  model, where the base-change map is a morphism \(\alpha\) of quasi-coherent
  \((\operatorname{Spec} R')\)-modules.

  \medskip
  \emph{Passage to global sections.} The functor
  \(\widetilde{(-)} : \operatorname{ModuleCat}(R') \to
  (\operatorname{Spec} R').\mathrm{Modules}\) is fully faithful with essential image
  the quasi-coherent modules; equivalently the counit
  \(\widetilde{\Gamma(\mathcal{N})} \to \mathcal{N}\) is an isomorphism exactly when
  \(\mathcal{N}\) is quasi-coherent. Source and target of \(\alpha\) are
  quasi-coherent, so both counits are isomorphisms and, by their naturality,
  \[
    \mathrm{IsIso}\,\alpha
    \;\Longleftrightarrow\;
    \mathrm{IsIso}\bigl(\Gamma(\alpha)\bigr),
  \]
  where \(\Gamma(\alpha)\) is a concrete \(R'\)-linear map in
  \(\operatorname{ModuleCat}(R')\). The whole computation therefore stays at the level
  of modules.

  \medskip
  \emph{Section-level identification (the second obligation).} By the section-level
  computation Lemma~\ref{lem:pushforward_base_change_mate_cancelBaseChange}, the map
  \(\Gamma(\alpha)\), with its domain \(\Gamma(g^*(f_*\widetilde{M}))\) identified with
  \(R' \otimes_R M\) and its codomain \(\Gamma(f'_*(g')^*\widetilde{M})\) identified
  with \((R' \otimes_R A) \otimes_A M\) through the two dictionaries
  Lemmas~\ref{lem:pullback_spec_tilde_iso} and~\ref{lem:pushforward_spec_tilde_iso},
  \emph{is} the canonical cancellation isomorphism
  \[
    \operatorname{cancelBaseChange} :
    (R' \otimes_R A) \otimes_A M \;\xrightarrow{\ \sim\ }\; R' \otimes_R M,
    \qquad (r' \otimes a) \otimes m \;\mapsto\; r' \otimes (a \cdot m)
  \]
  of \(R'\)-modules (Lemma~\ref{lem:cancelBaseChange_mathlib}). This is the entire
  content of the affine close, and it is performed \emph{directly on the module of
  global sections} — there is no separate abstract adjoint-mate identification at the
  level of sheaves. Because \(\operatorname{cancelBaseChange}\) is an isomorphism with
  no flatness hypothesis whatsoever, \(\Gamma(\alpha)\) is an isomorphism, hence so is
  \(\alpha\) on the affine chart, and by the locality reduction so is the global
  base-change map. Flatness of \(g\) plays no role in the affine lemma; it enters only
  in the global statement Theorem~\ref{thm:flat_base_change_pushforward}, where it is
  used to commute \(- \otimes_A B\) with the finite equalizer computing \(H^0\).
\end{proof}

\section{Flat base change for the pushforward}

The global statement is assembled from the affine lemma and the single homological
input that flatness preserves the finite equalizer computing \(H^0\). That input is
supplied verbatim by Mathlib, recorded here as a dependency anchor.

\begin{lemma}[Flat base change preserves finite equalizers]
  \label{lem:flat_preserves_equalizer_mathlib}
  \lean{LinearMap.tensorEqLocusEquiv}
  \mathlibok
  \textit{Provided by Mathlib.}
  Let \(A \to B\) be a flat ring homomorphism (more generally, let \(B\) be a flat
  \(A\)-module that is also an \(A\)-algebra, with the scalar-tower compatibility),
  and let \(f, g : N \rightrightarrows P\) be a pair of \(A\)-linear maps. Write
  \(\operatorname{eq}(f, g) = \{x \in N : f(x) = g(x)\}\) for their equalizer. Then the
  canonical \(B\)-linear comparison map
  \[
    B \otimes_A \operatorname{eq}(f, g)
    \;\xrightarrow{\ \sim\ }\;
    \operatorname{eq}\bigl(\mathbf{1}_B \otimes f,\ \mathbf{1}_B \otimes g\bigr)
  \]
  is an isomorphism; that is, base change along a flat algebra commutes with the
  formation of the finite equalizer (equivalently, of the kernel: the case
  \(g = 0\) is \texttt{LinearMap.tensorKerEquiv}, underlain by
  \texttt{Module.Flat.ker\_lTensor\_eq} and \texttt{Module.Flat.eqLocus\_lTensor\_eq}).
  Since \(- \otimes_A B\) also preserves finite products, it follows that
  \(- \otimes_A B\) carries the degree-\(0\)/\(1\) equalizer
  \(\prod_i \Gamma(U_i, \mathcal{F}) \rightrightarrows
    \prod_{i,j} \Gamma(U_{ij}, \mathcal{F})\) of a finite cover to the corresponding
  equalizer of the base-changed cover.
\end{lemma}

\begin{lemma}[Sheaf condition as an equalizer of products]
  \label{lem:sheaf_equalizer_products_mathlib}
  \lean{TopCat.Presheaf.isSheaf_iff_isSheafEqualizerProducts}
  \mathlibok
  \textit{Provided by Mathlib.}
  Let \(\mathcal{F}\) be a presheaf with values in a category with products on a
  topological space, and let \(\mathcal{U} = \{U_i\}_{i \in \iota}\) be a family of
  opens with union \(U = \bigcup_i U_i\). Then \(\mathcal{F}\) satisfies the sheaf
  condition on \(\mathcal{U}\) if and only if the fork
  \[
    \mathcal{F}(U)
    \;\xrightarrow{\;\mathrm{res}\;}\;
    \prod_i \mathcal{F}(U_i)
    \;\underset{\mathrm{rightRes}}{\overset{\mathrm{leftRes}}{\rightrightarrows}}\;
    \prod_{(i,j)} \mathcal{F}(U_i \cap U_j)
  \]
  is a limit (an equalizer), where the two parallel maps restrict the \(U_i\)- and
  \(U_j\)-components to \(U_i \cap U_j\). This is Mathlib's
  \texttt{isSheaf\_iff\_isSheafEqualizerProducts} together with the fork
  \texttt{SheafConditionEqualizerProducts.fork}; it is exactly Čech degree
  \(0 \to 1\) and requires no separatedness.
\end{lemma}

\begin{lemma}\leanok
[Finite affine cover with quasi-compact overlaps]
  \label{lem:finite_affine_cover_qcqs}
  \lean{AlgebraicGeometry.Scheme.exists_finite_affineCover_inter_isQuasiCompact}
  A quasi-compact scheme \(X\) admits a \emph{finite} affine open cover
  \(\mathcal{U} : X = U_1 \cup \dots \cup U_t\). If moreover \(X\) is
  quasi-separated, then each pairwise intersection \(U_{ij} = U_i \cap U_j\) is
  quasi-compact; and if \(X\) is separated, each \(U_{ij}\) is in fact affine.
  \begin{proof}
    The affine opens form a basis, so \(X\) is covered by affine opens; quasi-compactness
    extracts a finite subcover \(U_1, \dots, U_t\) (by Mathlib's affine cover and the
    finite-subcover property of a quasi-compact space). Quasi-separatedness says the
    intersection of two quasi-compact opens is quasi-compact, so each \(U_{ij}\) is
    quasi-compact. When \(X\) is separated the diagonal is a closed immersion, whence
    the intersection of two affine opens is affine (Mathlib, for separated schemes).
  \end{proof}
\end{lemma}

\begin{lemma}\leanok
[$\Gamma(X,\mathcal{F})$ as a finite sheaf-condition equalizer]
  \label{lem:gamma_finite_equalizer}
  \lean{AlgebraicGeometry.Modules.gammaIsLimitSheafConditionFork}
  % NOTE: The formal statement is strictly more general than this prose:
  % it takes any open family over any scheme without quasi-compactness, finiteness, or
  % affinity hypotheses, and produces an is-limit datum for the corresponding fork. The
  % finite affine cover of a quasi-compact quasi-separated scheme is the intended
  % application; the consolidated existential version has its own block,
  % Lemma~\ref{lem:gamma_finite_equalizer_cover}.
  \uses{lem:sheaf_equalizer_products_mathlib, lem:finite_affine_cover_qcqs}
  Let \(X\) be a quasi-compact, quasi-separated scheme and \(\mathcal{F}\) a sheaf of
  \(\mathcal{O}_X\)-modules. Fix a finite affine cover \(\mathcal{U} = \{U_i\}_{1 \le i \le t}\)
  of \(X\) as in Lemma~\ref{lem:finite_affine_cover_qcqs}. Then the module of global
  sections is the equalizer of the two restriction maps of the cover,
  \[
    \Gamma(X, \mathcal{F})
    \;=\;
    \operatorname{eq}\Bigl(
      \textstyle\prod_{i} \Gamma(U_i, \mathcal{F})
      \;\rightrightarrows\;
      \textstyle\prod_{i,j} \Gamma(U_{ij}, \mathcal{F})
    \Bigr),
    \qquad
    s \mapsto (s_i|_{U_{ij}}), \quad s \mapsto (s_j|_{U_{ij}}),
  \]
  and the products are \emph{finite} (indexed by the finite sets \(\{i\}\) and
  \(\{(i,j)\}\)).
  \begin{proof}
    \uses{lem:sheaf_equalizer_products_mathlib, lem:finite_affine_cover_qcqs}
    Apply the equalizer-products form of the sheaf condition
    (Lemma~\ref{lem:sheaf_equalizer_products_mathlib}) to the underlying sheaf of
    \(\mathcal{F}\) and the cover \(\mathcal{U}\): since \(\bigcup_i U_i = X\), the
    fork displayed there computes \(\mathcal{F}(X) = \Gamma(X, \mathcal{F})\) as the
    equalizer of \(\prod_i \Gamma(U_i, \mathcal{F}) \rightrightarrows \prod_{i,j}
    \Gamma(U_{ij}, \mathcal{F})\). The cover is finite by
    Lemma~\ref{lem:finite_affine_cover_qcqs}, so both products are finite. The
    module structure is inherited objectwise: the forgetful functor to abelian
    groups preserves and reflects the limit, so the equalizer is computed in
    \(\operatorname{ModuleCat}\) over the ground ring.
  \end{proof}
\end{lemma}



% ============================================================================
% Decomposition of lem:base_changed_equalizer_diagram into three components:
% (a) per-chart module iso, (b) restriction-naturality, and
% (c) tensor commutes with finite products. The proof below assembles these.
% ============================================================================

\begin{lemma}\leanok
[Per-chart base-change module isomorphism]
  \label{lem:base_change_chart_tensor_iso}
  \lean{AlgebraicGeometry.baseChange_chart_tensorIso}
  \uses{lem:pullback_spec_tilde_iso}
  % SOURCE: [Stacks Project], Schemes, Lemma `lemma-widetilde-pullback', Tag 01I9
  % (\S\,Quasi-coherent sheaves on affines) (read from references/stacks-schemes.tex,
  % L1242--1256).
  % SOURCE QUOTE: "Let
  % $(X, \mathcal{O}_X) = (\Spec(S), \mathcal{O}_{\Spec(S)})$,
  % $(Y, \mathcal{O}_Y) = (\Spec(R), \mathcal{O}_{\Spec(R)})$
  % be affine schemes.
  % Let $\psi : (X, \mathcal{O}_X) \to (Y, \mathcal{O}_Y)$ be a
  % morphism of affine schemes, corresponding to the ring map
  % $\psi^\sharp : R \to S$ (see Lemma \ref{lemma-category-affine-schemes}).
  % \begin{enumerate}
  % \item We have $\psi^* \widetilde M = \widetilde{S \otimes_R M}$
  % functorially in the $R$-module $M$.
  % \item We have $\psi_* \widetilde N = \widetilde{N_R}$ functorially
  % in the $S$-module $N$.
  % \end{enumerate}"
  \textit{Source: Stacks Project, Schemes, Lemma ``\(\widetilde M\) and pullback''
  (Tag 01I9), part (1), applied to the base-change ring map of one affine chart.}
  Let \(A \to B\) be a ring map, \(X\) a scheme over \(A\) (\(A = \Gamma(X,\mathcal{O}_X)\)
  the ground ring), \(\mathcal{F}\) a quasi-coherent \(\mathcal{O}_X\)-module, and
  \(X_B = X \times_{\operatorname{Spec} A}\operatorname{Spec} B\) with projection
  \(g' : X_B \to X\) and pulled-back module \(\mathcal{F}_B = (g')^{*}\mathcal{F}\). Let
  \(V \subseteq X\) be an \emph{affine} open, \(V = \operatorname{Spec} R\) with
  \(R = \Gamma(V,\mathcal{O}_X)\) and \(\Gamma(V,\mathcal{F}) = M\) the associated
  \(R\)-module (\(\mathcal{F}|_V = \widetilde M\)). Then \(V_B = g'^{-1}(V)\) is the affine
  \(\operatorname{Spec}(B \otimes_A R)\), and there is a canonical \(B\)-linear isomorphism
  \[
    \Gamma\bigl(V_B, \mathcal{F}_B\bigr)
    \;\cong\;
    \Gamma(V,\mathcal{F}) \otimes_A B .
  \]
  \begin{proof}
    \uses{lem:pullback_spec_tilde_iso}
    Base change preserves affineness, so \(V_B = V \times_{\operatorname{Spec} A}
    \operatorname{Spec} B = \operatorname{Spec}(B \otimes_A R)\) is affine, and the chart
    base-change map is \(\operatorname{Spec}\varphi\) for the ring map
    \(\varphi : R \to B \otimes_A R\). Restricting \(\mathcal{F}_B = (g')^{*}\mathcal{F}\)
    to \(V_B\) is pulling \(\mathcal{F}|_V = \widetilde M\) back along
    \(\operatorname{Spec}\varphi\), so the affine pullback dictionary
    Lemma~\ref{lem:pullback_spec_tilde_iso} (Stacks Tag 01I9, part~(1)) identifies it with
    the tilde of the base-changed module:
    \(\mathcal{F}_B|_{V_B} \cong \widetilde{(B \otimes_A R) \otimes_R M}\). Taking global
    sections over the affine \(V_B\) gives
    \(\Gamma(V_B,\mathcal{F}_B) \cong (B \otimes_A R)\otimes_R M\). Finally the cancellation
    of base changes \((B \otimes_A R)\otimes_R M \cong B \otimes_A M\) (the associativity
    isomorphism \(\operatorname{cancelBaseChange}\), provided by Mathlib, carries no
    flatness hypothesis) rewrites this as \(\Gamma(V,\mathcal{F}) \otimes_A B\). Every step
    is affine: no abstract base-change map, no flatness, and no adjoint-mate identification
    enters.
  \end{proof}
\end{lemma}

% --- (b)-scaffolding helpers ---

\begin{definition}\leanok
[Chart restriction as an algebra map]
  \label{def:chartRestrictionAlgHom}
  \lean{AlgebraicGeometry.chartRestrictionAlgHom}
  The restriction ring map \(\rho : R = \Gamma(V,\mathcal{O}_X)\to R' = \Gamma(W,\mathcal{O}_X)\) of an
  affine open immersion \(j : W\hookrightarrow V\), repackaged as an \(A\)-algebra homomorphism
  \(R\to R'\) via the chart algebra structures.
\end{definition}

\begin{definition}\leanok
[Base change of the chart restriction map]
  \label{def:chartBaseChangeRingMap}
  \lean{AlgebraicGeometry.chartBaseChangeRingMap}
  \uses{def:chartRestrictionAlgHom}
  The base-changed ring map \(\rho_B = \rho\otimes_A\mathrm{id}_B : R\otimes_A B\to R'\otimes_A B\),
  built by tensoring \cref{def:chartRestrictionAlgHom} with the identity on \(B\).
\end{definition}

\begin{lemma}\leanok
[(b1) Base-change ring square]
  \label{lem:chartBaseChange_ring_square}
  \lean{AlgebraicGeometry.chartBaseChange_ring_square}
  \uses{def:chartBaseChangeRingMap}
  The square of ring maps commutes: \(\iota_R \mathbin{;}\rho_B = \rho\mathbin{;}\iota_{R'}\), where
  \(\iota_R : R\to R\otimes_A B\) and \(\iota_{R'} : R'\to R'\otimes_A B\) are the includeLeft maps.
  \begin{proof}
    Immediate from the functoriality of the tensor-product map with the include-left maps.
  \end{proof}
\end{lemma}

\begin{lemma}\leanok
[Functor-level geometric base-change comparison]
  \label{lem:chartBaseChangeGeometricComparisonNat}
  \lean{AlgebraicGeometry.chartBaseChangeGeometricComparisonNat}
  \uses{lem:chartBaseChange_ring_square}
  The functor-level form of the geometric comparison: the natural isomorphism of iterated pullback
  functors
  \[
    \mathrm{pullback}(\operatorname{Spec}\iota_R) \ggg \mathrm{pullback}(\operatorname{Spec}\rho_B)
      \;\cong\;
    \mathrm{pullback}(\operatorname{Spec}\rho) \ggg \mathrm{pullback}(\operatorname{Spec}\iota_{R'}),
  \]
  built by composing the two pullback pseudofunctor coherences \(\mathrm{pullbackComp}\) with the
  pullback congruence \(\mathrm{pullbackCongr}\) on the \(\operatorname{Spec}\)-image of the
  base-change ring square \cref{lem:chartBaseChange_ring_square}. Its component at \(\widetilde M\) is
  the per-module geometric comparison \cref{lem:chartBaseChangeGeometricComparison}.
  \begin{proof}
    \uses{lem:chartBaseChange_ring_square}
    By \(\operatorname{Spec}\)-contravariance the ring square \cref{lem:chartBaseChange_ring_square}
    yields the geometric square \(\operatorname{Spec}\rho_B \ggg \operatorname{Spec}\iota_R =
    \operatorname{Spec}\iota_{R'} \ggg \operatorname{Spec}\rho\); compose
    \(\mathrm{pullbackComp}(\operatorname{Spec}\rho_B,\operatorname{Spec}\iota_R)\), the
    \(\mathrm{pullbackCongr}\) of that square, and
    \(\mathrm{pullbackComp}(\operatorname{Spec}\iota_{R'},\operatorname{Spec}\rho)^{-1}\).
  \end{proof}
\end{lemma}

\begin{lemma}\leanok
[Geometric base-change comparison of an open immersion]
  \label{lem:chartBaseChangeGeometricComparison}
  \lean{AlgebraicGeometry.chartBaseChangeGeometricComparison}
  \uses{lem:chartBaseChange_ring_square, lem:chartBaseChangeGeometricComparisonNat}
  Pure pullback pseudofunctoriality: \(j_B^{*}(\mathcal{F}_B|_{V_B})\cong\mathcal{F}_B|_{W_B}\). It is
  the component at \(\widetilde M\) of the functor-level geometric comparison
  \cref{lem:chartBaseChangeGeometricComparisonNat},
  \[
    \mathrm{chartBaseChangeGeometricComparison}\,M
      = \bigl(\mathrm{chartBaseChangeGeometricComparisonNat}\bigr).\mathrm{app}\,\widetilde M,
  \]
  definitionally equal to the distributed three-factor form (two pullback-composition isomorphisms
  composed with the congruence isomorphism induced by \cref{lem:chartBaseChange_ring_square} via
  \(\operatorname{Spec}\)-contravariance). Presenting it as this component makes the geometric leg of
  the ring-square glue syntactically the natural-transformation comparison.
  \begin{proof}
    \uses{lem:chartBaseChangeGeometricComparisonNat}
    Evaluate the functor-level comparison \cref{lem:chartBaseChangeGeometricComparisonNat} at
    \(\widetilde M\); by construction this is the composite of the two pullback-composition
    isomorphisms with the congruence isomorphism induced by the ring square
    \cref{lem:chartBaseChange_ring_square}.
  \end{proof}
\end{lemma}

\begin{lemma}\leanok
[(b2-alg) Module reassociation]
  \label{lem:chartBaseChangeModuleReassoc}
  \lean{AlgebraicGeometry.chartBaseChangeModuleReassoc}
  \uses{lem:chartBaseChange_ring_square}
  The base-change module reassociation
  \(\mathrm{ext}_{\rho_B}(\mathrm{ext}_{\iota_R}M)\cong\mathrm{ext}_{\iota_{R'}}(\mathrm{ext}_{\rho}M)\),
  obtained by reassociating two extend-scalars composites along the ring equality
  \cref{lem:chartBaseChange_ring_square}.
  \begin{proof}
    Reassociate the two extend-scalars composites and rewrite along the ring equality
    \cref{lem:chartBaseChange_ring_square}.
  \end{proof}
\end{lemma}

\begin{lemma}


\leanok
[(b2-alg) Component bridge for the algebraic reassociation]
  \label{lem:chartBaseChangeModuleReassoc_eq_natApp}
  \lean{AlgebraicGeometry.chartBaseChangeModuleReassoc_eq_natApp}
  \uses{lem:chartBaseChangeModuleReassoc}
  The per-module algebraic reassociation \cref{lem:chartBaseChangeModuleReassoc} is the \(M\)-component
  of its natural-transformation form \(\mathrm{chartBaseChangeModuleReassocNat}\): for every module
  \(M\),
  \[
    \mathrm{chartBaseChangeModuleReassoc}\,M
      = \bigl(\mathrm{chartBaseChangeModuleReassocNat}\bigr).\mathrm{app}\,M .
  \]
  This identifies the pointwise reassociation isomorphism with the component at \(M\) of its packaged
  natural-transformation form, so a natural-transformation-level identity built from
  \(\mathrm{chartBaseChangeModuleReassocNat}\) may be evaluated at \(M\) without re-deriving the
  reassociation.
  \begin{proof}
    \uses{lem:chartBaseChangeModuleReassoc}
    Both sides are isomorphisms of extend-scalars modules, so it suffices to compare their forward
    maps. The natural transformation \(\mathrm{chartBaseChangeModuleReassocNat}\) is obtained by
    transporting \cref{lem:chartBaseChangeModuleReassoc} along the base-change ring equality through
    the canonical isomorphism across that equality; evaluating the component at \(M\) distributes over
    this composite of an isomorphism, its inverse, and the equality-isomorphism, the last contributing
    only the canonical map across the equal objects at \(M\). The goal reduces to the equality of the
    two ring-square witnesses, which holds by proof irrelevance. Since the carriers are extend-scalars
    modules — no sheafification is involved — the identification stays purely algebraic.
  \end{proof}
\end{lemma}

\begin{lemma}


\leanok
[(b2) Geometric-leg coherence: the geometric comparison is the mate of the pushforward coherence]
  \label{lem:chartBaseChangeGeometricComparison_mate}
  \lean{AlgebraicGeometry.chartBaseChangeGeometricComparison_mate}
  \uses{lem:chartBaseChangeGeometricComparison, lem:chartBaseChange_ring_square}
  The geometric comparison \cref{lem:chartBaseChangeGeometricComparison} is assembled from two
  pullback pseudofunctor coherences (\(\mathrm{pullbackComp}\)) and the pullback congruence on the
  base-change ring square \cref{lem:chartBaseChange_ring_square}. Transported through the composite
  pullback--pushforward adjunction transpose, its forward map agrees with the corresponding composite
  of the two \emph{pushforward} composition coherences (\(\mathrm{pushforwardComp}\)) and the
  pushforward congruence on the same ring square: that is, the pullback coherences from which
  \cref{lem:chartBaseChangeGeometricComparison} is built are the adjoint \emph{mates} of the
  matching pushforward coherences. This is the geometric leg of the ring-square mate-glue
  \cref{lem:pullback_spec_tilde_iso_ring_square_mate_glue}.
  \begin{proof}
    \uses{lem:chartBaseChangeGeometricComparison, lem:chartBaseChange_ring_square,
      lem:conjugateEquiv_pullbackCongr, lem:pushforwardCongr_inv_eq}
    Unfold \cref{lem:chartBaseChangeGeometricComparison} as
    \(\mathrm{pullbackComp}(\operatorname{Spec}\rho_B,\operatorname{Spec}\iota_R)\) followed by the
    pullback congruence on \cref{lem:chartBaseChange_ring_square} followed by
    \(\mathrm{pullbackComp}(\operatorname{Spec}\iota_{R'},\operatorname{Spec}\rho)^{-1}\). Each
    pullback composition isomorphism is, under the pullback--pushforward adjunction, the adjoint
    mate of the corresponding pushforward composition isomorphism: this per-factor coherence is the
    geometric analogue of the algebraic \cref{lem:conjugateEquiv_extendScalarsComp}. Splitting the
    conjugate of the composite into its three factors by the composition law for conjugates and
    discharging each by that coherence (twice, once inverted), together with a single-pair congruence
    reconciliation \cref{lem:conjugateEquiv_pullbackCongr} for the pullback congruence and the matching
    pushforward congruence on \cref{lem:chartBaseChange_ring_square}, yields the matching composite of the two
    \emph{pushforward} composition coherences. Both composites are identity-on-sections over
    \(\operatorname{Spec} R\).
  \end{proof}
\end{lemma}

The geometric leg above uses two small project-bespoke reconciliation helpers, recorded here
for the dependency graph.

\begin{lemma}\leanok
[Mate of a pullback congruence is the pushforward congruence]
  \label{lem:conjugateEquiv_pullbackCongr}
  \lean{AlgebraicGeometry.conjugateEquiv_pullbackCongr}
  \uses{lem:chartBaseChange_ring_square}
  For two equal morphisms \(f = g\) of affine schemes, the adjoint conjugate (mate) of the pullback
  congruence iso \(\mathrm{pullbackCongr}(f=g)\) along the pullback--pushforward adjunctions equals the
  pushforward congruence iso \(\mathrm{pushforwardCongr}(f=g)\).
  \begin{proof}
    Substitute the equality \(f = g\); both congruence isos collapse to identities and the conjugate of
    an identity is the identity (\texttt{conjugateEquiv\_id}).
  \end{proof}
\end{lemma}

\begin{lemma}\leanok
[Inverse of a pushforward congruence is the reversed congruence]
  \label{lem:pushforwardCongr_inv_eq}
  \lean{AlgebraicGeometry.pushforwardCongr_inv_eq}
  The inverse of \(\mathrm{pushforwardCongr}(f=g)\) equals \(\mathrm{pushforwardCongr}(g=f)\).
  \begin{proof}
    Substitute the equality; both sides are identities.
  \end{proof}
\end{lemma}

The algebraic leg below is realized through two reusable mate facts: a
category-agnostic mate-uniqueness criterion and the per-piece conjugate of the
extend-scalars composition coherence. We record them first, together with the four
Mathlib results they rest on as dependency anchors. We also record one bespoke
\(\mathrm{eqToHom}\) reconciliation helper used by the algebraic leg.

\begin{lemma}\leanok
[Mate of an extend-scalars \(\mathrm{eqToHom}\) is the restrict-scalars \(\mathrm{eqToHom}\)]
  \label{lem:conjugateEquiv_extendScalars_eqToHom}
  \lean{AlgebraicGeometry.conjugateEquiv_extendScalars_eqToHom}
  \uses{lem:conjugateEquiv_extendScalarsComp}
  The adjoint conjugate of the \(\mathrm{eqToHom}\) between two extend-scalars functors (induced by a
  ring equality) equals the \(\mathrm{eqToHom}\) between the corresponding restrict-scalars functors.
  \begin{proof}
    Substitute the ring equality; \(\mathrm{eqToHom}\) collapses to the identity and the conjugate of an
    identity is the identity (\texttt{eqToHom\_refl} then \texttt{conjugateEquiv\_id}).
  \end{proof}
\end{lemma}

\begin{lemma}[Right triangle identity of an adjunction]
  \label{lem:right_triangle_components_mathlib}
  \lean{CategoryTheory.Adjunction.right_triangle_components}
  \mathlibok
  \textit{Provided by Mathlib.}
  For an adjunction \(F \dashv G\) with unit \(\eta\) and counit \(\varepsilon\), and
  for every object \(Y\) of the codomain category, the right triangle identity holds:
  the unit at \(G Y\) followed by \(G\) applied to the counit at \(Y\) is the identity,
  \[
    \eta_{G Y} \;\text{followed by}\; G(\varepsilon_Y) \;=\; \mathrm{id}_{G Y}.
  \]
\end{lemma}

\begin{lemma}[Naturality of a natural transformation]
  \label{lem:natTrans_naturality_mathlib}
  \lean{CategoryTheory.NatTrans.naturality}
  \mathlibok
  \textit{Provided by Mathlib.}
  For a natural transformation \(\theta : F \Rightarrow G\) and a morphism
  \(f : X \to Y\), the naturality square commutes: \(F(f)\) followed by \(\theta_Y\)
  equals \(\theta_X\) followed by \(G(f)\).
\end{lemma}

\begin{lemma}[Unit transpose of a conjugate transformation]
  \label{lem:unit_conjugateEquiv_mathlib}
  \lean{CategoryTheory.unit_conjugateEquiv}
  \mathlibok
  \textit{Provided by Mathlib.}
  Let \(\mathrm{adj}_1 : L_1 \dashv R_1\) and \(\mathrm{adj}_2 : L_2 \dashv R_2\) be
  adjunctions with units \(\eta^1, \eta^2\), and let \(\alpha : L_2 \Rightarrow L_1\)
  be a transformation of the left adjoints. Then the conjugate (adjoint mate)
  \(\mathrm{conjugateEquiv}(\alpha) : R_1 \Rightarrow R_2\) satisfies, at every object
  \(c\), the unit-transpose relation
  \[
    \eta^1_c \;\text{followed by}\; \mathrm{conjugateEquiv}(\alpha)_{L_1 c}
    \;=\; \eta^2_c \;\text{followed by}\; R_2(\alpha_c).
  \]
\end{lemma}

\begin{lemma}[Iso-conjugation equivalence of adjoint mates]
  \label{lem:conjugateIsoEquiv_mathlib}
  \lean{CategoryTheory.conjugateIsoEquiv}
  \mathlibok
  \textit{Provided by Mathlib.}
  Let \(\mathrm{adj}_1 : L_1 \dashv R_1\) and \(\mathrm{adj}_2 : L_2 \dashv R_2\) be adjunctions. The
  adjoint-mate (conjugation) operation upgrades to an equivalence between isomorphisms of the left
  adjoints and isomorphisms of the right adjoints,
  \[
    \mathrm{conjugateIsoEquiv} : (L_2 \cong L_1) \;\simeq\; (R_1 \cong R_2),
  \]
  whose action on the forward map of an iso is the conjugate transformation \(\mathrm{conjugateEquiv}\).
\end{lemma}

\begin{lemma}[Forward map of an iso-conjugate]
  \label{lem:conjugateIsoEquiv_apply_hom_mathlib}
  \lean{CategoryTheory.conjugateIsoEquiv_apply_hom}
  \mathlibok
  \textit{Provided by Mathlib.}
  For adjunctions \(\mathrm{adj}_1 : L_1 \dashv R_1\), \(\mathrm{adj}_2 : L_2 \dashv R_2\) and an
  isomorphism \(\alpha : L_2 \cong L_1\) of the left adjoints, the forward map of the iso-conjugate
  \cref{lem:conjugateIsoEquiv_mathlib} is the conjugate (mate) of the forward map of \(\alpha\):
  \[
    \bigl(\mathrm{conjugateIsoEquiv}\,\mathrm{adj}_1\,\mathrm{adj}_2\,\alpha\bigr).\mathrm{hom}
      = \mathrm{conjugateEquiv}\,\mathrm{adj}_1\,\mathrm{adj}_2\,(\alpha.\mathrm{hom}).
  \]
\end{lemma}

\begin{lemma}[Mates equivalence commutes with vertical composition]
  \label{lem:mateEquiv_vcomp_mathlib}
  \lean{CategoryTheory.mateEquiv_vcomp}
  \mathlibok
  \textit{Provided by Mathlib.}
  Let \(\mathrm{adj}_1 : L_1 \dashv R_1\), \(\mathrm{adj}_2 : L_2 \dashv R_2\),
  \(\mathrm{adj}_3 : L_3 \dashv R_3\) be adjunctions, and let \(\alpha\) be a
  \(\mathrm{TwoSquare}\,G_1\,L_1\,L_2\,H_1\) and \(\beta\) a
  \(\mathrm{TwoSquare}\,G_2\,L_2\,L_3\,H_2\). The mate of the horizontal composite square
  equals the vertical composite of the two single-step mates:
  \[
    \mathrm{mateEquiv}\,\mathrm{adj}_1\,\mathrm{adj}_3\,(\alpha \mathbin{\circ_h} \beta)
      \;=\;
    \bigl(\mathrm{mateEquiv}\,\mathrm{adj}_1\,\mathrm{adj}_2\,\alpha\bigr)
      \mathbin{\circ_v}
    \bigl(\mathrm{mateEquiv}\,\mathrm{adj}_2\,\mathrm{adj}_3\,\beta\bigr),
  \]
  where \(\mathbin{\circ_h}\) and \(\mathbin{\circ_v}\) are the horizontal and vertical
  composition of \(\mathrm{TwoSquare}\)s.
\end{lemma}

\begin{lemma}[Mates equivalence commutes with horizontal composition]
  \label{lem:mateEquiv_hcomp_mathlib}
  \lean{CategoryTheory.mateEquiv_hcomp}
  \mathlibok
  \textit{Provided by Mathlib.}
  Let \(\mathrm{adj}_1 : L_1 \dashv R_1\), \(\mathrm{adj}_2 : L_2 \dashv R_2\),
  \(\mathrm{adj}_3 : L_3 \dashv R_3\), \(\mathrm{adj}_4 : L_4 \dashv R_4\) be adjunctions, and
  let \(\alpha\) be a \(\mathrm{TwoSquare}\,G\,L_1\,L_2\,H\) and \(\beta\) a
  \(\mathrm{TwoSquare}\,H\,L_3\,L_4\,K\). The mate of the vertical composite square over the
  composite adjunctions equals the horizontal composite of the two single-step mates:
  \[
    \mathrm{mateEquiv}\,(\mathrm{adj}_1.\mathrm{comp}\,\mathrm{adj}_3)\,
      (\mathrm{adj}_2.\mathrm{comp}\,\mathrm{adj}_4)\,(\alpha \mathbin{\circ_v} \beta)
      \;=\;
    \bigl(\mathrm{mateEquiv}\,\mathrm{adj}_3\,\mathrm{adj}_4\,\beta\bigr)
      \mathbin{\circ_h}
    \bigl(\mathrm{mateEquiv}\,\mathrm{adj}_1\,\mathrm{adj}_2\,\alpha\bigr).
  \]
\end{lemma}

\begin{lemma}[Conjugate of a composite is the reverse composite of conjugates]
  \label{lem:conjugateEquiv_symm_comp_mathlib}
  \lean{CategoryTheory.conjugateEquiv_symm_comp}
  \mathlibok
  \textit{Provided by Mathlib.}
  Let \(\mathrm{adj}_1 : L_1 \dashv R_1\), \(\mathrm{adj}_2 : L_2 \dashv R_2\),
  \(\mathrm{adj}_3 : L_3 \dashv R_3\) be adjunctions between the same pair of categories, and
  let \(\alpha : R_1 \Rightarrow R_2\), \(\beta : R_2 \Rightarrow R_3\) be transformations of
  the right adjoints. The inverse conjugates compose contravariantly:
  \[
    (\mathrm{conjugateEquiv}\,\mathrm{adj}_2\,\mathrm{adj}_3)^{-1}(\beta)
      \;\mathbin{;}\;
    (\mathrm{conjugateEquiv}\,\mathrm{adj}_1\,\mathrm{adj}_2)^{-1}(\alpha)
      \;=\;
    (\mathrm{conjugateEquiv}\,\mathrm{adj}_1\,\mathrm{adj}_3)^{-1}(\alpha \mathbin{;} \beta).
  \]
\end{lemma}

\begin{lemma}[Iterated mate of a four-adjoint square is the conjugate]
  \label{lem:iterated_mateEquiv_conjugateEquiv_mathlib}
  \lean{CategoryTheory.iterated_mateEquiv_conjugateEquiv}
  \mathlibok
  \textit{Provided by Mathlib.}
  Suppose all four functors of a square are left adjoints, with adjunctions
  \(\mathrm{adj}_1 : L_1 \dashv R_1\), \(\mathrm{adj}_2 : L_2 \dashv R_2\),
  \(\mathrm{adj}_3 : F_1 \dashv U_1\), \(\mathrm{adj}_4 : F_2 \dashv U_2\), and let \(\alpha\)
  be a \(\mathrm{TwoSquare}\,F_1\,L_1\,L_2\,F_2\). Then the mate may be iterated, and the
  underlying natural transformation of the iterated mate equals the conjugate over the two
  composite adjunctions:
  \[
    \bigl(\mathrm{mateEquiv}\,\mathrm{adj}_4\,\mathrm{adj}_3\,
      (\mathrm{mateEquiv}\,\mathrm{adj}_1\,\mathrm{adj}_2\,\alpha)\bigr).\mathrm{natTrans}
      \;=\;
    \mathrm{conjugateEquiv}\,(\mathrm{adj}_1.\mathrm{comp}\,\mathrm{adj}_4)\,
      (\mathrm{adj}_3.\mathrm{comp}\,\mathrm{adj}_2)\,\alpha,
  \]
  the underlying transformation being recovered through \(\mathrm{TwoSquare.equivNatTrans}\)
  (the \(\mathrm{.natTrans}\) accessor). In particular the iterated mate is an isomorphism iff
  \(\alpha\) is.
\end{lemma}

\begin{lemma}[Unit transpose of the extend-scalars composition coherence]
  \label{lem:homEquiv_extendScalarsComp_mathlib}
  \lean{ModuleCat.homEquiv_extendScalarsComp}
  \mathlibok
  \textit{Provided by Mathlib.}
  For ring homomorphisms \(f : R_1 \to R_2\), \(g : R_2 \to R_3\) and an
  \(R_1\)-module \(M\), the adjunction transpose, under the extend--restrict
  adjunction for \(g \circ f\), of the component at \(M\) of the extend-scalars
  composition isomorphism \(\mathrm{extendScalarsComp}\,f\,g\) equals the explicit
  unit composite assembled from the units of the extend--restrict adjunctions for
  \(f\) and \(g\) and the inverse of the restrict-scalars composition isomorphism
  \(\mathrm{restrictScalarsComp}\,f\,g\).
\end{lemma}

\begin{lemma}\leanok
[Mate-uniqueness by the unit transpose]
  \label{lem:natTrans_ext_of_unit}
  \lean{AlgebraicGeometry.natTrans_ext_of_unit}
  \uses{lem:right_triangle_components_mathlib, lem:natTrans_naturality_mathlib}
  Let \(\mathrm{adj}_1 : L_1 \dashv R_1\) and \(\mathrm{adj}_2 : L_2 \dashv R_2\) be
  adjunctions between the same pair of categories, with units \(\eta^1, \eta^2\), and
  let \(\alpha : L_2 \Rightarrow L_1\) be a fixed transformation of the left adjoints.
  If two natural transformations \(\beta, \gamma : R_1 \Rightarrow R_2\) of the right
  adjoints have the same unit transpose against \(\alpha\) — that is, for every object
  \(c\),
  \[
    \eta^1_c \;\text{followed by}\; \beta_{L_1 c}
    \;=\; \eta^2_c \;\text{followed by}\; R_2(\alpha_c)
    \;=\; \eta^1_c \;\text{followed by}\; \gamma_{L_1 c}
  \]
  — then \(\beta = \gamma\).
\end{lemma}
\begin{proof}
  \uses{lem:right_triangle_components_mathlib, lem:natTrans_naturality_mathlib}
  Any transformation \(\delta\) satisfying the unit-transpose relation is pinned to the
  single explicit formula
  \[
    \delta_Y \;=\; \eta^2_{R_1 Y} \;\text{followed by}\; R_2(\alpha_{R_1 Y})
      \;\text{followed by}\; R_2(\varepsilon^1_Y),
  \]
  where \(\varepsilon^1\) is the counit of \(\mathrm{adj}_1\). Indeed, inserting the
  right triangle identity \cref{lem:right_triangle_components_mathlib} of
  \(\mathrm{adj}_1\) at \(Y\) writes \(\delta_Y\) as a composite to which naturality
  \cref{lem:natTrans_naturality_mathlib} of \(\delta\) and the unit-transpose
  hypothesis apply, collapsing it to the formula. Since both \(\beta\) and \(\gamma\)
  equal this one formula, they coincide.
\end{proof}

\begin{lemma}\leanok
[Per-piece algebraic mate: conjugate of the extend-scalars coherence]
  \label{lem:conjugateEquiv_extendScalarsComp}
  \lean{AlgebraicGeometry.conjugateEquiv_extendScalarsComp}
  \uses{lem:natTrans_ext_of_unit, lem:unit_conjugateEquiv_mathlib,
        lem:homEquiv_extendScalarsComp_mathlib}
  Let \(f : R \to S\) and \(g : S \to T\) be homomorphisms of commutative rings. Under
  the composite extend--restrict adjunction
  \((\mathrm{adj}_f) \,\square\, (\mathrm{adj}_g)\), transported against the single
  extend--restrict adjunction for \(g \circ f\), the adjoint mate (conjugate) of the
  extend-scalars composition isomorphism \(\mathrm{extendScalarsComp}\,f\,g\) is the
  inverse of the restrict-scalars composition isomorphism:
  \[
    \mathrm{conjugateEquiv}\bigl(\mathrm{extendScalarsComp}\,f\,g\bigr)
    \;=\; \bigl(\mathrm{restrictScalarsComp}\,f\,g\bigr)^{-1}.
  \]
\end{lemma}
\begin{proof}
  \uses{lem:natTrans_ext_of_unit, lem:unit_conjugateEquiv_mathlib,
        lem:homEquiv_extendScalarsComp_mathlib}
  Both sides are natural transformations of the right adjoints, and they share the same
  unit transpose against \(\alpha = \mathrm{extendScalarsComp}\,f\,g\), so they agree by
  the mate-uniqueness criterion \cref{lem:natTrans_ext_of_unit}. For the left-hand side
  this unit transpose is the defining relation of the conjugate,
  \cref{lem:unit_conjugateEquiv_mathlib}. For the right-hand side, the unit transpose of
  \(\bigl(\mathrm{restrictScalarsComp}\,f\,g\bigr)^{-1}\), read off the composite-adjunction
  unit, is exactly the extend-scalars composition coherence
  \cref{lem:homEquiv_extendScalarsComp_mathlib}. The two unit transposes coincide.
\end{proof}

\begin{lemma}


\leanok
[(b2) Algebraic-leg coherence: the module reassociation through \texttt{extendScalarsComp}]
  \label{lem:chartBaseChangeModuleReassoc_extendScalarsComp}
  \lean{AlgebraicGeometry.chartBaseChangeModuleReassoc_extendScalarsComp}
  \uses{lem:chartBaseChangeModuleReassoc, lem:chartBaseChange_ring_square}
  The algebraic reassociation \cref{lem:chartBaseChangeModuleReassoc} is the composite of the two
  extend-scalars composition coherences with the \(\mathrm{eqToIso}\) of the base-change ring square
  \cref{lem:chartBaseChange_ring_square}. Under the algebraic extend--restrict adjunction its forward
  map is the adjoint mate of the corresponding restrict-scalars composition coherence. This is the
  algebraic leg of the ring-square mate-glue
  \cref{lem:pullback_spec_tilde_iso_ring_square_mate_glue}.
  \begin{proof}
    \uses{lem:chartBaseChangeModuleReassoc, lem:chartBaseChange_ring_square,
          lem:natTrans_ext_of_unit, lem:conjugateEquiv_extendScalarsComp,
          lem:conjugateEquiv_extendScalars_eqToHom,
          lem:unit_conjugateEquiv_mathlib, lem:homEquiv_extendScalarsComp_mathlib}
    By definition \cref{lem:chartBaseChangeModuleReassoc} is
    \((\mathrm{extendScalarsComp}\,\iota_R\,\rho_B)^{-1}\) followed by the \(\mathrm{eqToIso}\) of the
    ring equality \cref{lem:chartBaseChange_ring_square} followed by
    \(\mathrm{extendScalarsComp}\,\rho\,\iota_{R'}\), each factor being Mathlib's
    \texttt{ModuleCat.extendScalarsComp}; the matching restrict-scalars reassociation is the
    analogous composite of the \(\mathrm{restrictScalarsComp}\) coherences with the same
    \(\mathrm{eqToIso}\). The claim is that the adjoint mate (conjugate) of the former, under the
    composite extend--restrict adjunctions, is the latter. Both are natural transformations of the
    right adjoints sharing the same unit transpose against \(\alpha = \) the extend-scalars
    reassociation, so they agree by the mate-uniqueness criterion
    \cref{lem:natTrans_ext_of_unit}: the conjugate side is its defining unit transpose
    \cref{lem:unit_conjugateEquiv_mathlib}, while the restrict side is the residual
    \emph{composite-unit coherence} — the three-factor analogue of the per-piece mate fact
    \cref{lem:conjugateEquiv_extendScalarsComp}, namely one extend/restrict-scalars composition
    coherence \cref{lem:homEquiv_extendScalarsComp_mathlib} per factor together with the
    \(\mathrm{eqToIso}\) of the ring square. All constituents are identity-on-elements.
  \end{proof}
\end{lemma}

% ============================================================================
% (b2 glue) Decomposition of the iterated-mate assembly into nat-trans-level steps.
% Each sub-lemma below is closable WITHOUT reassembling the four-leg telescope: the
% pointwise `.app M` route forces the composite-adjunction unit to weak-head-normal
% form (the documented kernel timeout). Steps 1-2 are definitional whiskering lifts
% that fold `.app M` back into whiskered nat-transes; step 3 recognises a SINGLE
% dictionary leg as an iterated single-step mate; steps 4-5 wire the two already-closed
% per-leg mate coherences to the nat-trans level by a single TwoSquare pasting each;
% step 6 is the assembly skeleton, performed entirely at the nat-trans/TwoSquare level
% so no component is ever evaluated.
% ============================================================================

\begin{lemma}


\leanok
[(b2 glue) Whisker-right lift of a pushforward-wrapped leg]
  \label{lem:ring_square_glue_whiskerRight_lift}
  \lean{AlgebraicGeometry.ring_square_glue_whiskerRight_lift}
  Let \(\Phi : F \Rightarrow G\) be a natural transformation between functors
  \(\mathrm{ModuleCat}\,R \to (\operatorname{Spec} S).\mathrm{Modules}\), and let \(g : Y \to Z\) be a
  morphism of schemes. For every object \(M\) the pushforward of the component,
  \(\mathrm{pushforward}(g).\mathrm{map}(\Phi_M)\), equals the \(M\)-component of the right-whiskered
  transformation \((\Phi \rhd \mathrm{pushforward}(g))_M\).
  \begin{proof}
    This is the definitional component formula for right whiskering, Mathlib's
    \texttt{CategoryTheory.whiskerRight\_app}: \((\Phi \rhd H)_M = H.\mathrm{map}(\Phi_M)\) by
    definition. No content beyond unfolding the whiskering; in particular no component of \(\Phi\) is
    expanded, so the leg's composite unit stays folded.
  \end{proof}
\end{lemma}

\begin{lemma}


\leanok
[(b2 glue) Whisker-left lift of an extend-scalars-wrapped leg]
  \label{lem:ring_square_glue_whiskerLeft_lift}
  \lean{AlgebraicGeometry.ring_square_glue_whiskerLeft_lift}
  Let \(\Phi : F \Rightarrow G\) be a natural transformation of functors out of
  \(\mathrm{ModuleCat}\,R'\), and let \(\mathrm{extendScalars}\,\varphi : \mathrm{ModuleCat}\,R \to
  \mathrm{ModuleCat}\,R'\). For every \(M\) the component
  \(\Phi_{(\mathrm{extendScalars}\,\varphi)(M)}\) equals the \(M\)-component of the left-whiskered
  transformation \((\mathrm{extendScalars}\,\varphi \lhd \Phi)_M\).
  \begin{proof}
    This is the definitional component formula for left whiskering, Mathlib's
    \texttt{CategoryTheory.whiskerLeft\_app}: \((F \lhd \Phi)_M = \Phi_{F(M)}\) by definition. No
    content; the wrapped leg becomes a whiskered nat-trans with no component forced.
  \end{proof}
\end{lemma}

\begin{lemma}


\leanok
[(b2 glue) Iterated single-step-mate recognition of one dictionary leg]
  \label{lem:ring_square_glue_pst_iterated_mate}
  \lean{AlgebraicGeometry.ring_square_glue_pst_iterated_mate}
  \uses{lem:pullback_spec_tilde_iso, lem:gammaPushforwardNatIso}
  For a ring map \(\varphi : R \to R'\), the inverse natural transformation underlying the affine
  dictionary \cref{lem:pullback_spec_tilde_iso} --- the conjugate
  \((\mathrm{conjugateEquiv}\,\mathrm{adj}_L\,\mathrm{adj}_R)^{-1}(\mathrm{gammaPushforwardNatIso}\,\varphi)\)
  over the binary composite adjunctions
  \(\mathrm{adj}_L = (\widetilde{(-)}\dashv\Gamma).\mathrm{comp}\,(\mathrm{pullback}\dashv\mathrm{pushforward})\)
  and
  \(\mathrm{adj}_R = (\mathrm{extendScalars}\dashv\mathrm{restrictScalars}).\mathrm{comp}\,(\widetilde{(-)}\dashv\Gamma)\)
  --- is the \emph{iterated single-step mate} of \(\mathrm{gammaPushforwardNatIso}\,\varphi\).
  \begin{proof}
    \uses{lem:pullback_spec_tilde_iso, lem:gammaPushforwardNatIso}
    The recognition is an instance of the \emph{mate-of-a-vertically-composed-adjunction} law: the mate
    of a horizontally pasted square taken over a vertical composite of adjunctions equals the vertical
    composite of the two single-step mates of its factors. This is Mathlib's
    \texttt{CategoryTheory.mateEquiv\_vcomp}: for a horizontal paste \(\alpha \ast \beta\) over the
    composite adjunctions, \(\mathrm{mateEquiv}\,\mathrm{adj}_1\,\mathrm{adj}_3\,(\alpha \ast \beta)\)
    is the vertical composite of \(\mathrm{mateEquiv}\,\mathrm{adj}_1\,\mathrm{adj}_2\,\alpha\) and
    \(\mathrm{mateEquiv}\,\mathrm{adj}_2\,\mathrm{adj}_3\,\beta\). Specialised to the two binary
    composite adjunctions \(\mathrm{adj}_L,\mathrm{adj}_R\) and read in the inverse (\(\mathrm{.symm}\))
    direction, it exhibits the conjugate
    \((\mathrm{conjugateEquiv}\,\mathrm{adj}_L\,\mathrm{adj}_R)^{-1}(\mathrm{gammaPushforwardNatIso}\,\varphi)\)
    as the vertical composite of the two single-step mates of its factors; the conjugate--mate vcomp
    coherences \texttt{CategoryTheory.conjugateEquiv\_mateEquiv\_vcomp} and
    \texttt{CategoryTheory.mateEquiv\_conjugateEquiv\_vcomp} re-express each factor's conjugate as a
    mate, so the pasting closes. In the project's toolchain this composite identity is packaged exactly
    as \texttt{CategoryTheory.iterated\_mateEquiv\_conjugateEquiv} (its \texttt{\_symm} branch for the
    inverse direction), whose underlying natural transformation is recovered through
    \texttt{TwoSquare.equivNatTrans} (the \texttt{.natTrans} accessor). Because the whole step is a
    structural identity of \(\mathrm{TwoSquare}\) pastings, touching only the \emph{single} composite
    adjunction of one leg, it exposes the iterated-mate structure without evaluating any component, so
    the composite unit is never forced to weak-head normal form.
  \end{proof}
\end{lemma}

\begin{lemma}


\leanok
[(b2 glue) Geometric leg at the natural-transformation level]
  \label{lem:ring_square_glue_geom_leg_nat}
  \lean{AlgebraicGeometry.ring_square_glue_geom_leg_nat}
  \uses{lem:chartBaseChangeGeometricComparison_mate, lem:conjugateIsoEquiv_apply_hom_mathlib}
  The hom-level reading of the closed geometric mate coherence
  \cref{lem:chartBaseChangeGeometricComparison_mate}: the adjoint conjugate \(\mathrm{conjugateEquiv}\),
  over the composite pullback--pushforward adjunctions, of the forward map of the functor-level
  geometric comparison \(\mathrm{chartBaseChangeGeometricComparisonNat}\) (built from
  \(\mathrm{pullbackComp}\)) equals the forward map of the pushforward comparison
  \(\mathrm{chartBaseChangePushforwardNat}\) (built from \(\mathrm{pushforwardComp}\)):
  \[
    \mathrm{conjugateEquiv}\,\mathrm{adj}\,\mathrm{adj}'\,
      \bigl(\mathrm{chartBaseChangeGeometricComparisonNat}\bigr).\mathrm{hom}
      = \bigl(\mathrm{chartBaseChangePushforwardNat}\bigr).\mathrm{hom}.
  \]
  This is the \(\mathrm{conjugateEquiv}\)-on-\(\mathrm{hom}\) form of the Iso-level mate leg
  \cref{lem:chartBaseChangeGeometricComparison_mate}, consumed by the four-leg telescope of the
  ring-square glue; no component is evaluated.
  \begin{proof}
    \uses{lem:chartBaseChangeGeometricComparison_mate, lem:conjugateIsoEquiv_apply_hom_mathlib}
    Take the forward-map (\(\mathrm{congrArg}\;\mathrm{hom}\)) image of the closed Iso-level mate leg
    \cref{lem:chartBaseChangeGeometricComparison_mate}, an equality of isomorphisms. Its left side is
    the forward map of an iso-conjugate, which by \cref{lem:conjugateIsoEquiv_apply_hom_mathlib} is the
    conjugate \(\mathrm{conjugateEquiv}\) of the forward map; rewriting along that identity yields the
    displayed natural-transformation equation.
  \end{proof}
\end{lemma}

\begin{lemma}


\leanok
[(b2 glue) Algebraic leg at the natural-transformation level]
  \label{lem:ring_square_glue_alg_leg_nat}
  \lean{AlgebraicGeometry.ring_square_glue_alg_leg_nat}
  \uses{lem:chartBaseChangeModuleReassoc_extendScalarsComp, lem:conjugateIsoEquiv_apply_hom_mathlib}
  The hom-level reading of the closed algebraic mate coherence
  \cref{lem:chartBaseChangeModuleReassoc_extendScalarsComp}: the adjoint conjugate
  \(\mathrm{conjugateEquiv}\), over the composite extend--restrict adjunctions, of the forward map of
  the extend-scalars reassociation \(\mathrm{chartBaseChangeModuleReassocNat}\) equals the forward map
  of the restrict-scalars reassociation \(\mathrm{chartBaseChangeRestrictReassocNat}\):
  \[
    \mathrm{conjugateEquiv}\,\mathrm{adj}\,\mathrm{adj}'\,
      \bigl(\mathrm{chartBaseChangeModuleReassocNat}\bigr).\mathrm{hom}
      = \bigl(\mathrm{chartBaseChangeRestrictReassocNat}\bigr).\mathrm{hom}.
  \]
  This is the \(\mathrm{conjugateEquiv}\)-on-\(\mathrm{hom}\) form of the Iso-level mate leg
  \cref{lem:chartBaseChangeModuleReassoc_extendScalarsComp}, consumed by the four-leg telescope of the
  ring-square glue.
  \begin{proof}
    \uses{lem:chartBaseChangeModuleReassoc_extendScalarsComp, lem:conjugateIsoEquiv_apply_hom_mathlib}
    Take the forward-map (\(\mathrm{congrArg}\;\mathrm{hom}\)) image of the closed Iso-level mate leg
    \cref{lem:chartBaseChangeModuleReassoc_extendScalarsComp}, an equality of isomorphisms. Its left
    side is the forward map of an iso-conjugate, which by
    \cref{lem:conjugateIsoEquiv_apply_hom_mathlib} is the conjugate \(\mathrm{conjugateEquiv}\) of the
    forward map; rewriting along that identity yields the displayed natural-transformation equation.
  \end{proof}
\end{lemma}

\begin{lemma}

\leanok
[(b2 glue) Abstract mate cocycle of the ring square]
  \label{lem:ring_square_cocycle}
  \lean{AlgebraicGeometry.ring_square_cocycle}
  % NOTE (refuted route, slated for excision): SUPERSEDED by the re-scoped extend-scalars cocycle
  % \cref{lem:pullback_spec_tilde_iso_ring_square_natural_extendScalarsCocycle}; retained only
  % because Lean still references this declaration.
  % NOTE: This block states the FAITHFUL abstract mate identity, carrying the four hypotheses
  % (i)-(iv). The realization that folds this abstract telescope directly onto the concrete
  % heavy ring-square goal (whether by a bare instantiation or by interposing identity-transport
  % bridges at the telescope boundaries) is refuted: the compounded boundary defeq cannot be
  % eliminated, only relocated. Realization is delegated to the two routes of
  % \cref{lem:pullback_spec_tilde_iso_ring_square_natural_catCommSq} and
  % \cref{lem:pullback_spec_tilde_iso_ring_square_natural_concrete} below.
  \uses{lem:mateEquiv_vcomp_mathlib, lem:mateEquiv_hcomp_mathlib,
        lem:conjugateEquiv_symm_comp_mathlib, lem:iterated_mateEquiv_conjugateEquiv_mathlib,
        lem:conjugateIsoEquiv_mathlib}
  A purely category-theoretic cocycle, stated and proved over \emph{abstract} categories, functors
  and adjunctions: \emph{no} sheafification \(\widetilde{(-)}\), pullback
  \(\mathrm{pullback}(\operatorname{Spec}-)\), or extend-scalars functor occurs anywhere in the
  statement or the proof, so that nothing can be forced to weak-head normal form. Fix abstract
  categories and four columns of vertically composable adjunctions arranged as the two
  composite-left-adjoint legs of a commuting square of categories: a geometric leg
  \(L_{\mathrm{geom}} = T \ggg P\) with right adjoint \(R_{\mathrm{geom}}\) and an algebraic leg
  \(L_{\mathrm{alg}} = E \ggg T'\) with right adjoint \(R_{\mathrm{alg}}\), together with a base
  comparison family \(\gamma\) (the abstract stand-in for \(\mathrm{gammaPushforwardNatIso}\)) and the
  per-leg comparison transformations of the square. Suppose:
  \begin{enumerate}
    \item[(i)] (\emph{per-leg unit triangle} --- abstract form of
      \cref{lem:pullback_spec_tilde_iso_inv_unit_triangle}) each dictionary leg's inverse is the
      adjoint conjugate, over its composite adjunctions, of the corresponding base comparison
      \(\gamma\), as recorded by the defining unit triangle;
    \item[(ii)] (\emph{geometric-leg mate equation} --- abstract form of
      \cref{lem:ring_square_glue_geom_leg_nat}) the conjugate of the geometric comparison equals the
      pushforward comparison;
    \item[(iii)] (\emph{algebraic-leg mate equation} --- abstract form of
      \cref{lem:ring_square_glue_alg_leg_nat}) the conjugate of the extend-scalars reassociation equals
      the restrict-scalars reassociation;
    \item[(iv)] (\emph{composition coherence} --- abstract form of
      \cref{lem:gammaPushforwardNatIso_comp}) the base comparison \(\gamma\) is multiplicative across
      the two ways of composing the square's edges.
  \end{enumerate}
  Then the iterated mates of the two legs agree, as an equation of conjugate/mate isomorphisms
  \[
    \mathrm{conjugateIsoEquiv}\,\mathrm{adj}_{\mathrm{geom}}\,\mathrm{adj}_{\mathrm{geom}}'\,
      (\Phi_{\mathrm{geom}})
      \;\ggg\;
      (\cdots)
      \;=\;
    \mathrm{conjugateIsoEquiv}\,\mathrm{adj}_{\mathrm{alg}}\,\mathrm{adj}_{\mathrm{alg}}'\,
      (\Phi_{\mathrm{alg}})
      \;\ggg\;
      (\cdots),
  \]
  an identity between two \(\ggg\)-telescopes of abstract conjugate/mate isomorphisms in which no
  concrete functor is named.

  \medskip
  \emph{Status of the realization (open).} The abstract cocycle below is an identity of the
  \(\mathrm{TwoSquare}\)/conjugate calculus and is established on its own terms; its hypotheses
  (i)--(iv) are each discharged by the concrete coherences cited above. What is \emph{open} is the
  \emph{realization} of the concrete ring-square naturality
  \cref{lem:pullback_spec_tilde_iso_ring_square_natural} by instantiating this cocycle at the heavy
  data. Folding the abstract telescope onto the concrete goal by a single instantiation forces a
  compounded definitional reduction of the four telescope boundaries (the source of each factor must
  be matched against the target of the previous one, and those boundaries are definitionally but not
  syntactically equal). Interposing identity-transport bridges at the boundaries relocates but does
  not eliminate that compounded reduction, so the direct fold does not realize the crux. The
  realization is therefore carried out by the two sibling routes below — the functor-paste route
  \cref{lem:pullback_spec_tilde_iso_ring_square_natural_catCommSq}, which keeps adjacency at the
  level of named functors, and the concrete tilde route
  \cref{lem:pullback_spec_tilde_iso_ring_square_natural_concrete}, which bypasses the mate apparatus
  entirely.
  \begin{proof}
    \uses{lem:mateEquiv_vcomp_mathlib, lem:mateEquiv_hcomp_mathlib,
          lem:conjugateEquiv_symm_comp_mathlib, lem:iterated_mateEquiv_conjugateEquiv_mathlib,
          lem:conjugateIsoEquiv_mathlib}
    Reduce the equality of isomorphisms to the underlying inverse natural-transformation equality.
    Transpose each leg through its composite adjunctions: by \cref{lem:conjugateEquiv_symm_comp_mathlib}
    the inverse conjugate of a composite is the (contravariant) composite of the inverse conjugates, so
    the two routings telescope into vertical chains of per-edge conjugates. Recognise each edge's
    iterated single-step mate as a genuine conjugate by
    \cref{lem:iterated_mateEquiv_conjugateEquiv_mathlib} (its underlying natural transformation recovered
    through \(\mathrm{TwoSquare.equivNatTrans}\), the \(\mathrm{.natTrans}\) accessor), and fuse the
    per-leg \(\mathrm{TwoSquare}\) pastings with \cref{lem:mateEquiv_vcomp_mathlib} and
    \cref{lem:mateEquiv_hcomp_mathlib} --- the mate of a composite square is the composite of the mates,
    so the four-edge telescope assembles structurally without touching a component. Discharge the four
    legs by the four hypotheses: hypothesis (i) replaces each dictionary leg by the conjugate of the
    base comparison; hypotheses (ii) and (iii) discharge the geometric and algebraic comparison edges;
    hypothesis (iv) closes the residual on the multiplicativity of the base comparison across the
    square's two edge-orders. Re-upgrading to isomorphisms by \cref{lem:conjugateIsoEquiv_mathlib} yields
    the displayed \(\ggg\)-equation. Because the entire argument lives in the abstract
    \(\mathrm{TwoSquare}\)/conjugate calculus --- no carrier appears in the statement or in any
    intermediate term --- no weak-head-normal-form reduction of a heavy carrier can ever be triggered.
  \end{proof}
\end{lemma}

\begin{lemma}


\leanok
[(b2) Iterated-mate glue of the ring-square naturality]
  \label{lem:pullback_spec_tilde_iso_ring_square_mate_glue}
  \lean{AlgebraicGeometry.pullback_spec_tilde_iso_ring_square_mate_glue}
  % NOTE (refuted route, slated for excision): SUPERSEDED by the re-scoped extend-scalars cocycle
  % \cref{lem:pullback_spec_tilde_iso_ring_square_natural_extendScalarsCocycle}; retained only
  % because Lean still references this declaration.
  \uses{lem:ring_square_cocycle, lem:ring_square_glue_geom_leg_nat,
        lem:ring_square_glue_alg_leg_nat, lem:pullback_spec_tilde_iso_inv_unit_triangle,
        lem:gammaPushforwardNatIso_comp, lem:pullback_spec_tilde_iso}
  After the ring-square equation of \cref{lem:pullback_spec_tilde_iso_ring_square_natural} is
  transposed through the three adjunctions of the composite left adjoint
  \(L_{\mathrm{tot}} = \widetilde{(-)} \ggg \mathrm{pullback}(\operatorname{Spec}\iota_R) \ggg
  \mathrm{pullback}(\operatorname{Spec}\rho_B)\), it becomes a single equation in
  \(\mathrm{ModuleCat}\,R\) between two maps \(M \to \mathrm{restr}_{\iota_R}\,\mathrm{restr}_{\rho_B}\,N\).
  That equation holds: each \cref{lem:pullback_spec_tilde_iso} leg (for the four ring maps
  \(\iota_R,\rho,\rho_B,\iota_{R'}\)) transposes to \cref{lem:gammaPushforwardNatIso} by
  \cref{lem:pullback_spec_tilde_iso_inv_unit_triangle}, the geometric comparison transposes by
  \cref{lem:chartBaseChangeGeometricComparison_mate}, and the algebraic reassociation transposes by
  \cref{lem:chartBaseChangeModuleReassoc_extendScalarsComp}; the three are then assembled through the
  nested \(\mathrm{pushforward}.\mathrm{map}\)/\(\Gamma.\mathrm{map}\) wrappers and closed on the
  composition coherence \cref{lem:gammaPushforwardNatIso_comp}.
  \begin{proof}
    \uses{lem:ring_square_cocycle, lem:ring_square_glue_geom_leg_nat,
          lem:ring_square_glue_alg_leg_nat, lem:pullback_spec_tilde_iso_inv_unit_triangle,
          lem:gammaPushforwardNatIso_comp, lem:pullback_spec_tilde_iso}
    All of the telescoping content has been extracted into the abstract cocycle
    \cref{lem:ring_square_cocycle}, whose statement and proof name no carrier. It therefore suffices to
    instantiate that cocycle at the present heavy data, supplying the four hypotheses
    \cref{lem:pullback_spec_tilde_iso_inv_unit_triangle} (the per-leg unit triangle, instantiated at the
    four ring maps \(\iota_R,\rho,\rho_B,\iota_{R'}\) through the dictionary
    \cref{lem:pullback_spec_tilde_iso}), \cref{lem:ring_square_glue_geom_leg_nat} (the geometric leg),
    \cref{lem:ring_square_glue_alg_leg_nat} (the algebraic leg), and
    \cref{lem:gammaPushforwardNatIso_comp} (the composition coherence). This instantiation of the abstract
    cocycle at the heavy data is, however, \emph{not} a direct instantiation: a single unification of the
    abstract telescope against the concrete goal forces a compounded weak-head-normal-form reduction of
    the four boundary identifications (the source of each factor must be matched against the target of the
    previous one, and those boundaries are definitionally-but-not-syntactically equal), which is exactly
    the reduction that times out.

    \medskip
    \emph{Status (open).} This iterated-mate fold is \emph{not} a viable realization of the crux.
    Interposing identity-transport bridges at the four telescope boundaries — so that adjacency
    becomes syntactic and each boundary identity is an isolated obligation — relocates the compounded
    boundary reduction but does not eliminate it, so the closing instantiation still forces the
    carrier reductions it was meant to avoid. The ring-square naturality
    \cref{lem:pullback_spec_tilde_iso_ring_square_natural} is therefore realized not through this
    glue but through the two sibling routes below: the functor-paste route
    \cref{lem:pullback_spec_tilde_iso_ring_square_natural_catCommSq} and the concrete tilde route
    \cref{lem:pullback_spec_tilde_iso_ring_square_natural_concrete}. This block is retained as the
    abstract record of the mate decomposition; its instantiation remains open.
  \end{proof}
\end{lemma}

\begin{lemma}\leanok
[(b2) Ring-square naturality of the tilde-pullback dictionary]
  \label{lem:pullback_spec_tilde_iso_ring_square_natural}
  \lean{AlgebraicGeometry.pullback_spec_tilde_iso_ring_square_natural}
  \uses{lem:pullback_spec_tilde_iso, lem:gammaPushforwardNatIso, lem:chartBaseChange_ring_square,
        lem:chartBaseChangeGeometricComparison, lem:chartBaseChangeModuleReassoc}
  \textit{Source: Stacks Project, Tag 01I9, part~(1), the word ``functorially''
  --- the verbatim quote is in \cref{lem:pullback_tilde_restriction_natural}.}
  The explicit affine tilde-pullback isomorphism \cref{lem:pullback_spec_tilde_iso} is natural across
  the base-change ring square \cref{lem:chartBaseChange_ring_square}: the forward map of the
  geometric comparison \cref{lem:chartBaseChangeGeometricComparison} and the forward map of the
  algebraic reassociation \cref{lem:chartBaseChangeModuleReassoc} agree once transported through the two
  per-chart dictionaries of \cref{lem:pullback_spec_tilde_iso} for \(V\) and for \(W\). This is the
  residual crux of \cref{lem:pullback_tilde_restriction_natural}.
  \begin{proof}
    \uses{lem:pullbackSpecTildeNatIso, lem:pullback_spec_tilde_iso,
          lem:chartBaseChange_extendScalars_cocycle,
          lem:pullback_spec_tilde_iso_ring_square_natural_extendScalarsCocycle}
    This equation is realized by the re-scoped \emph{extend-scalars cocycle} route
    \cref{lem:pullback_spec_tilde_iso_ring_square_natural_extendScalarsCocycle} of the following
    subsection, which is the live discharge of the crux. In outline: rewrite each of the four
    geometric legs by the proven single-map dictionary
    \cref{lem:pullbackSpecTildeNatIso}/\cref{lem:pullback_spec_tilde_iso} — turning each
    \(\Gamma(\mathrm{pullback}(\operatorname{Spec}\varphi_i))(\widetilde M) \cong
    \widetilde{\mathrm{ext}_{\varphi_i} M}\) \emph{before} any composition is formed, so that the
    composite pullback coherence \(\Gamma(\mathrm{pullbackComp})\), although present in the geometric
    comparison as an \emph{abstract} natural isomorphism, is never \emph{unfolded} — never forced to
    weak-head normal form on an object; the
    residual coherence is then the purely algebraic base-change cocycle
    \cref{lem:chartBaseChange_extendScalars_cocycle}, an isomorphism of light \(\mathrm{ModuleCat}\)
    carriers, and the sheaf-level equation follows by applying \(\widetilde{(-)}\) once and invoking its
    functoriality. No flatness, no adjoint mate of the base-change map, and no evaluation of
    \(\Gamma(\mathrm{pullbackComp})\) enter. The earlier iterated-mate and concrete-section routes
    (see the refuted-routes subsection below) are superseded and slated for excision.
  \end{proof}
\end{lemma}

\subsection{Re-scoped crux: the base-change cocycle at the extendScalars level}

The decisive structural fact about the ring-square crux
\cref{lem:pullback_spec_tilde_iso_ring_square_natural} is that the obstruction in every adjoint-mate
or section-level discharge is the evaluation of the pullback coherence
\(\Gamma(\mathrm{pullbackComp})\) of a \emph{composite} of \(\operatorname{Spec}\)-pullbacks: there
is no comparison \((\operatorname{Spec}\varphi)^{*}\widetilde M \cong \widetilde{\mathrm{ext}_\varphi M}\)
available as \(\widetilde{(-)}\) of an explicit module map at the level of a composite, so once a
composite pullback coherence is formed the identification with an extend-scalars reassociation has
no exit. The re-scoped route removes that obstruction at the source by never forming a composite
geometric coherence at all. It proceeds in three steps.

\begin{enumerate}
  \item \emph{Rewrite each geometric leg first.} Each of the four ring maps
    \(\varphi_i \in \{\iota_R,\rho,\rho_B,\iota_{R'}\}\) of the base-change ring square
    \cref{lem:chartBaseChange_ring_square} is handled \emph{singly} by the proven affine dictionary
    \cref{lem:pullbackSpecTildeNatIso}/\cref{lem:pullback_spec_tilde_iso}, which identifies
    \(\widetilde{(-)} \ggg \mathrm{pullback}(\operatorname{Spec}\varphi_i) \cong
    \mathrm{ext}_{\varphi_i} \ggg \widetilde{(-)}\) for a \emph{single} map. Performing these four
    rewrites \emph{before} any composition means the composite pullback coherence
    \(\Gamma(\mathrm{pullbackComp})\) of \(\operatorname{Spec}\)-pullbacks — the term whose weak-head
    reduction is the documented wall — is kept \emph{abstract}: it is present in the geometric
    comparison only as a natural isomorphism and is never unfolded on an object. The live constraint on
    the realization is instead the \(200\)k-heartbeat per-declaration elaboration budget, which the fine
    leg-by-leg decomposition exists to beat.
  \item \emph{Discharge the residual cocycle algebraically.} What remains after the four single-map
    rewrites is the purely algebraic base-change reassociation
    \[
      (R'\otimes_A B) \otimes_{R\otimes_A B}\bigl((R\otimes_A B)\otimes_R M\bigr)
        \;\cong\;
      (R'\otimes_A B) \otimes_{R'}\bigl(R'\otimes_R M\bigr),
    \]
    i.e. the module reassociation \cref{lem:chartBaseChangeModuleReassoc} of the ring square, now
    \emph{computed} through the extend-scalars pseudofunctor: its associativity coherence
    \cref{lem:extendScalars_assoc_mathlib} (the \(\mathrm{mapComp}\)-associativity of
    \cref{lem:moduleCatExtendScalarsPseudofunctor_mathlib}, built on the composition isomorphism
    \cref{lem:extendScalarsComp_mathlib}) together with the cancellation of an iterated base change
    \cref{lem:isPushout_cancelBaseChange_mathlib} / \cref{lem:isBaseChange_mathlib}. Every carrier
    here is a light \(\mathrm{ModuleCat}\), so no sheafification or pullback functor is forced to
    weak-head normal form.
  \item \emph{Apply \(\widetilde{(-)}\) once.} The assembled algebraic isomorphism is mapped by the
    sheafification functor \(\widetilde{(-)}\) (\(\mathrm{tilde.functor}\)) a \emph{single} time;
    functoriality \(\widetilde{(-)}(f \ggg g) = \widetilde{(-)}(f) \ggg \widetilde{(-)}(g)\)
    (\(\mathrm{Functor.map\_comp}\)) distributes it across the legs and produces the sheaf-level
    cocycle for free.
\end{enumerate}

This is exactly Stacks Tag 01I9 part~(1)'s assertion that the identification
\(\psi^{*}\widetilde M = \widetilde{S\otimes_R M}\) holds \emph{functorially in the \(R\)-module
\(M\)}: the functoriality, applied per single ring map and then closed by algebraic base-change
associativity, is the cocycle. We first record the Mathlib results the algebraic cocycle stands on
as dependency anchors, then state the algebraic cocycle and the re-scoped realization of the crux.

\begin{lemma}[Extend-scalars composition isomorphism]
  \label{lem:extendScalarsComp_mathlib}
  \lean{ModuleCat.extendScalarsComp}
  \mathlibok
  \textit{Provided by Mathlib.}
  For ring homomorphisms \(f : R_1 \to R_2\) and \(g : R_2 \to R_3\), the extend-scalars functor of
  the composite agrees with the composite of the extend-scalars functors, as a natural isomorphism
  \[
    \mathrm{ext}_{g \circ f} \;\cong\; \mathrm{ext}_f \ggg \mathrm{ext}_g,
  \]
  with component at \(M\) the canonical \(R_3\)-linear isomorphism
  \(R_3 \otimes_{R_1} M \cong R_3 \otimes_{R_2} (R_2 \otimes_{R_1} M)\).
\end{lemma}

\begin{lemma}[Associativity coherence of extend-scalars composition]
  \label{lem:extendScalars_assoc_mathlib}
  \lean{ModuleCat.extendScalars_assoc'}
  \mathlibok
  \uses{lem:extendScalarsComp_mathlib}
  \textit{Provided by Mathlib.}
  For ring homomorphisms \(f_{12} : R_1 \to R_2\), \(f_{23} : R_2 \to R_3\),
  \(f_{34} : R_3 \to R_4\), the two ways of reassociating a triple extend-scalars composite agree:
  the composition isomorphisms \cref{lem:extendScalarsComp_mathlib} for the two bracketings, threaded
  through the functor associator, compose to the identity on \(\mathrm{ext}_{f_{34}\circ f_{23}\circ
  f_{12}}\). This is the pentagon/associativity coherence underlying the extend-scalars pseudofunctor
  \cref{lem:moduleCatExtendScalarsPseudofunctor_mathlib}.
\end{lemma}

\begin{lemma}[Extend-scalars pseudofunctor on commutative rings]
  \label{lem:moduleCatExtendScalarsPseudofunctor_mathlib}
  \lean{CommRingCat.moduleCatExtendScalarsPseudofunctor}
  \mathlibok
  \uses{lem:extendScalarsComp_mathlib}
  \textit{Provided by Mathlib.}
  The assignment \(R \mapsto \mathrm{ModuleCat}\,R\), \(f \mapsto \mathrm{ext}_f\) is a pseudofunctor
  from the locally-discrete bicategory on \(\mathrm{CommRingCat}\) to \(\mathrm{Cat}\), with
  \(\mathrm{mapComp}\) the composition isomorphism \cref{lem:extendScalarsComp_mathlib} and
  \(\mathrm{mapId}\) the identity-extension isomorphism. It packages the extend-scalars composition
  coherences — in particular the associativity \cref{lem:extendScalars_assoc_mathlib} — as
  pseudofunctor data.
\end{lemma}

\begin{lemma}[Cancellation of an iterated base change]
  \label{lem:isPushout_cancelBaseChange_mathlib}
  \lean{Algebra.IsPushout.cancelBaseChange}
  \mathlibok
  \textit{Provided by Mathlib.}
  Let \(R \to S\) and \(R \to A\) be ring maps with pushout \(B = S \otimes_R A\) (an
  \(\mathrm{Algebra.IsPushout}\) square), and let \(M\) be an \(A\)-module. Then the iterated base
  change cancels:
  \[
    B \otimes_A M \;\cong\; S \otimes_R M
  \]
  as \(S\)-modules, with \(1 \otimes_A m \mapsto 1 \otimes_R m\). Instantiated at the affine--affine
  base-change model it identifies \((R'\otimes_A B)\otimes_{R\otimes_A B}\bigl((R\otimes_A B)\otimes_R
  M\bigr)\) with \((R'\otimes_A B)\otimes_{R'}\bigl(R'\otimes_R M\bigr)\).
\end{lemma}

\begin{lemma}[Base-change property of a linear map]
  \label{lem:isBaseChange_mathlib}
  \lean{IsBaseChange}
  \mathlibok
  \textit{Provided by Mathlib.}
  For an \(R\)-algebra \(S\) and an \(R\)-linear map \(f : M \to N\) with \(N\) an \(S\)-module,
  \(\mathrm{IsBaseChange}\,S\,f\) is the property that \((S, f)\) exhibits \(N\) as the base change
  \(S \otimes_R M\): every \(S\)-module receives a unique \(S\)-linear factorization through \(f\).
  It is the universal-property packaging of \(\mathrm{cancelBaseChange}\)
  \cref{lem:isPushout_cancelBaseChange_mathlib}, used to identify the carriers of the algebraic
  cocycle without choosing tensor presentations.
\end{lemma}

\begin{lemma}\leanok
[Algebraic base-change cocycle through the extend-scalars pseudofunctor]
  \label{lem:chartBaseChange_extendScalars_cocycle}
  \lean{AlgebraicGeometry.chartBaseChange_extendScalars_cocycle}
  \uses{lem:chartBaseChangeModuleReassoc, lem:chartBaseChange_ring_square,
        lem:extendScalarsComp_mathlib, lem:extendScalars_assoc_mathlib,
        lem:moduleCatExtendScalarsPseudofunctor_mathlib,
        lem:isPushout_cancelBaseChange_mathlib, lem:isBaseChange_mathlib}
  % NOTE (iter-031 review): Lean declaration now EXISTS and is sorry-free (created + proved this
  % iter; cold build green). PROSE-vs-LEAN DIVERGENCE: the prose below describes the aspirational
  % "natively, rather than an eqToIso transport" route via cancelBaseChange/IsBaseChange; the
  % landed Lean realizes the SAME iso via `extendScalarsComp ∘ eqToIso(ring-square) ∘ extendScalarsComp`
  % (one unavoidable eqToIso transport along the ring square). Same iso, sorry-free; the cancelBaseChange
  % refinement is optional polish (task_results). Planner: reconcile the "rather than eqToIso" prose.
  The base-change module reassociation \cref{lem:chartBaseChangeModuleReassoc}
  \[
    \mathrm{ext}_{\rho_B}\bigl(\mathrm{ext}_{\iota_R} M\bigr)
      \;\cong\;
    \mathrm{ext}_{\iota_{R'}}\bigl(\mathrm{ext}_{\rho} M\bigr)
  \]
  of the ring square \cref{lem:chartBaseChange_ring_square}, presented \emph{natively} through the
  extend-scalars pseudofunctor rather than through an \(\mathrm{eqToIso}\) transport. Both sides are
  the value at \(M\) of a triple extend-scalars composite over the two routings of the square; their
  identity is the associativity coherence \cref{lem:extendScalars_assoc_mathlib} of the pseudofunctor
  \cref{lem:moduleCatExtendScalarsPseudofunctor_mathlib} together with the base-change cancellation
  \cref{lem:isPushout_cancelBaseChange_mathlib}. Every carrier is a \(\mathrm{ModuleCat}\); no
  sheafification or pullback functor occurs.
  \begin{proof}
    \uses{lem:chartBaseChangeModuleReassoc, lem:chartBaseChange_ring_square,
          lem:extendScalarsComp_mathlib, lem:extendScalars_assoc_mathlib,
          lem:moduleCatExtendScalarsPseudofunctor_mathlib,
          lem:isPushout_cancelBaseChange_mathlib, lem:isBaseChange_mathlib}
    Both routings of the ring square \cref{lem:chartBaseChange_ring_square}
    \(\iota_R \mathbin{;}\rho_B = \rho\mathbin{;}\iota_{R'}\) give a triple of composable ring maps
    out of \(R\) ending at \(R'\otimes_A B\); their extend-scalars images are reassociated by the
    composition isomorphisms \cref{lem:extendScalarsComp_mathlib} for the two bracketings, and these
    agree by the pseudofunctor associativity \cref{lem:extendScalars_assoc_mathlib}. The remaining
    identification of the two middle carriers — \((R\otimes_A B)\otimes_R M\) against
    \(R'\otimes_R M\) re-based to \(R'\otimes_A B\) — is the cancellation of an iterated base change
    \cref{lem:isPushout_cancelBaseChange_mathlib}, whose universal property
    \cref{lem:isBaseChange_mathlib} fixes it without a choice of tensor presentation. The composite
    is the reassociation \cref{lem:chartBaseChangeModuleReassoc}, now exhibited as a cocycle of light
    module isomorphisms.
  \end{proof}
\end{lemma}

% ============================================================================
% (Re-scoped realization, FINE decomposition — iter-032 effort-breaker)
% The realization `…_extendScalarsCocycle` below is decomposed into a \uses-linked chain so that
% every sub-statement elaborates within the 200k-heartbeat per-declaration budget. The wall is NOT
% that `pullbackComp` is built (it is — inside the geometric comparison); the wall is the ELABORATION
% cost of a single monolithic functor-level glue whose heavy `extendScalars`/ModuleCat-diamond whisker
% carriers blow the budget (the monolithic form elaborated only at 800k). The decomposition keeps
% `pullbackComp` ABSTRACT — present as a natural iso, never evaluated on an object (object-level
% evaluation is the dead `_concrete` wall) — and splits the glue into: a single-map dictionary fold
% `pullbackSpecTildeComp`, its multiplicativity `pullbackSpecTildeComp_eq_composite` (the one place
% `pullbackComp` enters, abstractly), the geometric-comparison unfold, and the functor-level master
% glue. The realization is then a single `.app M` of the master glue. The two Mathlib pullback
% coherences are recorded as anchors first.
% ============================================================================

\begin{lemma}[Pullback-composition coherence of the modules pseudofunctor]
  \label{lem:pullbackComp_mathlib}
  \lean{AlgebraicGeometry.Scheme.Modules.pullbackComp}
  \mathlibok
  \textit{Provided by Mathlib.}
  For composable scheme morphisms \(f : X \to Y\), \(g : Y \to Z\), the modules-pullback pseudofunctor
  is coherent under composition: a natural isomorphism
  \[
    \mathrm{pullback}(f \ggg g) \;\cong\; \mathrm{pullback}\,g \ggg \mathrm{pullback}\,f
  \]
  of functors \(Z.\mathrm{Modules} \to X.\mathrm{Modules}\). It is used here only as an \emph{abstract}
  natural isomorphism: it is never evaluated on an object (the object-level evaluation of
  \(\Gamma(\mathrm{pullbackComp})\) of a composite of \(\operatorname{Spec}\)-pullbacks is the
  documented kernel wall).
\end{lemma}

\begin{lemma}[Pullback congruence along an equality of morphisms]
  \label{lem:pullbackCongr_mathlib}
  \lean{AlgebraicGeometry.Scheme.Modules.pullbackCongr}
  \mathlibok
  \textit{Provided by Mathlib.}
  For parallel scheme morphisms with \(f = g : X \to Y\), the modules-pullback functors are naturally
  isomorphic, \(\mathrm{pullback}\,f \cong \mathrm{pullback}\,g\), the congruence isomorphism induced
  by the equality. It supplies the middle factor of the geometric comparison
  \cref{lem:chartBaseChangeGeometricComparisonNat} from the \(\operatorname{Spec}\)-image of the
  base-change ring square.
\end{lemma}

\begin{lemma}[Functoriality of a functor over composition]
  \label{lem:functor_map_comp_mathlib}
  \lean{CategoryTheory.Functor.map_comp}
  \mathlibok
  \textit{Provided by Mathlib.}
  For a functor \(F\) and composable morphisms \(f, g\), the image of the composite is the composite of
  the images, \(F(f \ggg g) = F(f) \ggg F(g)\). It is used here only with \(F = \widetilde{(-)}\)
  (\(\mathrm{tilde.functor}\)): the single sheafification of the assembled algebraic isomorphism
  distributes across the leg-by-leg factorization, so \(\widetilde{(-)}\) is applied to the algebraic
  cocycle exactly \emph{once} (its target carriers are light \(\mathrm{ModuleCat}\) maps; no composite
  pullback coherence is materialized by this step).
\end{lemma}

\begin{lemma}
[Defining unfold of the functor-level geometric comparison]
  \label{lem:chartBaseChangeGeometricComparisonNat_unfold}
  \lean{AlgebraicGeometry.chartBaseChangeGeometricComparisonNat_unfold}
  \uses{lem:chartBaseChangeGeometricComparisonNat, lem:pullbackComp_mathlib,
        lem:pullbackCongr_mathlib, lem:chartBaseChange_ring_square}
  The functor-level geometric comparison \cref{lem:chartBaseChangeGeometricComparisonNat} is, as an
  abstract natural isomorphism, the three-factor composite
  \[
    \mathrm{chartBaseChangeGeometricComparisonNat}
      = \mathrm{pullbackComp}(\operatorname{Spec}\rho_B,\operatorname{Spec}\iota_R)
        \;\ggg\; \mathrm{pullbackCongr}(\mathrm{hsq})
        \;\ggg\; \mathrm{pullbackComp}(\operatorname{Spec}\iota_{R'},\operatorname{Spec}\rho)^{-1},
  \]
  where \(\mathrm{hsq} : \operatorname{Spec}\rho_B \ggg \operatorname{Spec}\iota_R =
  \operatorname{Spec}\iota_{R'} \ggg \operatorname{Spec}\rho\) is the \(\operatorname{Spec}\)-image of
  the base-change ring square \cref{lem:chartBaseChange_ring_square}. Recording this defining equation
  as its own block lets the heavier glue statements quote the geometric comparison as an \emph{opaque}
  symbol and unfold it only inside a localized step, so the \(\mathrm{pullbackComp}\) coherences are
  never forced to weak-head normal form during elaboration of those statements.
  \begin{proof}
    \uses{lem:chartBaseChangeGeometricComparisonNat, lem:chartBaseChange_ring_square}
    This is the definition of \cref{lem:chartBaseChangeGeometricComparisonNat}: by
    \(\operatorname{Spec}\)-contravariance the ring square \cref{lem:chartBaseChange_ring_square} gives
    \(\mathrm{hsq}\), and the comparison is by construction the displayed composite of the two
    \(\mathrm{pullbackComp}\) coherences \cref{lem:pullbackComp_mathlib} with the
    \(\mathrm{pullbackCongr}\) \cref{lem:pullbackCongr_mathlib} of \(\mathrm{hsq}\); the equation holds
    by unfolding.
  \end{proof}
\end{lemma}

\begin{lemma}

\leanok
[Two-leg fold of the affine pullback dictionary]
  \label{lem:pullbackSpecTildeComp}
  \lean{AlgebraicGeometry.pullbackSpecTildeComp}
  \uses{lem:pullbackSpecTildeNatIso}
  For composable ring maps \(\varphi : R \to S\) and \(\varphi' : S \to T\), the two single-map
  dictionaries \cref{lem:pullbackSpecTildeNatIso} for \(\varphi\) and \(\varphi'\) assemble, by
  whiskering, into the \emph{two-leg} natural isomorphism
  \[
    \mathrm{pullbackSpecTildeComp}\,\varphi\,\varphi'
      : \widetilde{(-)} \ggg \mathrm{pullback}(\operatorname{Spec}\varphi)
          \ggg \mathrm{pullback}(\operatorname{Spec}\varphi')
        \;\cong\;
        \mathrm{ext}_{\varphi} \ggg \mathrm{ext}_{\varphi'} \ggg \widetilde{(-)},
  \]
  defined as
  \(\bigl(\mathrm{pullbackSpecTildeNatIso}\,\varphi \rhd
      \mathrm{pullback}(\operatorname{Spec}\varphi')\bigr) \ggg
    \bigl(\mathrm{ext}_{\varphi} \lhd \mathrm{pullbackSpecTildeNatIso}\,\varphi'\bigr)\). Each leg is a
  \emph{single} \(\operatorname{Spec}\)-pullback dictionary; the two \(\operatorname{Spec}\)-pullbacks
  appear only as the functor composite \(\ggg\), never folded into a single
  \(\mathrm{pullback}(\operatorname{Spec}(\varphi'\circ\varphi))\). Project-local.
  \begin{proof}
    \uses{lem:pullbackSpecTildeNatIso}
    Whisker the first dictionary on the right by \(\mathrm{pullback}(\operatorname{Spec}\varphi')\) and
    the second on the left by \(\mathrm{ext}_{\varphi}\), then compose; the source and target functors
    match by construction.
  \end{proof}
\end{lemma}

\begin{lemma}
[Multiplicativity of the affine pullback dictionary over composition]
  \label{lem:pullbackSpecTildeComp_eq_composite}
  \lean{AlgebraicGeometry.pullbackSpecTildeComp_eq_composite}
  \uses{lem:pullbackSpecTildeComp, lem:pullbackSpecTildeNatIso, lem:pullback_spec_tilde_iso,
        lem:gammaPushforwardNatIso_comp, lem:pullbackComp_mathlib, lem:extendScalarsComp_mathlib}
  The single-map dictionary \cref{lem:pullbackSpecTildeNatIso} is a \emph{pseudonatural} transformation
  of the \(\operatorname{Spec}\)-pullback and extend-scalars pseudofunctors: the two-leg fold
  \cref{lem:pullbackSpecTildeComp} equals the single-map dictionary of the \emph{composite} ring map
  \(\varphi'\circ\varphi\), conjugated by the pullback and extend-scalars composition coherences,
  \[
    \mathrm{pullbackSpecTildeComp}\,\varphi\,\varphi'
      = \bigl(\widetilde{(-)} \lhd \mathrm{pullbackComp}(\operatorname{Spec}\varphi',
            \operatorname{Spec}\varphi)\bigr)
        \;\ggg\; \mathrm{pullbackSpecTildeNatIso}(\varphi'\circ\varphi)
        \;\ggg\; \bigl(\mathrm{extendScalarsComp}\,\varphi\,\varphi' \rhd \widetilde{(-)}\bigr).
  \]
  The right-hand side names a \emph{single} \(\operatorname{Spec}\)-pullback (of the composite ring
  map) and uses \(\mathrm{pullbackComp}\) \cref{lem:pullbackComp_mathlib} only as an abstract natural
  iso; no \(\operatorname{Spec}\)-pullback of a composite is ever evaluated on an object. This is the
  single place the geometric pullback coherence enters the re-scoped route.
  \begin{proof}
    \uses{lem:pullbackSpecTildeNatIso, lem:pullback_spec_tilde_iso, lem:gammaPushforwardNatIso_comp,
          lem:pullbackComp_mathlib, lem:extendScalarsComp_mathlib}
    By construction \cref{lem:pullbackSpecTildeNatIso} each dictionary
    \(\mathrm{pullbackSpecTildeNatIso}\,\varphi\) is the conjugate (\(\mathrm{conjugateIsoEquiv}\)) of
    the base comparison \(\mathrm{gammaPushforwardNatIso}\,\varphi\) over the composite
    \(\widetilde{(-)}\dashv\Gamma\) / \(\mathrm{pullback}\dashv\mathrm{pushforward}\) and
    \(\mathrm{ext}\dashv\mathrm{restr}\) / \(\widetilde{(-)}\dashv\Gamma\) adjunctions. The base
    comparison is multiplicative across composites: \(\mathrm{gammaPushforwardNatIso}(\varphi'\circ
    \varphi)\) is the composite of \(\mathrm{gammaPushforwardNatIso}\,\varphi\) and
    \(\mathrm{gammaPushforwardNatIso}\,\varphi'\) glued by the pushforward composition coherence
    (\cref{lem:gammaPushforwardNatIso_comp}). Transporting that composition identity through the
    conjugation turns the pushforward/pullback composition coherence into the
    \(\mathrm{pullbackComp}\) conjugator on the left and the algebraic composition coherence
    \cref{lem:extendScalarsComp_mathlib} into the \(\mathrm{extendScalarsComp}\) conjugator on the
    right, exhibiting the two-leg fold as the displayed conjugate of the composite dictionary. The
    whole argument is at the natural-transformation level, so \(\mathrm{pullbackComp}\) is never
    evaluated on an object.
  \end{proof}
\end{lemma}

% ----------------------------------------------------------------------------
% (L1a--L1d) FINE per-leg single-map rewrites — iter-032 effort-breaker.
% Each is the M-component of a SINGLE whiskered single-map dictionary
% `pullbackSpecTildeNatIso φᵢ` (one ring map, one Spec-pullback). No composite
% Spec-pullback, hence no `Γ(pullbackComp)` of a composite, is ever formed in
% these statements — they are the smallest object-level rewrite pieces of the
% assembly, each elaborating well inside the 200k per-declaration budget. They
% are the definitional factors of the two-leg folds
% \cref{lem:pullbackSpecTildeComp}: L1a;L1b is the ι_R;ρ_B routing fold and
% L1c;L1d is the ρ;ι_{R'} routing fold.
% ----------------------------------------------------------------------------

\begin{lemma}
[(L1a) Per-leg rewrite of the \(\iota_R\) geometric leg]
  \label{lem:chartBaseChange_geomLeg_iotaR_app}
  \lean{AlgebraicGeometry.chartBaseChange_geomLeg_iotaR_app}
  % planned, not yet in Lean
  \uses{lem:pullbackSpecTildeNatIso, lem:pullbackSpecTildeNatIso_app, lem:pullback_spec_tilde_iso,
        lem:ring_square_glue_whiskerRight_lift}
  For every \(R\)-module \(M\), the \(M\)-component of the single-map dictionary
  \cref{lem:pullbackSpecTildeNatIso} for \(\iota_R\), right-whiskered by
  \(\mathrm{pullback}(\operatorname{Spec}\rho_B)\), is the outer pullback functor applied to the
  per-module dictionary \cref{lem:pullback_spec_tilde_iso} of \(\iota_R\):
  \[
    \bigl(\mathrm{pullbackSpecTildeNatIso}\,\iota_R \rhd
      \mathrm{pullback}(\operatorname{Spec}\rho_B)\bigr).\mathrm{app}\,M
      = \mathrm{pullback}(\operatorname{Spec}\rho_B).\mathrm{map}
        \bigl(\mathrm{pullback\_spec\_tilde\_iso}\,\iota_R\,M\bigr).
  \]
  Only the \emph{single} \(\operatorname{Spec}\)-pullback of \(\iota_R\) and one outer pullback functor
  occur; no composite pullback coherence is built. Project-local.
  \begin{proof}
    \uses{lem:ring_square_glue_whiskerRight_lift, lem:pullbackSpecTildeNatIso_app}
    By the right-whiskering component formula \cref{lem:ring_square_glue_whiskerRight_lift},
    \((\Phi \rhd H).\mathrm{app}\,M = H.\mathrm{map}(\Phi.\mathrm{app}\,M)\); apply it with
    \(\Phi = \mathrm{pullbackSpecTildeNatIso}\,\iota_R\) and
    \(H = \mathrm{pullback}(\operatorname{Spec}\rho_B)\), then rewrite the component by
    \cref{lem:pullbackSpecTildeNatIso_app}.
  \end{proof}
\end{lemma}

\begin{lemma}
[(L1b) Per-leg rewrite of the \(\rho_B\) geometric leg]
  \label{lem:chartBaseChange_geomLeg_rhoB_app}
  \lean{AlgebraicGeometry.chartBaseChange_geomLeg_rhoB_app}
  % planned, not yet in Lean
  \uses{lem:pullbackSpecTildeNatIso, lem:pullbackSpecTildeNatIso_app, lem:pullback_spec_tilde_iso,
        lem:ring_square_glue_whiskerLeft_lift}
  For every \(R\)-module \(M\), the \(M\)-component of the single-map dictionary
  \cref{lem:pullbackSpecTildeNatIso} for \(\rho_B\), left-whiskered by \(\mathrm{ext}_{\iota_R}\), is the
  per-module dictionary \cref{lem:pullback_spec_tilde_iso} of \(\rho_B\) at the base-changed module:
  \[
    \bigl(\mathrm{ext}_{\iota_R} \lhd \mathrm{pullbackSpecTildeNatIso}\,\rho_B\bigr).\mathrm{app}\,M
      = \mathrm{pullback\_spec\_tilde\_iso}\,\rho_B\,\bigl(\mathrm{ext}_{\iota_R} M\bigr).
  \]
  A \emph{single} \(\operatorname{Spec}\)-pullback (of \(\rho_B\)) only; no composite. Project-local.
  \begin{proof}
    \uses{lem:ring_square_glue_whiskerLeft_lift, lem:pullbackSpecTildeNatIso_app}
    By the left-whiskering component formula \cref{lem:ring_square_glue_whiskerLeft_lift},
    \((F \lhd \Phi).\mathrm{app}\,M = \Phi.\mathrm{app}\,(F\,M)\); apply it with
    \(F = \mathrm{ext}_{\iota_R}\) and \(\Phi = \mathrm{pullbackSpecTildeNatIso}\,\rho_B\), then rewrite
    the component by \cref{lem:pullbackSpecTildeNatIso_app}.
  \end{proof}
\end{lemma}

\begin{lemma}
[(L1c) Per-leg rewrite of the \(\rho\) geometric leg]
  \label{lem:chartBaseChange_geomLeg_rho_app}
  \lean{AlgebraicGeometry.chartBaseChange_geomLeg_rho_app}
  % planned, not yet in Lean
  \uses{lem:pullbackSpecTildeNatIso, lem:pullbackSpecTildeNatIso_app, lem:pullback_spec_tilde_iso,
        lem:ring_square_glue_whiskerRight_lift}
  For every \(R\)-module \(M\), the \(M\)-component of the single-map dictionary
  \cref{lem:pullbackSpecTildeNatIso} for \(\rho\), right-whiskered by
  \(\mathrm{pullback}(\operatorname{Spec}\iota_{R'})\), is the outer pullback functor applied to the
  per-module dictionary \cref{lem:pullback_spec_tilde_iso} of \(\rho\):
  \[
    \bigl(\mathrm{pullbackSpecTildeNatIso}\,\rho \rhd
      \mathrm{pullback}(\operatorname{Spec}\iota_{R'})\bigr).\mathrm{app}\,M
      = \mathrm{pullback}(\operatorname{Spec}\iota_{R'}).\mathrm{map}
        \bigl(\mathrm{pullback\_spec\_tilde\_iso}\,\rho\,M\bigr).
  \]
  Only the \emph{single} \(\operatorname{Spec}\)-pullback of \(\rho\) and one outer pullback functor
  occur. Project-local.
  \begin{proof}
    \uses{lem:ring_square_glue_whiskerRight_lift, lem:pullbackSpecTildeNatIso_app}
    As \cref{lem:chartBaseChange_geomLeg_iotaR_app}, with the right-whiskering component formula
    \cref{lem:ring_square_glue_whiskerRight_lift} for
    \(\Phi = \mathrm{pullbackSpecTildeNatIso}\,\rho\),
    \(H = \mathrm{pullback}(\operatorname{Spec}\iota_{R'})\), then
    \cref{lem:pullbackSpecTildeNatIso_app}.
  \end{proof}
\end{lemma}

\begin{lemma}
[(L1d) Per-leg rewrite of the \(\iota_{R'}\) geometric leg]
  \label{lem:chartBaseChange_geomLeg_iotaRPrime_app}
  \lean{AlgebraicGeometry.chartBaseChange_geomLeg_iotaRPrime_app}
  % planned, not yet in Lean
  \uses{lem:pullbackSpecTildeNatIso, lem:pullbackSpecTildeNatIso_app, lem:pullback_spec_tilde_iso,
        lem:ring_square_glue_whiskerLeft_lift}
  For every \(R\)-module \(M\), the \(M\)-component of the single-map dictionary
  \cref{lem:pullbackSpecTildeNatIso} for \(\iota_{R'}\), left-whiskered by \(\mathrm{ext}_{\rho}\), is the
  per-module dictionary \cref{lem:pullback_spec_tilde_iso} of \(\iota_{R'}\) at the base-changed module:
  \[
    \bigl(\mathrm{ext}_{\rho} \lhd \mathrm{pullbackSpecTildeNatIso}\,\iota_{R'}\bigr).\mathrm{app}\,M
      = \mathrm{pullback\_spec\_tilde\_iso}\,\iota_{R'}\,\bigl(\mathrm{ext}_{\rho} M\bigr).
  \]
  A \emph{single} \(\operatorname{Spec}\)-pullback (of \(\iota_{R'}\)) only; no composite. Project-local.
  \begin{proof}
    \uses{lem:ring_square_glue_whiskerLeft_lift, lem:pullbackSpecTildeNatIso_app}
    As \cref{lem:chartBaseChange_geomLeg_rhoB_app}, with the left-whiskering component formula
    \cref{lem:ring_square_glue_whiskerLeft_lift} for \(F = \mathrm{ext}_{\rho}\),
    \(\Phi = \mathrm{pullbackSpecTildeNatIso}\,\iota_{R'}\), then
    \cref{lem:pullbackSpecTildeNatIso_app}.
  \end{proof}
\end{lemma}

\begin{lemma}
[(L2) Algebraic-cocycle pivot of the two routings at \(M\)]
  \label{lem:chartBaseChange_algLeg_cocycle_app}
  \lean{AlgebraicGeometry.chartBaseChange_algLeg_cocycle_app}
  % planned, not yet in Lean
  \uses{lem:chartBaseChangeModuleReassoc_eq_natApp, lem:chartBaseChangeModuleReassoc,
        lem:chartBaseChange_extendScalars_cocycle, lem:ring_square_glue_whiskerRight_lift,
        lem:functor_map_comp_mathlib}
  After the four per-leg rewrites \cref{lem:chartBaseChange_geomLeg_iotaR_app},
  \cref{lem:chartBaseChange_geomLeg_rhoB_app}, \cref{lem:chartBaseChange_geomLeg_rho_app},
  \cref{lem:chartBaseChange_geomLeg_iotaRPrime_app} the residual identification of the two routings is
  purely algebraic and is the single sheafification of the base-change cocycle: for every \(R\)-module
  \(M\),
  \[
    \bigl(\mathrm{chartBaseChangeModuleReassocNat} \rhd \widetilde{(-)}\bigr).\mathrm{app}\,M
      = \widetilde{(-)}\bigl(\mathrm{chartBaseChangeModuleReassoc}\,M\bigr),
  \]
  and \(\mathrm{chartBaseChangeModuleReassoc}\,M\) is the proven extend-scalars cocycle
  \cref{lem:chartBaseChange_extendScalars_cocycle} identifying the two bracketings
  \(\mathrm{ext}_{\rho_B}(\mathrm{ext}_{\iota_R}M)\cong\mathrm{ext}_{\iota_{R'}}(\mathrm{ext}_{\rho}M)\).
  Every carrier is a light \(\mathrm{ModuleCat}\); no sheafification beyond the single outer
  \(\widetilde{(-)}\) and no pullback functor is forced to weak-head normal form. Project-local.
  \begin{proof}
    \uses{lem:ring_square_glue_whiskerRight_lift, lem:chartBaseChangeModuleReassoc_eq_natApp,
          lem:chartBaseChange_extendScalars_cocycle}
    The right-whiskering component formula \cref{lem:ring_square_glue_whiskerRight_lift} with
    \(\Phi = \mathrm{chartBaseChangeModuleReassocNat}\), \(H = \widetilde{(-)}\) gives
    \((\mathrm{chartBaseChangeModuleReassocNat} \rhd \widetilde{(-)}).\mathrm{app}\,M
      = \widetilde{(-)}(\mathrm{chartBaseChangeModuleReassocNat}.\mathrm{app}\,M)\); rewrite the
    component by the bridge \cref{lem:chartBaseChangeModuleReassoc_eq_natApp}. The reassociation is the
    algebraic cocycle \cref{lem:chartBaseChange_extendScalars_cocycle}, whose carriers are extend-scalars
    \(\mathrm{ModuleCat}\)s, so the equality stays purely algebraic.
  \end{proof}
\end{lemma}

\begin{lemma}
[(b2 master) Geometric comparison transported through the dictionary equals the algebraic reassociation]
  \label{lem:chartBaseChangeGeometricComparison_dictionary_glue}
  \lean{AlgebraicGeometry.chartBaseChangeGeometricComparison_dictionary_glue}
  % NOTE (iter-032 effort-breaker): OPTIONAL AGGREGATE — no longer on the critical path of the
  % realization \cref{lem:pullback_spec_tilde_iso_ring_square_natural_extendScalarsCocycle}, which now
  % assembles the per-leg pieces L1a--L1d + L2 directly. This monolithic functor-level equation is the
  % residual 200k-budget RISK (the historical `ring_square_glue_natTrans' elaborated only at 800k); it
  % is retained as a packaged functor-level statement but a prover should prefer the fine pieces.
  \uses{lem:pullbackSpecTildeComp, lem:pullbackSpecTildeComp_eq_composite,
        lem:chartBaseChangeGeometricComparisonNat, lem:chartBaseChangeGeometricComparisonNat_unfold,
        lem:chartBaseChangeModuleReassoc_eq_natApp, lem:chartBaseChange_extendScalars_cocycle,
        lem:chartBaseChange_ring_square}
  The functor-level master identity of the re-scoped route: the whiskered geometric comparison,
  followed by the two-leg dictionary fold of the \(\rho;\iota_{R'}\) routing, equals the two-leg
  dictionary fold of the \(\iota_R;\rho_B\) routing, followed by the whiskered algebraic
  reassociation,
  \[
    \bigl(\widetilde{(-)} \lhd \mathrm{chartBaseChangeGeometricComparisonNat}\bigr)
      \;\ggg\; \mathrm{pullbackSpecTildeComp}\,\rho\,\iota_{R'}
      \;=\;
    \mathrm{pullbackSpecTildeComp}\,\iota_R\,\rho_B
      \;\ggg\; \bigl(\mathrm{chartBaseChangeModuleReassocNat} \rhd \widetilde{(-)}\bigr),
  \]
  an equality of natural isomorphisms
  \(\widetilde{(-)} \ggg \mathrm{pullback}(\operatorname{Spec}\iota_R) \ggg
    \mathrm{pullback}(\operatorname{Spec}\rho_B) \cong \mathrm{ext}_{\rho} \ggg \mathrm{ext}_{\iota_{R'}}
    \ggg \widetilde{(-)}\). Project-local.
  \begin{proof}
    \uses{lem:pullbackSpecTildeComp_eq_composite, lem:chartBaseChangeGeometricComparisonNat_unfold,
          lem:chartBaseChange_extendScalars_cocycle, lem:chartBaseChange_ring_square}
    Apply the multiplicativity \cref{lem:pullbackSpecTildeComp_eq_composite} to both routings. Each
    two-leg fold becomes the single-map dictionary of its composite ring map conjugated by
    \(\mathrm{pullbackComp}\) and \(\mathrm{extendScalarsComp}\): the \(\iota_R;\rho_B\) routing folds
    through \(\mathrm{pullbackSpecTildeNatIso}(\rho_B\circ\iota_R)\) and the \(\rho;\iota_{R'}\) routing
    through \(\mathrm{pullbackSpecTildeNatIso}(\iota_{R'}\circ\rho)\). The two composite ring maps are
    \emph{equal}, \(\rho_B\circ\iota_R = \iota_{R'}\circ\rho\), by the base-change ring square
    \cref{lem:chartBaseChange_ring_square}; the geometric comparison, unfolded by
    \cref{lem:chartBaseChangeGeometricComparisonNat_unfold} as
    \(\mathrm{pullbackComp} \ggg \mathrm{pullbackCongr}(\mathrm{hsq}) \ggg \mathrm{pullbackComp}^{-1}\),
    is precisely the abstract \(\mathrm{pullbackComp}\)/\(\mathrm{pullbackCongr}\) reassociation
    matching those two equal composite \(\operatorname{Spec}\)-pullbacks, so it cancels the two
    \(\mathrm{pullbackComp}\) conjugators against the common single-map dictionary core. What remains
    is the difference of the two \(\mathrm{extendScalarsComp}\) bracketings of the equal composite ring
    map, which is exactly the algebraic reassociation \(\mathrm{chartBaseChangeModuleReassocNat}\)
    realized natively as the extend-scalars cocycle \cref{lem:chartBaseChange_extendScalars_cocycle}
    (its \(M\)-component being \cref{lem:chartBaseChangeModuleReassoc_eq_natApp}). No carrier is
    evaluated: the geometric comparison and the dictionaries enter only through their abstract
    natural-iso unfoldings, so \(\mathrm{pullbackComp}\) is never forced to weak-head normal form.
  \end{proof}
\end{lemma}

\begin{lemma}

\leanok
[(Re-scoped) extendScalars-cocycle realization of the ring-square naturality]
  \label{lem:pullback_spec_tilde_iso_ring_square_natural_extendScalarsCocycle}
  \lean{AlgebraicGeometry.pullback_spec_tilde_iso_ring_square_natural_extendScalarsCocycle}
  \uses{lem:chartBaseChangeGeometricComparisonNat_unfold, lem:chartBaseChange_geomLeg_iotaR_app,
        lem:chartBaseChange_geomLeg_rhoB_app, lem:chartBaseChange_geomLeg_rho_app,
        lem:chartBaseChange_geomLeg_iotaRPrime_app, lem:chartBaseChange_algLeg_cocycle_app,
        lem:functor_map_comp_mathlib, lem:pullbackSpecTildeComp_eq_composite,
        lem:pullbackSpecTildeNatIso_app, lem:pullback_spec_tilde_iso,
        lem:chartBaseChangeGeometricComparison, lem:chartBaseChange_extendScalars_cocycle}
  % NOTE (iter-032 effort-breaker): FINE decomposition landed. The realization now assembles the
  % per-leg single-map rewrites L1a--L1d (\cref{lem:chartBaseChange_geomLeg_iotaR_app},
  % \cref{lem:chartBaseChange_geomLeg_rhoB_app}, \cref{lem:chartBaseChange_geomLeg_rho_app},
  % \cref{lem:chartBaseChange_geomLeg_iotaRPrime_app}) + the algebraic-cocycle pivot L2
  % \cref{lem:chartBaseChange_algLeg_cocycle_app} + the controlled-unfold L0
  % \cref{lem:chartBaseChangeGeometricComparisonNat_unfold}, DIRECTLY — it no longer routes through the
  % monolithic functor-level master glue \cref{lem:chartBaseChangeGeometricComparison_dictionary_glue}
  % (kept as an optional aggregate; it is the residual 200k-budget risk, not on this critical path).
  % Lean declaration EXISTS as a typed-sorry scaffold (statement = crux conclusion verbatim); the proof
  % fill is the iter-032+ target. This is the live discharge of the crux
  % \cref{lem:pullback_spec_tilde_iso_ring_square_natural}; it supersedes the refuted mate, vComp, and
  % concrete-section routes of the following subsection.
  % SOURCE: [Stacks Project], Schemes, Lemma `lemma-widetilde-pullback', Tag 01I9
  % (\S\,Quasi-coherent sheaves on affines), part (1), the word "functorially"
  % (read from references/stacks-schemes.tex, L1241--1256).
  % SOURCE QUOTE: "\begin{enumerate}
  % \item We have $\psi^* \widetilde M = \widetilde{S \otimes_R M}$
  % functorially in the $R$-module $M$.
  % \item We have $\psi_* \widetilde N = \widetilde{N_R}$ functorially
  % in the $S$-module $N$.
  % \end{enumerate}"
  \textit{Source: Stacks Project, Schemes, Lemma ``\(\widetilde M\) and pullback'' (Tag 01I9),
  part~(1), the word ``functorially'' — read as the naturality of the per-map tilde-transport
  dictionary, closed by algebraic base-change associativity.}
  The ring-square naturality \cref{lem:pullback_spec_tilde_iso_ring_square_natural} holds, realized
  by the re-scoped extend-scalars cocycle, assembled \emph{leg by leg} from small kernel-light pieces.
  Each geometric leg is rewritten singly by the proven single-map dictionary
  \cref{lem:pullbackSpecTildeNatIso}/\cref{lem:pullback_spec_tilde_iso} —
  \(\widetilde{(-)} \ggg \mathrm{pullback}(\operatorname{Spec}\varphi_i) \cong \mathrm{ext}_{\varphi_i}
  \ggg \widetilde{(-)}\) for a \emph{single} ring map \(\varphi_i\), never a composite — via the four
  per-leg rewrites L1a--L1d (\cref{lem:chartBaseChange_geomLeg_iotaR_app},
  \cref{lem:chartBaseChange_geomLeg_rhoB_app}, \cref{lem:chartBaseChange_geomLeg_rho_app},
  \cref{lem:chartBaseChange_geomLeg_iotaRPrime_app}); the geometric comparison enters only through its
  controlled unfold \cref{lem:chartBaseChangeGeometricComparisonNat_unfold} and the abstract
  pseudonaturality \cref{lem:pullbackSpecTildeComp_eq_composite} (the one place
  \(\mathrm{pullbackComp}\) appears, as an abstract natural isomorphism); and the residual is the
  algebraic cocycle pivot L2 \cref{lem:chartBaseChange_algLeg_cocycle_app}, one outer
  \(\widetilde{(-)}\) of the proven extend-scalars cocycle
  \cref{lem:chartBaseChange_extendScalars_cocycle} distributed by functoriality
  \cref{lem:functor_map_comp_mathlib}. No flatness and no adjoint mate enter, and the pullback
  coherence \(\mathrm{pullbackComp}\) is present only \emph{abstractly} (it \emph{is} built — inside the
  geometric comparison — but is never unfolded / forced to weak-head normal form on an object: the dead
  \(\_concrete\) wall is exactly that object-level \(\Gamma(\mathrm{pullbackComp})\) evaluation). The
  live constraint the fine decomposition exists to beat is the \(200\)k-heartbeat per-declaration
  elaboration budget: the monolithic functor-level glue elaborated only at \(800\)k, so each piece here
  is sized to clear \(200\)k on its own.
  % SOURCE QUOTE PROOF: "The first assertion follows from the identification in
  % Lemma \ref{lemma-compare-constructions}
  % and the result of Modules, Lemma \ref{modules-lemma-restrict-quasi-coherent}."
  \begin{proof}
    \uses{lem:chartBaseChangeGeometricComparisonNat_unfold, lem:chartBaseChange_geomLeg_iotaR_app,
          lem:chartBaseChange_geomLeg_rhoB_app, lem:chartBaseChange_geomLeg_rho_app,
          lem:chartBaseChange_geomLeg_iotaRPrime_app, lem:chartBaseChange_algLeg_cocycle_app,
          lem:functor_map_comp_mathlib, lem:pullbackSpecTildeComp_eq_composite,
          lem:pullbackSpecTildeNatIso_app, lem:pullback_spec_tilde_iso,
          lem:chartBaseChangeGeometricComparison, lem:chartBaseChange_extendScalars_cocycle}
    \textit{Source: Stacks Project, Tag 01I9, proof of part~(1).}
    Assemble the fine pieces directly, never materializing the monolithic master glue. Present the
    geometric leg as \((\mathrm{chartBaseChangeGeometricComparisonNat}).\mathrm{app}\,\widetilde M\)
    (\cref{lem:chartBaseChangeGeometricComparison}), kept opaque via the controlled unfold
    \cref{lem:chartBaseChangeGeometricComparisonNat_unfold}. Rewrite the four legs singly by L1a--L1d
    (\cref{lem:chartBaseChange_geomLeg_iotaR_app}, \cref{lem:chartBaseChange_geomLeg_rhoB_app},
    \cref{lem:chartBaseChange_geomLeg_rho_app}, \cref{lem:chartBaseChange_geomLeg_iotaRPrime_app}), each
    an explicit \(\mathrm{pullback\_spec\_tilde\_iso}\,\varphi_i\) (\cref{lem:pullbackSpecTildeNatIso_app},
    \cref{lem:pullback_spec_tilde_iso}) of a single \(\operatorname{Spec}\)-pullback; match the two
    routings across the comparison by the abstract pseudonaturality
    \cref{lem:pullbackSpecTildeComp_eq_composite}. The residual is L2
    \cref{lem:chartBaseChange_algLeg_cocycle_app} — one \(\widetilde{(-)}\) of the cocycle
    \cref{lem:chartBaseChange_extendScalars_cocycle} distributed by
    \cref{lem:functor_map_comp_mathlib}; an iso cancellation of the trailing
    \(\rho;\iota_{R'}\)-dictionary factor gives the conclusion. Only \(\mathrm{.app}\,M\) of abstract
    functor isos is taken, so no composite pullback coherence is forced to weak-head normal form.
  \end{proof}
\end{lemma}

\subsection{Refuted routes for the ring-square crux (slated for excision)}

% NOTE: The two route blocks of this subsection (the functor-paste / vComp route
% \cref{lem:pullback_spec_tilde_iso_ring_square_natural_catCommSq} and the concrete-section route
% \cref{lem:pullback_spec_tilde_iso_ring_square_natural_concrete}), together with the abstract mate
% cocycle \cref{lem:ring_square_cocycle} and its glue
% \cref{lem:pullback_spec_tilde_iso_ring_square_mate_glue} recorded above, are SUPERSEDED by the
% re-scoped extend-scalars cocycle route
% \cref{lem:pullback_spec_tilde_iso_ring_square_natural_extendScalarsCocycle} of the preceding
% subsection. They are retained only because Lean still references their declarations; they are
% slated for excision once the re-scoped scaffold lands.
The blocks below record the routes that were attempted for the ring-square crux
\cref{lem:pullback_spec_tilde_iso_ring_square_natural} before it was re-scoped. Each is walled at a
heavy-carrier reduction — the functor-paste route at the \(\mathrm{vComp}\)-glue mate and the
component extraction at \(M\), the concrete-section route at the evaluation of
\(\Gamma(\mathrm{pullbackComp})\) for which Mathlib supplies no
\((\operatorname{Spec}\varphi)^{*}\widetilde M \cong \widetilde{\mathrm{ext}_\varphi M}\) reduction.
They are \emph{superseded} by the re-scoped route
\cref{lem:pullback_spec_tilde_iso_ring_square_natural_extendScalarsCocycle} and slated for excision;
they are kept here only because the Lean development still names their declarations. We first record
the \(\mathrm{Cat}\)-level square algebra that the functor-paste route consumes.

\begin{lemma}[Commuting square of categories]
  \label{lem:catCommSq_mathlib}
  \lean{CategoryTheory.CatCommSq}
  \mathlibok
  \textit{Provided by Mathlib.}
  For functors \(T : \mathcal{C}_1 \to \mathcal{C}_2\), \(V_1 : \mathcal{C}_1 \to \mathcal{C}_4\),
  \(V_2 : \mathcal{C}_2 \to \mathcal{C}_5\) and \(B : \mathcal{C}_4 \to \mathcal{C}_5\), the structure
  \(\mathrm{CatCommSq}\,T\,V_1\,V_2\,B\) packages a chosen natural isomorphism
  \(T \ggg V_2 \cong V_1 \ggg B\) filling the square of categories. It is the data of a commuting
  square at the level of functors, matched on the named functors at its corners rather than on the
  objects they produce.
\end{lemma}

\begin{lemma}[Horizontal pasting of commuting squares of categories]
  \label{lem:catCommSq_hComp_mathlib}
  \lean{CategoryTheory.CatCommSq.hComp}
  \mathlibok
  \uses{lem:catCommSq_mathlib}
  \textit{Provided by Mathlib.}
  Given \(\mathrm{CatCommSq}\,T_1\,V_1\,V_2\,B_1\) and \(\mathrm{CatCommSq}\,T_2\,V_2\,V_3\,B_2\)
  sharing the middle vertical functor \(V_2\), their horizontal paste is a
  \(\mathrm{CatCommSq}\,(T_1 \ggg T_2)\,V_1\,V_3\,(B_1 \ggg B_2)\). The composition adjacency is the
  \emph{shared functor} \(V_2\), matched by congruence of functor composition on named functors — no
  identification of the underlying objects is performed.
\end{lemma}

\begin{lemma}[Vertical pasting of commuting squares of categories]
  \label{lem:catCommSq_vComp_mathlib}
  \lean{CategoryTheory.CatCommSq.vComp}
  \mathlibok
  \uses{lem:catCommSq_mathlib}
  \textit{Provided by Mathlib.}
  Given \(\mathrm{CatCommSq}\,H_1\,L_1\,R_1\,H_2\) and \(\mathrm{CatCommSq}\,H_2\,L_2\,R_2\,H_3\)
  sharing the middle horizontal functor \(H_2\), their vertical paste is a
  \(\mathrm{CatCommSq}\,H_1\,(L_1 \ggg L_2)\,(R_1 \ggg R_2)\,H_3\). As with the horizontal paste, the
  adjacency is the shared named functor \(H_2\).
\end{lemma}

\begin{lemma}

\leanok
[(Route A) Functor-paste realization of the ring-square naturality]
  \label{lem:pullback_spec_tilde_iso_ring_square_natural_catCommSq}
  \lean{AlgebraicGeometry.pullback_spec_tilde_iso_ring_square_natural_catCommSq}
  % NOTE (refuted route, slated for excision): SUPERSEDED by the re-scoped extend-scalars cocycle
  % \cref{lem:pullback_spec_tilde_iso_ring_square_natural_extendScalarsCocycle}; retained only
  % because Lean still references this declaration. Walls: the vComp-glue mate proof and the >350k-hb
  % `.app M` extraction.
  \uses{lem:pullbackSpecTildeNatIso, lem:chartBaseChangeGeometricComparisonNat,
        lem:chartBaseChangeModuleReassoc, lem:chartBaseChangeModuleReassoc_eq_natApp,
        lem:chartBaseChange_ring_square, lem:catCommSq_mathlib, lem:catCommSq_hComp_mathlib,
        lem:catCommSq_vComp_mathlib}
  The ring-square naturality \cref{lem:pullback_spec_tilde_iso_ring_square_natural} holds, realized by
  pasting commuting squares of categories. Each affine pullback dictionary
  \cref{lem:pullbackSpecTildeNatIso} for a ring map \(\varphi\) is the filling isomorphism
  \(\widetilde{(-)} \ggg \mathrm{pullback}(\operatorname{Spec}\varphi) \cong
  \mathrm{ext}_\varphi \ggg \widetilde{(-)}\) of a commuting square of categories
  \cref{lem:catCommSq_mathlib}, with horizontal edges the two sheafification functors and vertical
  edges \(\mathrm{ext}_\varphi\) and \(\mathrm{pullback}(\operatorname{Spec}\varphi)\). Registering
  the four dictionaries for the ring maps \(\iota_R,\rho,\rho_B,\iota_{R'}\) of the base-change ring
  square \cref{lem:chartBaseChange_ring_square}, together with the geometric comparison
  \cref{lem:chartBaseChangeGeometricComparisonNat} (the square of pullback functors) and the algebraic
  reassociation \cref{lem:chartBaseChangeModuleReassoc} (the square of extend-scalars functors) as
  commuting squares of categories, the two ways around the ring square assemble by horizontal and
  vertical pasting \cref{lem:catCommSq_hComp_mathlib}, \cref{lem:catCommSq_vComp_mathlib} into two
  commuting squares of categories with the \emph{same} corner functors. Their filling isomorphisms
  agree, and evaluating that equality of natural isomorphisms at a single module \(M\) — extracted once
  — is exactly the asserted ring-square naturality.
  \begin{proof}
    \uses{lem:pullbackSpecTildeNatIso, lem:chartBaseChangeGeometricComparisonNat,
          lem:chartBaseChangeModuleReassoc, lem:chartBaseChangeModuleReassoc_eq_natApp,
          lem:chartBaseChange_ring_square, lem:catCommSq_mathlib, lem:catCommSq_hComp_mathlib,
          lem:catCommSq_vComp_mathlib}
    Wrap each of the six legs as a commuting square of categories: the four dictionary squares
    \cref{lem:pullbackSpecTildeNatIso} (one per ring map), the geometric comparison square
    \cref{lem:chartBaseChangeGeometricComparisonNat}, and the algebraic reassociation square
    \cref{lem:chartBaseChangeModuleReassoc} (presented at the natural-transformation level through
    \cref{lem:chartBaseChangeModuleReassoc_eq_natApp}). Paste the two routings of the base-change ring
    square \cref{lem:chartBaseChange_ring_square}: the route through \(\iota_R\) then \(\rho_B\) and the
    route through \(\rho\) then \(\iota_{R'}\), each glued to the corresponding comparison square. By the
    pasting laws \cref{lem:catCommSq_hComp_mathlib} and \cref{lem:catCommSq_vComp_mathlib} both routings
    produce a commuting square of categories with identical corner functors — the corners are built by
    \emph{functor composition on named functors} (\(\widetilde{(-)}\),
    \(\mathrm{pullback}(\operatorname{Spec}-)\), \(\mathrm{ext}\)), so adjacency never forces an
    identification of the underlying heavy objects. The two pastes' filling isomorphisms therefore agree
    as natural isomorphisms; applying that equality at \(M\) once recovers the equation of
    \cref{lem:pullback_spec_tilde_iso_ring_square_natural}. The per-component formulas
    \(\mathrm{hComp\_iso\_hom\_app}\) / \(\mathrm{vComp\_iso\_hom\_app}\) are \emph{not} expanded against
    the goal, since their right-hand sides re-cross the heavy carriers; the extraction at \(M\) is a
    single structural evaluation of the already-established equality of fillings.
  \end{proof}
\end{lemma}

\begin{lemma}

\leanok
[(Route B) Concrete tilde realization of the ring-square naturality]
  \label{lem:pullback_spec_tilde_iso_ring_square_natural_concrete}
  \lean{AlgebraicGeometry.pullback_spec_tilde_iso_ring_square_natural_concrete}
  % NOTE (refuted route, slated for excision): SUPERSEDED by the re-scoped extend-scalars cocycle
  % \cref{lem:pullback_spec_tilde_iso_ring_square_natural_extendScalarsCocycle}; retained only
  % because Lean still references this declaration. Wall: the residual `tilde.preimage` of the
  % geometric `pullbackComp` leg has no Mathlib reduction, so this route is structurally dead.
  \uses{lem:base_change_chart_tensor_iso, lem:pullback_spec_tilde_iso, lem:chartBaseChange_ring_square,
        lem:gammaPushforwardIsoAt, lem:tildeRestriction_isLocalizedModule,
        lem:powers_restrictScalars, lem:cancelBaseChange_mathlib}
  % SOURCE: [Stacks Project], Schemes, Lemma `lemma-widetilde-pullback', Tag 01I9
  % (\S\,Quasi-coherent sheaves on affines) (read from references/stacks-schemes.tex,
  % L1241--1256).
  % SOURCE QUOTE: "Let
  % $(X, \mathcal{O}_X) = (\Spec(S), \mathcal{O}_{\Spec(S)})$,
  % $(Y, \mathcal{O}_Y) = (\Spec(R), \mathcal{O}_{\Spec(R)})$
  % be affine schemes.
  % Let $\psi : (X, \mathcal{O}_X) \to (Y, \mathcal{O}_Y)$ be a
  % morphism of affine schemes, corresponding to the ring map
  % $\psi^\sharp : R \to S$ (see Lemma \ref{lemma-category-affine-schemes}).
  % \begin{enumerate}
  % \item We have $\psi^* \widetilde M = \widetilde{S \otimes_R M}$
  % functorially in the $R$-module $M$.
  % \item We have $\psi_* \widetilde N = \widetilde{N_R}$ functorially
  % in the $S$-module $N$.
  % \end{enumerate}"
  \textit{Source: Stacks Project, Schemes, Lemma ``\(\widetilde M\) and pullback''
  (Tag 01I9), part~(1), the word ``functorially'' — read as the naturality of the
  tilde-transport isomorphism under restriction.}
  The ring-square naturality \cref{lem:pullback_spec_tilde_iso_ring_square_natural} holds, realized
  concretely in the tilde register without any adjoint-mate identification. By the affine pullback
  dictionary \cref{lem:pullback_spec_tilde_iso} (Stacks Tag 01I9, part~(1)) each leg is the
  tilde-transport isomorphism \((\operatorname{Spec}\varphi)^{*}\widetilde M \cong
  \widetilde{R' \otimes_R M}\), an isomorphism of \emph{honest module maps computed by localization};
  part~(1)'s assertion that this identification is \emph{functorial} in \(M\) is exactly the naturality
  that the square asserts. Concretely, restriction along the base-changed open immersion
  \(j_B : W_B \hookrightarrow V_B\) is, on global sections, localization of the dictionary at the
  base-change ring square \cref{lem:chartBaseChange_ring_square}; tracing both routings of that square
  through the per-chart isomorphisms \cref{lem:base_change_chart_tensor_iso} for \(V\) and for \(W\) and
  cancelling \((B \otimes_A R) \otimes_R M \cong B \otimes_A M\)
  (\cref{lem:cancelBaseChange_mathlib}), each sends a section \(s\) to \(\rho^{\#}(s) \otimes b\), so the
  two routings agree. This is the equation of
  \cref{lem:pullback_spec_tilde_iso_ring_square_natural}.
  % SOURCE QUOTE PROOF: "The first assertion follows from the identification in
  % Lemma \ref{lemma-compare-constructions}
  % and the result of Modules, Lemma \ref{modules-lemma-restrict-quasi-coherent}."
  \begin{proof}
    \uses{lem:base_change_chart_tensor_iso, lem:pullback_spec_tilde_iso, lem:chartBaseChange_ring_square,
          lem:gammaPushforwardIsoAt, lem:tildeRestriction_isLocalizedModule,
          lem:powers_restrictScalars, lem:cancelBaseChange_mathlib}
    \textit{Source: Stacks Project, Tag 01I9, proof of part~(1).}
    Work on a standard basic open \(D(a) \subseteq V\) (with \(a \in R\)); the basic opens form a basis,
    so it suffices to check the square there. Over \(D(a)\), the affine pullback dictionary
    \cref{lem:pullback_spec_tilde_iso} is, on sections, the open-indexed section comparison
    \cref{lem:gammaPushforwardIsoAt} evaluated at \(D(a)\), which identifies the pushforward/pullback
    sections with the localization \(M[1/\varphi a]\) read over \(R[1/a]\) by restriction of scalars; that
    the structure-sheaf restriction realizes this localization is
    \cref{lem:tildeRestriction_isLocalizedModule}, transported to the base ring by
    \cref{lem:powers_restrictScalars}. Each constituent of these isomorphisms is identity-on-carrier — a
    structure-sheaf restriction map or a restriction-of-scalars repackaging — so each leg of the square is
    a localization map. The base-change ring square \cref{lem:chartBaseChange_ring_square}
    \(\iota_R \mathbin{;} \rho_B = \rho \mathbin{;} \iota_{R'}\) localizes, on \(D(a)\), to the equality of
    the two routings; cancelling \((B \otimes_A R) \otimes_R M \cong B \otimes_A M\)
    \cref{lem:cancelBaseChange_mathlib} on each routing — the per-chart identification
    \cref{lem:base_change_chart_tensor_iso} — shows both send \(s\) to \(\rho^{\#}(s) \otimes b\). Hence the
    square commutes on every basic open, and therefore as the asserted equation of
    \cref{lem:pullback_spec_tilde_iso_ring_square_natural}. No flatness, no adjoint mate, and no abstract
    base-change map enter; every arrow is a concrete localization of modules.

    % NOTE: This route grounds the crux in Stacks Tag 01I9's "functorially" (part (1)): the
    % tilde-transport isomorphism is natural in the module, and restriction along an affine open
    % immersion localizes it. The Lean target named here is a planned declaration realizing this
    % concrete derivation; it bypasses the mate apparatus of
    % \cref{lem:pullback_spec_tilde_iso_ring_square_mate_glue} entirely.
  \end{proof}
\end{lemma}

\subsection{Restriction-naturality and the flat base-change pushforward theorem}

With the ring-square crux \cref{lem:pullback_spec_tilde_iso_ring_square_natural} in hand, we record
the restriction-naturality of the affine pullback dictionary it feeds, and assemble the global flat
base-change isomorphism for the pushforward.

\begin{lemma}\leanok
[Restriction-naturality of the affine pullback dictionary]
  \label{lem:pullback_tilde_restriction_natural}
  \lean{AlgebraicGeometry.pullback_spec_tilde_iso_restriction_naturality}
  \uses{lem:pullback_spec_tilde_iso, lem:base_change_chart_tensor_iso}
  % SOURCE: [Stacks Project], Schemes, Lemma `lemma-widetilde-pullback', Tag 01I9
  % (\S\,Quasi-coherent sheaves on affines), part (1), the word "functorially"
  % (read from references/stacks-schemes.tex, L1242--1256).
  % SOURCE QUOTE: "\item We have $\psi^* \widetilde M = \widetilde{S \otimes_R M}$
  % functorially in the $R$-module $M$."
  % SOURCE QUOTE PROOF: "The first assertion follows from the identification in
  % Lemma \ref{lemma-compare-constructions}
  % and the result of Modules, Lemma \ref{modules-lemma-restrict-quasi-coherent}."
  \textit{Source: Stacks Project, Schemes, Lemma ``\(\widetilde M\) and pullback''
  (Tag 01I9), part~(1) — the assertion that the identification is \emph{functorial},
  read here as naturality along an open immersion of affines.}
  Keep the data of Lemma~\ref{lem:base_change_chart_tensor_iso}, and let
  \(j : W \hookrightarrow V\) be an open immersion of \emph{affine} opens of \(X\)
  (\(V = \operatorname{Spec} R\), \(W = \operatorname{Spec} R'\), corresponding to the
  restriction ring map \(\rho : R = \Gamma(V,\mathcal{O}_X) \to R' = \Gamma(W,\mathcal{O}_X)\),
  with \(\rho^{\#}\) the section restriction
  \(\Gamma(V,\mathcal{F}) \to \Gamma(W,\mathcal{F})\)). Then the per-chart isomorphisms of
  Lemma~\ref{lem:base_change_chart_tensor_iso} for \(V\) and for \(W\) commute with
  restriction: the square
  \[
    \begin{array}{ccc}
      \Gamma\bigl(V_B,\mathcal{F}_B\bigr) & \xrightarrow{\ \operatorname{res}\ }
        & \Gamma\bigl(W_B,\mathcal{F}_B\bigr) \\[2pt]
      \cong\big\downarrow & & \big\downarrow\cong \\[2pt]
      \Gamma(V,\mathcal{F}) \otimes_A B & \xrightarrow{\ \rho^{\#}\otimes\,\mathrm{id}_B\ }
        & \Gamma(W,\mathcal{F}) \otimes_A B
    \end{array}
  \]
  commutes, where the top map is restriction along the base-changed open immersion
  \(j_B : W_B \hookrightarrow V_B\). Equivalently, the base change \((- \otimes_A B)\) of the
  fork restriction map \(\Gamma(V,\mathcal{F}) \to \Gamma(W,\mathcal{F})\) \emph{is} the
  restriction map \(\Gamma(V_B,\mathcal{F}_B) \to \Gamma(W_B,\mathcal{F}_B)\), transported
  through the two per-chart dictionaries.
  \begin{proof}
    \uses{lem:pullback_spec_tilde_iso, lem:base_change_chart_tensor_iso,
          lem:pullback_spec_tilde_iso_ring_square_natural,
          lem:chartBaseChangeGeometricComparison, lem:chartBaseChangeModuleReassoc}
    This is the \emph{concrete functoriality} of the explicit tilde-pullback isomorphism of
    Lemma~\ref{lem:pullback_spec_tilde_iso} — every arrow is an honest module map computed by
    localization. Restriction along
    \(j : W\hookrightarrow V\) is the ring map \(\rho : R\to R'\); its base change
    \(\rho\otimes_A B\) is the coordinate map of \(j_B : W_B\hookrightarrow V_B\) (base change
    preserves open immersions). The dictionary \((\operatorname{Spec}\varphi)^{*}\widetilde M
    \cong\widetilde{(B\otimes_A R)\otimes_R M}\) is built from the structural localization
    isomorphisms, hence natural in the base ring map: pulling back then restricting along
    \(j_B\) equals restricting along \(j\) then pulling back (on a standard open
    \(D(a)\subseteq V\) underlying \(W\) this is just localizing the dictionary at \(a\)).
    Tracing both legs through the cancellation
    \((B\otimes_A R)\otimes_R M\cong B\otimes_A M\)
    (Lemma~\ref{lem:base_change_chart_tensor_iso}), each sends \(s\) to
    \(\rho^{\#}(s)\otimes b\), so the square commutes. Naturality is in the ring map and
    module, so no flatness and no adjoint mate enter.
  \end{proof}
\end{lemma}

\begin{lemma}[Tensor commutes with a finite product]
  \label{lem:tensor_finite_product_comm}
  \lean{TensorProduct.piRight}
  \mathlibok
  \textit{Source: Mathlib, \texttt{TensorProduct.piRight}
  (\texttt{Mathlib.LinearAlgebra.TensorProduct.Pi}).}
  Let \(A \to B\) be a ring map, \(\iota\) a \emph{finite} index set, and \((M_i)_{i\in\iota}\)
  a family of \(A\)-modules. Then the canonical map is a \(B\)-linear isomorphism
  \[
    B \otimes_A \Bigl(\textstyle\prod_{i\in\iota} M_i\Bigr)
    \;\cong\;
    \textstyle\prod_{i\in\iota}\bigl(B \otimes_A M_i\bigr),
  \]
  and likewise for any finite index set \(\{(i,j)\}\). \textit{Provided by Mathlib.}
\end{lemma}

\begin{lemma}[Base change of the finite equalizer diagram]
  \label{lem:base_changed_equalizer_diagram}
  \lean{AlgebraicGeometry.baseChange_sheafConditionFork_tensorIso}
  \uses{lem:gamma_finite_equalizer, lem:base_change_chart_tensor_iso,
        lem:pullback_tilde_restriction_natural, lem:tensor_finite_product_comm}
  % NOTE: The three auxiliary lemmas require: separatedness of \(X\) (so that the
  % overlap charts \(U_{ij}\) are affine, enabling
  % \cref{lem:base_change_chart_tensor_iso,lem:pullback_tilde_restriction_natural});
  % finiteness of the cover index (so that \cref{lem:tensor_finite_product_comm} applies);
  % and quasi-coherence of \(\mathcal{F}\) (to write each restriction as a tilde module).
  % Flatness of \(A \to B\) is not used here; it enters only downstream at
  % \cref{lem:flat_base_change_separated}.
  % SOURCE: [Stacks Project], Tag 02KH ($i=0$ case), separated paragraph
  % (read from references/stacks-coherent.md; quoted verbatim in the
  % SOURCE QUOTE PROOF of thm:flat_base_change_pushforward below).
  % SOURCE QUOTE: "If $\mathcal{U}_B : X_B = (U_1)_B \cup \ldots \cup (U_t)_B$
  % we obtain by base change, then it is still the case that each $(U_i)_B$
  % is affine and that $X_B$ is separated. ... We have the following relation
  % between the {\v C}ech complexes
  % $\check{\mathcal{C}}^\bullet(\mathcal{U}_B, \mathcal{F}_B) =
  % \check{\mathcal{C}}^\bullet(\mathcal{U}, \mathcal{F}) \otimes_A B$
  % as follows from Lemma \ref{lemma-affine-base-change}."
  \textit{Source: Stacks Project, Tag 02KH (``Flat base change''), $i=0$ case,
  the relation \(\check{\mathcal{C}}^\bullet(\mathcal{U}_B,\mathcal{F}_B) =
  \check{\mathcal{C}}^\bullet(\mathcal{U},\mathcal{F}) \otimes_A B\).}
  Let \(A \to B\) be a ring map, \(X\) a separated scheme over \(A\) with finite
  affine cover \(\mathcal{U} = \{U_i\}\) (so each \(U_{ij}\) is affine), and
  \(\mathcal{F}\) quasi-coherent. Set \(X_B = X \times_{\operatorname{Spec} A}
  \operatorname{Spec} B\), \(\mathcal{F}_B\) the pullback, and
  \(\mathcal{U}_B = \{(U_i)_B\}\). Then \(\mathcal{U}_B\) is a finite affine cover of
  \(X_B\), and the sheaf-condition fork of Lemma~\ref{lem:gamma_finite_equalizer} for
  \((X_B, \mathcal{F}_B)\) is obtained from that for \((X, \mathcal{F})\) by applying
  \(- \otimes_A B\): there are compatible isomorphisms
  \[
    \Gamma\bigl((U_i)_B, \mathcal{F}_B\bigr) \cong \Gamma(U_i, \mathcal{F}) \otimes_A B,
    \qquad
    \Gamma\bigl((U_{ij})_B, \mathcal{F}_B\bigr) \cong \Gamma(U_{ij}, \mathcal{F}) \otimes_A B,
  \]
  intertwining the two restriction maps, so that the whole diagram
  \(\prod_i \Gamma((U_i)_B,\mathcal{F}_B) \rightrightarrows \prod_{i,j}
  \Gamma((U_{ij})_B,\mathcal{F}_B)\) is \((- \otimes_A B)\) applied to
  \(\prod_i \Gamma(U_i,\mathcal{F}) \rightrightarrows \prod_{i,j}
  \Gamma(U_{ij},\mathcal{F})\).
  \begin{proof}
    \uses{lem:gamma_finite_equalizer, lem:base_change_chart_tensor_iso,
          lem:pullback_tilde_restriction_natural, lem:tensor_finite_product_comm}
    Base change preserves affineness and open immersions, so each \((U_i)_B\) is affine and
    covers \(X_B\); as \(X\) is separated each \(U_{ij}\), hence \((U_{ij})_B\), is affine, so
    the fork of Lemma~\ref{lem:gamma_finite_equalizer} applies to \((X_B,\mathcal{F}_B)\). Now
    assemble three seams. \emph{(a)}~Lemma~\ref{lem:base_change_chart_tensor_iso} on each
    \(U_i\) and overlap \(U_{ij}\) gives the per-chart isomorphisms
    \(\Gamma((U_i)_B,\mathcal{F}_B)\cong\Gamma(U_i,\mathcal{F})\otimes_A B\), at the module
    level. \emph{(b)}~Lemma~\ref{lem:pullback_tilde_restriction_natural} on each
    \(U_{ij}\hookrightarrow U_i\) makes them commute with restriction, so they intertwine
    \(\mathrm{leftRes}\), \(\mathrm{rightRes}\) with their base changes. \emph{(c)}~The covers
    being finite, Lemma~\ref{lem:tensor_finite_product_comm} pushes \(-\otimes_A B\) through
    the products \(\prod_i\), \(\prod_{i,j}\). Thus the base-changed fork is term- and
    arrow-wise \((-\otimes_A B)\) of the original; passing to equalizers identifies the
    base-changed equalizer locus with \((-\otimes_A B)\) of the original.
  \end{proof}
\end{lemma}

\begin{lemma}\leanok
[Flat base change, separated case]
  \label{lem:flat_base_change_separated}
  \lean{AlgebraicGeometry.flatBaseChange_pushforward_isIso_of_isSeparated}
  \uses{lem:gamma_finite_equalizer, lem:base_changed_equalizer_diagram,
        lem:flat_preserves_equalizer_mathlib}
  % SOURCE: [Stacks Project], Tag 02KH ($i=0$ case), separated paragraph
  % (read from references/stacks-coherent.md; quoted verbatim in the
  % SOURCE QUOTE PROOF of thm:flat_base_change_pushforward below).
  % SOURCE QUOTE: "Since $A \to B$ is flat, the same thing remains true on
  % taking cohomology."
  \textit{Source: Stacks Project, Tag 02KH (``Flat base change''), $i=0$ case,
  separated paragraph.}
  Let \(A \to B\) be flat, \(X\) a separated, quasi-compact scheme over \(A\), and
  \(\mathcal{F}\) quasi-coherent. Then the comparison map
  \[
    \Gamma(X, \mathcal{F}) \otimes_A B \;\longrightarrow\; \Gamma(X_B, \mathcal{F}_B)
  \]
  is an isomorphism.
  \begin{proof}
    \uses{lem:gamma_finite_equalizer, lem:base_changed_equalizer_diagram,
          lem:flat_preserves_equalizer_mathlib}
    By Lemma~\ref{lem:gamma_finite_equalizer}, \(\Gamma(X,\mathcal{F})\) and
    \(\Gamma(X_B,\mathcal{F}_B)\) are the equalizers of the finite forks of the cover
    \(\mathcal{U}\) and its base change \(\mathcal{U}_B\). By
    Lemma~\ref{lem:base_changed_equalizer_diagram} the \(X_B\)-fork is the
    \(X\)-fork with \(- \otimes_A B\) applied. Since \(A \to B\) is flat,
    \(- \otimes_A B\) commutes with the formation of this finite equalizer
    (Lemma~\ref{lem:flat_preserves_equalizer_mathlib}); hence
    \(\Gamma(X,\mathcal{F}) \otimes_A B = \operatorname{eq}(\cdots) \otimes_A B
    \xrightarrow{\sim} \operatorname{eq}(\cdots \otimes_A B) = \Gamma(X_B,\mathcal{F}_B)\),
    and this iso is the comparison map.
  \end{proof}
\end{lemma}

\begin{lemma}\leanok
[Flat base change, Mayer--Vietoris reduction of the quasi-separated case]
  \label{lem:flat_base_change_mayer_vietoris}
  \lean{AlgebraicGeometry.flatBaseChange_pushforward_mayerVietoris}
  \uses{lem:flat_base_change_separated, lem:pullback_spec_tilde_iso,
        lem:flat_preserves_equalizer_mathlib, lem:finite_affine_cover_qcqs}
  % SOURCE: [Stacks Project], Tag 02KH ($i=0$ case), quasi-separated paragraph
  % (read from references/stacks-coherent.md; quoted verbatim in the
  % SOURCE QUOTE PROOF of thm:flat_base_change_pushforward below).
  % SOURCE QUOTE: "In case $X$ is quasi-separated, choose an affine open
  % covering $\mathcal{U} : X = U_1 \cup \ldots \cup U_t$. ... The reader who
  % wishes to avoid this spectral sequence can use Mayer-Vietoris and induction
  % on $t$ as in the proof of Lemma
  % \ref{lemma-quasi-coherence-higher-direct-images}."
  \textit{Source: Stacks Project, Tag 02KH (``Flat base change''), $i=0$ case,
  quasi-separated paragraph (Mayer--Vietoris induction on \(t\)).}
  Let \(A \to B\) be flat, \(X\) a quasi-compact, quasi-separated scheme over \(A\),
  and \(\mathcal{F}\) quasi-coherent. Then the comparison map
  \(\Gamma(X,\mathcal{F}) \otimes_A B \to \Gamma(X_B,\mathcal{F}_B)\) is an
  isomorphism.
  \begin{proof}
    \uses{lem:flat_base_change_separated, lem:pullback_spec_tilde_iso,
          lem:flat_preserves_equalizer_mathlib, lem:finite_affine_cover_qcqs}
    Choose a finite affine cover \(\mathcal{U} : X = U_1 \cup \dots \cup U_t\)
    (Lemma~\ref{lem:finite_affine_cover_qcqs}) and induct on \(t\). For \(t = 1\),
    \(X = U_1\) is affine and the comparison map is the concrete affine module base
    change \(\Gamma(X,\mathcal{F})\otimes_A B \cong \Gamma(X_B,\mathcal{F}_B)\) supplied
    directly by the tilde dictionary Lemma~\ref{lem:pullback_spec_tilde_iso} (Tag 01I9),
    no abstract base-change map required. For \(t > 1\) put \(V = U_t\)
    (affine) and \(W = U_1 \cup \dots \cup U_{t-1}\) (covered by \(t-1\) affines, hence
    quasi-compact and quasi-separated). Their intersection \(W \cap V =
    \bigcup_{i<t}(U_i \cap U_t)\) is covered by the \(t-1\) opens \(U_i \cap U_t\), each
    an open of the affine \(U_i\) and so quasi-compact and separated; thus \(W \cap V\)
    is separated and quasi-compact, and flat base change holds for it by the separated
    case Lemma~\ref{lem:flat_base_change_separated}. The two-member cover
    \(\{W, V\}\) of \(X\) gives the left-exact Mayer--Vietoris equalizer
    \[
      0 \to \Gamma(X,\mathcal{F}) \to
        \Gamma(W,\mathcal{F}) \oplus \Gamma(V,\mathcal{F}) \to \Gamma(W\cap V,\mathcal{F}),
    \]
    and likewise for \(X_B\) with \(\{W_B, V_B\}\). Flatness makes \(- \otimes_A B\)
    preserve this finite kernel/equalizer
    (Lemma~\ref{lem:flat_preserves_equalizer_mathlib}); since flat base change holds for
    \(V\) (affine), for \(W\) (inductive hypothesis, \(t-1\) members), and for
    \(W \cap V\) (separated), the five-lemma/equalizer comparison forces it for \(X\).
  \end{proof}
\end{lemma}

\begin{lemma}\leanok
[Reduction of the base-change map to the global-sections comparison]
  \label{lem:flat_base_change_reduce_global_sections}
  \lean{AlgebraicGeometry.flatBaseChange_isIso_iff_gammaTensorComparison}
  \uses{def:pushforward_base_change_map}
  % SOURCE: [Stacks Project], Tag 02KH (the reduction "part (1) follows from
  % part (2)") (read from references/stacks-coherent.md; quoted verbatim in
  % the SOURCE QUOTE PROOF of thm:flat_base_change_pushforward below).
  % SOURCE QUOTE: "since $R^if_*\mathcal{F}$ is quasi-coherent ..., it is the
  % quasi-coherent $\mathcal{O}_S$-module associated to the $A$-module
  % $H^0(S, R^if_*\mathcal{F}) = H^i(X, \mathcal{F})$ ... Since pullback by $g$
  % corresponds to $- \otimes_A B$ on modules ... we see that it suffices to
  % prove (2)."
  \textit{Source: Stacks Project, Tag 02KH (``Flat base change''), the reduction
  ``part (1) follows from part (2)''.}
  In the cartesian square with \(g\) flat and \(f\) quasi-compact and
  quasi-separated, the sheaf-level base-change map
  \(g^*(f_*\mathcal{F}) \to f'_*(g')^*\mathcal{F}\) of
  Definition~\ref{def:pushforward_base_change_map} is an isomorphism if and only if,
  after reducing to the affine bases \(S = \operatorname{Spec} A\),
  \(S' = \operatorname{Spec} B\), the module comparison map
  \(\Gamma(X,\mathcal{F}) \otimes_A B \to \Gamma(X_B, \mathcal{F}_B)\) is an
  isomorphism.
  \begin{proof}
    \uses{def:pushforward_base_change_map}
    Being an isomorphism is local on \(S'\), so we may take \(S = \operatorname{Spec} A\)
    and \(S' = \operatorname{Spec} B\) with \(A \to B\) flat. Because \(f\) is
    quasi-compact and quasi-separated, \(f_*\mathcal{F}\) is quasi-coherent on \(S\)
    (Stacks Tag 01XJ / Mathlib's quasi-coherence of the pushforward under a
    quasi-compact quasi-separated morphism), hence is the sheaf associated to the
    \(A\)-module \(\Gamma(S, f_*\mathcal{F}) = \Gamma(X,\mathcal{F})\); likewise
    \(f'_*\mathcal{F}'\) is associated to \(\Gamma(X_B,\mathcal{F}_B)\). Pullback along
    \(g\) corresponds to \(- \otimes_A B\) on modules, so under the
    tilde-equivalence the sheaf base-change map is the tilde of the comparison map
    \(\Gamma(X,\mathcal{F}) \otimes_A B \to \Gamma(X_B,\mathcal{F}_B)\); since
    \(\widetilde{(-)}\) is fully faithful on quasi-coherent modules, one is an
    isomorphism iff the other is.
  \end{proof}
\end{lemma}

\begin{theorem}\leanok
[Flat base change, $i=0$ case]
  \label{thm:flat_base_change_pushforward}
  \lean{AlgebraicGeometry.flatBaseChange_pushforward_isIso}
  \uses{def:pushforward_base_change_map,
        lem:flat_preserves_equalizer_mathlib,
        lem:flat_base_change_reduce_global_sections,
        lem:flat_base_change_mayer_vietoris}
  % SOURCE: [Stacks Project], Tag 02KH (Lemma "Flat base change", part for
  % $i = 0$) (read from references/stacks-coherent.tex, L947--970, and
  % references/stacks-coherent.md).
  % SOURCE QUOTE: "Consider a cartesian diagram of schemes
  % $$
  % \xymatrix{
  % X' \ar[d]_{f'} \ar[r]_{g'} & X \ar[d]^f \\
  % S' \ar[r]^g & S
  % }
  % $$
  % Let $\mathcal{F}$ be a quasi-coherent $\mathcal{O}_X$-module
  % with pullback $\mathcal{F}' = (g')^*\mathcal{F}$.
  % Assume that $g$ is flat and that $f$ is quasi-compact and quasi-separated.
  % For any $i \geq 0$
  % \begin{enumerate}
  % \item the base change map of
  % Cohomology, Lemma \ref{cohomology-lemma-base-change-map-flat-case}
  % is an isomorphism
  % $$
  % g^*R^if_*\mathcal{F} \longrightarrow R^if'_*\mathcal{F}',
  % $$
  % \item if $S = \Spec(A)$ and $S' = \Spec(B)$, then
  % $H^i(X, \mathcal{F}) \otimes_A B = H^i(X', \mathcal{F}')$.
  % \end{enumerate}"
  \textit{Source: Stacks Project, Tag 02KH (``Flat base change''), the
  \(i = 0\) case.}
  Consider the cartesian square
  \[
    \begin{array}{ccc}
      X' & \xrightarrow{\;g'\;} & X \\[2pt]
      {\scriptstyle f'}\big\downarrow & & \big\downarrow{\scriptstyle f} \\[2pt]
      S' & \xrightarrow{\;g\;} & S
    \end{array}
  \]
  with \(\mathcal{F}\) a quasi-coherent \(\mathcal{O}_X\)-module and
  \(\mathcal{F}' = (g')^*\mathcal{F}\). Assume that \(g\) is flat and that \(f\)
  is quasi-compact and quasi-separated. Then the base-change map of
  Definition~\ref{def:pushforward_base_change_map} is an isomorphism
  \[
    g^*\bigl(f_*\mathcal{F}\bigr) \;\xrightarrow{\;\sim\;}\;
    f'_*\bigl((g')^*\mathcal{F}\bigr) = f'_*\mathcal{F}' .
  \]
  Equivalently, in the affine situation \(S = \operatorname{Spec}(A)\),
  \(S' = \operatorname{Spec}(B)\) with \(A \to B\) flat, the comparison map
  \[
    H^0(X, \mathcal{F}) \otimes_A B \;\xrightarrow{\;\sim\;}\;
    H^0(X_B, \mathcal{F}_B)
  \]
  is an isomorphism, where \(X_B = X \times_{\operatorname{Spec}(A)}
  \operatorname{Spec}(B)\) and \(\mathcal{F}_B\) is the pullback of
  \(\mathcal{F}\).

  The hypothesis that \(\mathcal{F}\) is quasi-coherent is essential: it is what
  makes \(f_*\mathcal{F}\) quasi-coherent on the affine base and thus the tilde of
  a module, and it is what lets the affine pullback dictionary
  (Lemma~\ref{lem:pullback_spec_tilde_iso}) write \(\mathcal{F}\) as
  \(\widetilde{M}\) over each affine open. Without it the comparison map need not be
  an isomorphism.
\end{theorem}

% SOURCE QUOTE PROOF: "Having said this, part (1) follows from part (2). Namely,
% part (1) is local on $S'$ and hence we may assume $S$
% and $S'$ are affine. In other words, we have $S = \Spec(A)$
% and $S' = \Spec(B)$ as in (2).
% Then since $R^if_*\mathcal{F}$ is quasi-coherent
% (Lemma \ref{lemma-quasi-coherence-higher-direct-images}),
% it is the quasi-coherent $\mathcal{O}_S$-module associated to the
% $A$-module $H^0(S, R^if_*\mathcal{F}) = H^i(X, \mathcal{F})$
% (equality by
% Lemma \ref{lemma-quasi-coherence-higher-direct-images-application}).
% Similarly, $R^if'_*\mathcal{F}'$ is the quasi-coherent
% $\mathcal{O}_{S'}$-module associated to the $B$-module
% $H^i(X', \mathcal{F}')$. Since pullback by $g$ corresponds
% to $- \otimes_A B$ on modules
% (Schemes, Lemma \ref{schemes-lemma-widetilde-pullback})
% we see that it suffices to prove (2).
%
% Let $A \to B$ be a flat ring homomorphism.
% Let $X$ be a quasi-compact and quasi-separated scheme over $A$.
% Let $\mathcal{F}$ be a quasi-coherent $\mathcal{O}_X$-module.
% Set $X_B = X \times_{\Spec(A)} \Spec(B)$ and denote
% $\mathcal{F}_B$ the pullback of $\mathcal{F}$.
% We are trying to show that the map
% $$
% H^i(X, \mathcal{F}) \otimes_A B \longrightarrow H^i(X_B, \mathcal{F}_B)
% $$
% (given by the reference in the statement of the lemma)
% is an isomorphism.
%
% In case $X$ is separated, choose an affine open covering
% $\mathcal{U} : X = U_1 \cup \ldots \cup U_t$ and recall that
% $$
% \check{H}^p(\mathcal{U}, \mathcal{F}) = H^p(X, \mathcal{F}),
% $$
% see
% Lemma \ref{lemma-cech-cohomology-quasi-coherent}.
% If $\mathcal{U}_B : X_B = (U_1)_B \cup \ldots \cup (U_t)_B$ we obtain
% by base change, then it is still the case that each $(U_i)_B$ is affine
% and that $X_B$ is separated. Thus we obtain
% $$
% \check{H}^p(\mathcal{U}_B, \mathcal{F}_B) = H^p(X_B, \mathcal{F}_B).
% $$
% We have the following relation between the {\v C}ech complexes
% $$
% \check{\mathcal{C}}^\bullet(\mathcal{U}_B, \mathcal{F}_B) =
% \check{\mathcal{C}}^\bullet(\mathcal{U}, \mathcal{F}) \otimes_A B
% $$
% as follows from
% Lemma \ref{lemma-affine-base-change}.
% Since $A \to B$ is flat, the same thing remains true on taking cohomology.
%
% In case $X$ is quasi-separated, choose an affine open covering
% $\mathcal{U} : X = U_1 \cup \ldots \cup U_t$. We will use the
% {\v C}ech-to-cohomology spectral sequence
% Cohomology, Lemma \ref{cohomology-lemma-cech-spectral-sequence}.
% The reader who wishes to avoid this spectral sequence
% can use Mayer-Vietoris and induction on $t$ as in the proof of
% Lemma \ref{lemma-quasi-coherence-higher-direct-images}."
% NOTE: For the $i = 0$ case the argument is Čech-free: $H^0(X,\mathcal{F})$ is the
% sheaf-condition equalizer of a finite affine cover (Čech degree 0/1, NOT Čech
% cohomology -- no Leray, no spectral sequence), each AFFINE term base-changes by the
% concrete tilde dictionary lem:pullback_spec_tilde_iso (Tag 01I9), and flatness of
% $A \to B$ commutes $-\otimes_A B$ with that finite equalizer via
% lem:flat_preserves_equalizer_mathlib. The quasi-separated case is the Mayer--Vietoris
% induction on the number of cover members named in the source quote. The abstract
% affine lemma lem:affine_base_change_pushforward is not used in this proof.
\begin{proof}
  \uses{def:pushforward_base_change_map, lem:pullback_spec_tilde_iso,
        lem:flat_preserves_equalizer_mathlib,
        lem:flat_base_change_reduce_global_sections,
        lem:flat_base_change_mayer_vietoris}
  For \(i = 0\) the argument is {\v C}ech-\emph{free}: it uses only the sheaf
  condition (the degree-\(0\)/\(1\) part of {\v C}ech, i.e.\ an equalizer of a finite
  cover), the affine lemma, and the fact that flat base change commutes with that
  finite equalizer. There is no Leray sequence and no {\v C}ech-to-cohomology spectral
  sequence. The proof is the composite of the decomposition chain.

  \medskip
  \emph{Reduction to the global-sections comparison.} By
  Lemma~\ref{lem:flat_base_change_reduce_global_sections}, the base-change map of
  Definition~\ref{def:pushforward_base_change_map} is an isomorphism if and only if,
  after reducing to affine bases \(S = \operatorname{Spec}(A)\),
  \(S' = \operatorname{Spec}(B)\) with \(A \to B\) flat, the module comparison map
  \[
    \Gamma(X, \mathcal{F}) \otimes_A B \longrightarrow \Gamma(X_B, \mathcal{F}_B),
    \qquad X_B = X \times_{\operatorname{Spec}(A)} \operatorname{Spec}(B),
  \]
  is an isomorphism (this uses that \(f_*\mathcal{F}\) is quasi-coherent — \(f\) is
  quasi-compact and quasi-separated — so each side is the tilde of its module of
  global sections, and that pullback along \(g\) is \(- \otimes_A B\) on modules).

  \medskip
  \emph{The comparison map is an isomorphism.} Since \(f\) is quasi-compact and
  quasi-separated, \(X\) is quasi-compact and quasi-separated over \(A\), so
  Lemma~\ref{lem:flat_base_change_mayer_vietoris} applies and gives
  \[
    \Gamma(X, \mathcal{F}) \otimes_A B \;\xrightarrow{\ \sim\ }\; \Gamma(X_B, \mathcal{F}_B).
  \]
  Internally that lemma presents \(\Gamma(X,\mathcal{F})\) as the finite
  sheaf-condition equalizer of a finite affine cover
  (Lemma~\ref{lem:gamma_finite_equalizer}), base changes the equalizer diagram term by
  term via the concrete affine tilde dictionary Lemma~\ref{lem:pullback_spec_tilde_iso}
  (Lemma~\ref{lem:base_changed_equalizer_diagram}), and commutes \(- \otimes_A B\)
  past that finite equalizer by flatness
  (Lemma~\ref{lem:flat_preserves_equalizer_mathlib}), reducing the general
  quasi-separated case to the separated case
  (Lemma~\ref{lem:flat_base_change_separated}) by Mayer--Vietoris induction on the
  number of cover members. Combined with the reduction above, the sheaf-level
  base-change map is an isomorphism.
\end{proof}

\subsection{The \texorpdfstring{$A$}{A}-module \texorpdfstring{$\operatorname{eqLocus}$}{eqLocus} globalization (FBC-B)}

The separated-case argument is repackaged here at the level of \(A\)-modules, where
\(A = \Gamma(X,\mathcal{O}_X)\) is the ring of global sections of the structure sheaf, so that flat
base change becomes the Mathlib equivalence \(\operatorname{tensorEqLocusEquiv}\)
(Lemma~\ref{lem:flat_preserves_equalizer_mathlib}) applied to an honest
\(\operatorname{LinearMap.eqLocus}\). The blocks below construct that presentation concretely: every
module of sections \(\Gamma(U, M)\) is viewed as an \(A\)-module by restriction of scalars along the
structure-sheaf restriction \(A \to \Gamma(U,\mathcal{O}_X)\), the restriction maps of \(M\) become
\(A\)-linear, and the two product legs of the sheaf condition assemble into a pair of \(A\)-linear maps
whose equalizer locus is \(\Gamma(X, M)\). Throughout, \(M\) is a sheaf of
\(\mathcal{O}_X\)-modules and \(U : \iota \to \mathrm{Opens}(X)\) is an open cover with
\(\bigvee_i U_i = \top\).

The bijectivity of the keystone presentation rests on two element-level sheaf axioms of the underlying
abelian-group sheaf, recorded as Mathlib dependency anchors.

\begin{lemma}[Sheaf separatedness: local equality forces global equality]
  \label{lem:sheaf_eq_of_locally_eq_mathlib}
  \lean{TopCat.Sheaf.eq_of_locally_eq'}
  \mathlibok
  \textit{Provided by Mathlib.}
  Let \(\mathcal{G}\) be a sheaf on a topological space \(X\), \(U : \iota \to \mathrm{Opens}(X)\) a
  family covering an open \(V\) (i.e.\ \(V \le \bigvee_i U_i\)), and \(s, t \in \mathcal{G}(V)\) two
  sections. If the restrictions of \(s\) and \(t\) to every \(U_i\) agree, then \(s = t\). This is the
  separatedness half of the sheaf condition, stated element-wise.
\end{lemma}

\begin{lemma}[Sheaf gluing: a compatible family has a unique amalgamation]
  \label{lem:sheaf_existsUnique_gluing_mathlib}
  \lean{TopCat.Sheaf.existsUnique_gluing'}
  \mathlibok
  \textit{Provided by Mathlib.}
  Let \(\mathcal{G}\) be a sheaf on \(X\), \(U : \iota \to \mathrm{Opens}(X)\) a family covering an open
  \(V\), and \((s_i)_i\) a family of sections \(s_i \in \mathcal{G}(U_i)\) that is \emph{compatible}
  (the restrictions of \(s_i\) and \(s_j\) to \(U_i \cap U_j\) agree for all \(i, j\)). Then there is a
  unique \(s \in \mathcal{G}(V)\) whose restriction to each \(U_i\) is \(s_i\). This is the gluing half
  of the sheaf condition, stated element-wise.
\end{lemma}

\begin{definition}\leanok
[The ground ring of global sections]
  \label{def:fbcb_groundRing}
  \lean{AlgebraicGeometry.groundRing}
  For a scheme \(X\), the \emph{ground ring} is the commutative ring
  \(A = \Gamma(X, \mathcal{O}_X) = \mathcal{O}_X(\top)\) of global sections of the structure sheaf,
  taken from the \(\mathbf{CommRing}\)-valued structure presheaf so its ring and algebra structures
  resolve canonically.
\end{definition}

\begin{definition}\leanok
[The structure-sheaf restriction ring map \texorpdfstring{$A \to \Gamma(U,\mathcal{O}_X)$}{A to sections over U}]
  \label{def:fbcb_rhoU}
  \lean{AlgebraicGeometry.rhoU}
  \uses{def:fbcb_groundRing}
  For an open \(U \subseteq X\), the structure-sheaf restriction along \(U \subseteq \top\) is a ring
  homomorphism \(\rho_U : A = \Gamma(X,\mathcal{O}_X) \to \Gamma(U,\mathcal{O}_X)\). It is the ring map
  along which sections over \(U\) are regarded as \(A\)-modules.
\end{definition}

\begin{lemma}\leanok
[Transitivity of the restriction ring maps]
  \label{lem:fbcb_rhoU_comp}
  \lean{AlgebraicGeometry.rhoU_comp}
  \uses{def:fbcb_rhoU}
  For opens \(V \le U \subseteq X\), the structure-sheaf restriction
  \(\Gamma(U,\mathcal{O}_X) \to \Gamma(V,\mathcal{O}_X)\) composed with \(\rho_U\) equals \(\rho_V\):
  \[
    \bigl(\Gamma(U,\mathcal{O}_X) \to \Gamma(V,\mathcal{O}_X)\bigr) \circ \rho_U = \rho_V .
  \]
  \begin{proof}
    \uses{def:fbcb_rhoU}
    Both sides are the structure-sheaf restriction \(\Gamma(X,\mathcal{O}_X) \to \Gamma(V,\mathcal{O}_X)\)
    along \(V \subseteq \top\), factored as \(V \subseteq U \subseteq \top\); functoriality of the
    structure presheaf collapses the two-step restriction to the one-step restriction.
  \end{proof}
\end{lemma}

\begin{definition}\leanok
[Sections over \texorpdfstring{$U$}{U} as an \texorpdfstring{$A$}{A}-module]
  \label{def:fbcb_gammaModA}
  \lean{AlgebraicGeometry.gammaModA}
  \uses{def:fbcb_rhoU}
  The sections \(\Gamma(U, M)\), a module over \(\Gamma(U,\mathcal{O}_X)\), are regarded as an
  \(A\)-module by restriction of scalars along \(\rho_U\) (Definition~\ref{def:fbcb_rhoU}). This is the
  common \(A\)-module home in which the global-sections equalizer is formed.
\end{definition}

\begin{definition}\leanok
[The restriction \texorpdfstring{$\Gamma(U,M) \to \Gamma(V,M)$}{Gamma(U,M) to Gamma(V,M)} as an \texorpdfstring{$A$}{A}-module morphism]
  \label{def:fbcb_gammaResAHom}
  \lean{AlgebraicGeometry.gammaResAHom}
  \uses{def:fbcb_gammaModA, lem:fbcb_rhoU_comp}
  For \(V \le U\), the structure-sheaf restriction \(M.\mathrm{map} : \Gamma(U,M) \to \Gamma(V,M)\) is
  promoted to a morphism of \(A\)-modules \(\Gamma(U,M) \to \Gamma(V,M)\) (Definition~\ref{def:fbcb_gammaModA}).
  It is built from \(M.\mathrm{map}\) by restriction of scalars, using the transitivity
  \(\rho\)-coherence Lemma~\ref{lem:fbcb_rhoU_comp} to identify the two scalar actions.
\end{definition}

\begin{definition}\leanok
[The \texorpdfstring{$A$}{A}-linear restriction map]
  \label{def:fbcb_gammaResA}
  \lean{AlgebraicGeometry.gammaResA}
  \uses{def:fbcb_gammaResAHom}
  The underlying \(A\)-linear map \(\Gamma(U,M) \to_A \Gamma(V,M)\) of the \(A\)-module morphism of
  Definition~\ref{def:fbcb_gammaResAHom}. It is the building block of the two product legs of the
  equalizer.
\end{definition}

\begin{lemma}\leanok
[The \texorpdfstring{$A$}{A}-linear restriction is the underlying presheaf restriction]
  \label{lem:fbcb_gammaResA_apply}
  \lean{AlgebraicGeometry.gammaResA_apply}
  \uses{def:fbcb_gammaResA}
  On underlying elements, the \(A\)-linear restriction map of Definition~\ref{def:fbcb_gammaResA} acts
  as the presheaf restriction of \(M\): for \(x \in \Gamma(U,M)\) and \(V \le U\),
  \(\mathrm{gammaResA}(x) = (M.\mathrm{map})\,x\).
  \begin{proof}
    \uses{def:fbcb_gammaResA}
    Restriction of scalars and the composition with the comparison isomorphism of
    Definition~\ref{def:fbcb_gammaResAHom} leave the underlying function unchanged, so the map agrees
    with \(M.\mathrm{map}\) on elements.
  \end{proof}
\end{lemma}

\begin{lemma}\leanok
[Functoriality of the \texorpdfstring{$A$}{A}-linear restriction maps]
  \label{lem:fbcb_gammaResA_comp}
  \lean{AlgebraicGeometry.gammaResA_comp}
  \uses{def:fbcb_gammaResA, lem:fbcb_gammaResA_apply}
  For opens \(W \le V \le U\) and \(x \in \Gamma(U,M)\), the \(A\)-linear restrictions compose as the
  single restriction:
  \[
    \mathrm{gammaResA}_{V \le U}\!\bigl(x\bigr) \text{ then } \mathrm{gammaResA}_{W \le V}
      \;=\; \mathrm{gammaResA}_{W \le U}\!\bigl(x\bigr).
  \]
  \begin{proof}
    \uses{def:fbcb_gammaResA, lem:fbcb_gammaResA_apply}
    By Lemma~\ref{lem:fbcb_gammaResA_apply} each \(A\)-linear restriction is the underlying presheaf
    restriction of \(M\); functoriality of the presheaf \(M\) collapses the composite of the two
    restrictions \(W \subseteq V \subseteq U\) to the single restriction \(W \subseteq U\).
  \end{proof}
\end{lemma}

\begin{definition}\leanok
[The left restriction leg]
  \label{def:fbcb_leftRes}
  \lean{AlgebraicGeometry.leftRes}
  \uses{def:fbcb_gammaResA}
  The \(A\)-linear map
  \[
    \mathrm{leftRes} : \prod_i \Gamma(U_i, M) \longrightarrow \prod_{(i,j)} \Gamma(U_i \cap U_j, M),
  \]
  whose \((i,j)\)-component is the \(A\)-linear restriction \(\Gamma(U_i, M) \to \Gamma(U_i \cap U_j, M)\)
  of Definition~\ref{def:fbcb_gammaResA} applied to the \(i\)-th factor (restriction along
  \(U_i \cap U_j \subseteq U_i\)).
\end{definition}

\begin{definition}\leanok
[The right restriction leg]
  \label{def:fbcb_rightRes}
  \lean{AlgebraicGeometry.rightRes}
  \uses{def:fbcb_gammaResA}
  The \(A\)-linear map \(\mathrm{rightRes} : \prod_i \Gamma(U_i, M) \to \prod_{(i,j)} \Gamma(U_i \cap U_j, M)\)
  whose \((i,j)\)-component is the \(A\)-linear restriction of Definition~\ref{def:fbcb_gammaResA} applied
  to the \(j\)-th factor (restriction along \(U_i \cap U_j \subseteq U_j\)).
\end{definition}

\begin{definition}\leanok
[Restriction of a global section to the cover]
  \label{def:fbcb_toCover}
  \lean{AlgebraicGeometry.toCover}
  \uses{def:fbcb_gammaResA}
  The \(A\)-linear map \(\mathrm{toCover} : \Gamma(\top, M) \to \prod_i \Gamma(U_i, M)\) restricting a
  global section to each member of the cover, with \(i\)-th component the \(A\)-linear restriction
  \(\Gamma(\top, M) \to \Gamma(U_i, M)\) of Definition~\ref{def:fbcb_gammaResA}.
\end{definition}

\begin{lemma}\leanok
[The restricted family is compatible]
  \label{lem:fbcb_leftRes_toCover}
  \lean{AlgebraicGeometry.leftRes_toCover}
  \uses{def:fbcb_leftRes, def:fbcb_rightRes, def:fbcb_toCover, lem:fbcb_gammaResA_comp}
  For every global section \(s \in \Gamma(\top, M)\) the two legs agree on the restricted family:
  \[
    \mathrm{leftRes}\bigl(\mathrm{toCover}(s)\bigr) = \mathrm{rightRes}\bigl(\mathrm{toCover}(s)\bigr).
  \]
  \begin{proof}
    \uses{def:fbcb_leftRes, def:fbcb_rightRes, def:fbcb_toCover, lem:fbcb_gammaResA_comp}
    Both sides are, in the \((i,j)\)-component, the restriction of \(s\) from \(\top\) to
    \(U_i \cap U_j\): the left leg restricts \(\top \to U_i \to U_i \cap U_j\) and the right leg
    restricts \(\top \to U_j \to U_i \cap U_j\), and by the functoriality
    Lemma~\ref{lem:fbcb_gammaResA_comp} each two-step restriction collapses to the one-step restriction
    \(\top \to U_i \cap U_j\). Hence the two components coincide.
  \end{proof}
\end{lemma}

\begin{definition}\leanok
[The corestriction into the equalizer locus]
  \label{def:fbcb_toCoverEqLocus}
  \lean{AlgebraicGeometry.toCoverEqLocus}
  \uses{def:fbcb_toCover, lem:fbcb_leftRes_toCover}
  By the compatibility Lemma~\ref{lem:fbcb_leftRes_toCover}, the map \(\mathrm{toCover}\) corestricts to
  an \(A\)-linear map
  \[
    \mathrm{toCoverEqLocus} : \Gamma(\top, M)
      \longrightarrow \operatorname{eqLocus}(\mathrm{leftRes}, \mathrm{rightRes})
      = \{\, t \in \textstyle\prod_i \Gamma(U_i,M) : \mathrm{leftRes}(t) = \mathrm{rightRes}(t)\,\}.
  \]
\end{definition}

\begin{lemma}\leanok
[Global sections as an \texorpdfstring{$A$}{A}-module equalizer locus]
  \label{lem:fbcb_gammaTopEquivEqLocus}
  \lean{AlgebraicGeometry.gammaTopEquivEqLocus}
  \uses{def:fbcb_toCoverEqLocus, def:fbcb_gammaModA, def:fbcb_gammaResA, def:fbcb_leftRes,
        def:fbcb_rightRes, lem:sheaf_eq_of_locally_eq_mathlib, lem:sheaf_existsUnique_gluing_mathlib}
  % NOTE: bijectivity is established by the element-level sheaf axioms
  % (lem:sheaf_eq_of_locally_eq_mathlib for injectivity, lem:sheaf_existsUnique_gluing_mathlib for
  % surjectivity) on the underlying abelian-group sheaf; the route is direct and does not
  % pass through the sheaf-condition-fork limit (lem:sheaf_equalizer_products_mathlib).
  For \(M\) a sheaf of \(\mathcal{O}_X\)-modules and any open cover \(U\) of \(X\) (\(\bigvee_i U_i = \top\)), the
  corestriction \(\mathrm{toCoverEqLocus}\) of Definition~\ref{def:fbcb_toCoverEqLocus} is an
  \(A\)-linear isomorphism
  \[
    \Gamma(X, M) = \Gamma(\top, M) \;\cong\; \operatorname{eqLocus}(\mathrm{leftRes}, \mathrm{rightRes}),
  \]
  identifying the global sections, as a module over the ground ring \(A = \Gamma(X,\mathcal{O}_X)\), with
  the submodule on which the two restriction legs agree.
  \begin{proof}
    \uses{def:fbcb_toCoverEqLocus, lem:sheaf_eq_of_locally_eq_mathlib,
          lem:sheaf_existsUnique_gluing_mathlib}
    The map \(\mathrm{toCoverEqLocus}\) is \(A\)-linear by construction, so it suffices to prove it
    bijective on underlying sets. \emph{Injectivity:} if two global sections \(s, t\) have equal images,
    their restrictions to every \(U_i\) agree, so by sheaf separatedness
    (Lemma~\ref{lem:sheaf_eq_of_locally_eq_mathlib}) applied to the underlying abelian-group sheaf
    the underlying abelian-group sheaf, \(s = t\). \emph{Surjectivity:} an element of the locus is a family
    \((s_i)_i\) with \(\mathrm{leftRes} = \mathrm{rightRes}\), i.e.\ a compatible family on the cover; by
    the gluing axiom (Lemma~\ref{lem:sheaf_existsUnique_gluing_mathlib}) it amalgamates to a global
    section restricting to \((s_i)_i\). Hence \(\mathrm{toCoverEqLocus}\) is an \(A\)-linear bijection,
    so an isomorphism.
  \end{proof}
\end{lemma}

\begin{lemma}\leanok
[Flat base change commutes with the \texorpdfstring{$H^0$}{H0} equalizer]
  \label{lem:fbcb_baseChangeGammaEquiv}
  \lean{AlgebraicGeometry.baseChangeGammaEquiv}
  \uses{lem:fbcb_gammaTopEquivEqLocus, lem:flat_preserves_equalizer_mathlib}
  Let \(B\) be a flat \(A\)-algebra (\(A = \Gamma(X,\mathcal{O}_X)\)). Then base changing the global
  sections along \(A \to B\) is the equalizer of the base-changed restriction legs: there is a
  \(B\)-linear isomorphism
  \[
    B \otimes_A \Gamma(X, M) \;\cong\;
      \operatorname{eqLocus}\bigl(\mathrm{id}_B \otimes \mathrm{leftRes},\ \mathrm{id}_B \otimes \mathrm{rightRes}\bigr).
  \]
  \begin{proof}
    \uses{lem:fbcb_gammaTopEquivEqLocus, lem:flat_preserves_equalizer_mathlib}
    Tensoring the keystone isomorphism Lemma~\ref{lem:fbcb_gammaTopEquivEqLocus} by \(B\) over \(A\)
    gives \(B \otimes_A \Gamma(X,M) \cong B \otimes_A \operatorname{eqLocus}(\mathrm{leftRes},\mathrm{rightRes})\),
    and by flatness of \(B\) over \(A\) the Mathlib equivalence \(\operatorname{tensorEqLocusEquiv}\)
    (Lemma~\ref{lem:flat_preserves_equalizer_mathlib}) carries \(B \otimes_A (-)\) of the equalizer to the
    equalizer of the base-changed legs. Composing the two yields the stated \(B\)-linear isomorphism.
    By the same presentation applied to the pulled-back sheaf, the right-hand side is the global sections
    of \(M\) over the base-changed scheme, so this is the module-level core of the \(H^0\)
    flat-base-change comparison.
  \end{proof}
\end{lemma}

\section{FBC-B: the global \texorpdfstring{$H^0$}{H0} flat-base-change isomorphism via the finite affine-cover equalizer (direct route)}

The blocks of the preceding subsection assemble into a \emph{direct} proof of the named global
isomorphism that bypasses the adjoint-mate keystone of the affine/mate route. The route is
\emph{{\v C}ech-cohomology-free}: \(H^0\) is a finite limit — the equalizer of a finite affine cover —
and the flat functor \(-\otimes_A B\) preserves finite limits, so the global comparison isomorphism
follows from per-chart affine base change \emph{without} the abstract conjugate/mate identification and
without the Mayer--Vietoris induction.

Concretely, the module-level core
\(B \otimes_A \Gamma(X,\mathcal{F}) \cong \operatorname{eqLocus}(\mathrm{id}_B\otimes\mathrm{leftRes},
\mathrm{id}_B\otimes\mathrm{rightRes})\) (Lemma~\ref{lem:fbcb_baseChangeGammaEquiv}) already commutes flat
base change past the finite \(H^0\) equalizer of a cover \(\{U_i\}\) of \(X\). The single remaining step
is to recognise the right-hand base-changed equalizer, chart by chart, as the equalizer-locus
presentation (Lemma~\ref{lem:fbcb_gammaTopEquivEqLocus}) of the global sections of the pulled-back module
over the base-changed cover \(\{(U_i)_B\}\) of \(X' = X_B\). The theorem below states that assembly.

\begin{definition}\leanok
[Ground-ring algebra map of the base-change pullback]
  \label{def:fbcb_pullbackGroundRingAlg}
  \lean{AlgebraicGeometry.Modules.pullbackGroundRingAlg}
  For a flat \(A\)-algebra \(B\) (\(A = \Gamma(X,\OO_X)\)) and the base-change pullback
  \(X' = X \times_{\operatorname{Spec} A} \operatorname{Spec} B\), the projection
  \(f' : X' \to \operatorname{Spec} B\) induces a ring homomorphism
  \(B \to \Gamma(X',\OO_{X'})\), obtained as
  \(f'^{\sharp}_{\top} \circ (\Gamma\!\operatorname{Spec}\text{-iso})^{-1}\). This is the
  map along which \(\Gamma(X',\mathcal{F}')\) is viewed as a \(B\)-module (restriction of
  scalars) in Theorem~\ref{thm:fbcb_global_direct}; chasing \(\Gamma(X',\OO_{X'}) = B\)
  directly is avoided (that equality is the theorem at \(\mathcal{F} = \OO_X\)).
\end{definition}

\begin{lemma}



\leanok
[Per-chart identification of the base-changed equalizer locus]
  \label{lem:fbcb_baseChangeEqLocusToPullbackGamma}
  \lean{AlgebraicGeometry.Modules.baseChangeEqLocusToPullbackGamma}
  \uses{lem:fbcb_gammaTopEquivEqLocus, lem:pullback_spec_tilde_iso,
        def:fbcb_pullbackGroundRingAlg, lem:pullback_tilde_restriction_natural,
        lem:base_change_chart_tensor_iso, lem:tensor_finite_product_comm}
  Let \(X\) be a quasi-compact, quasi-separated (for the affine-overlap step, separated)
  scheme with ground ring \(A = \Gamma(X,\OO_X)\), \(B\) a flat \(A\)-algebra,
  \(X' = X_B\), \(\mathcal{F}' = (g')^*\mathcal{F}\), and \(\{U_i\}_{i}\) the finite affine
  cover of Lemma~\ref{lem:finite_affine_cover_qcqs} with base change
  \(\{U_i' = (g')^{-1}(U_i)\}\). Then there is a \(B\)-linear isomorphism
  \[
    \operatorname{eqLocus}\bigl(\mathrm{id}_B\otimes\mathrm{leftRes},\
      \mathrm{id}_B\otimes\mathrm{rightRes}\bigr)
    \;\cong\;
    (\operatorname{restr}_{\,\mathtt{pullbackGroundRingAlg}\,B})\,
      \bigl(\Gamma(X', \mathcal{F}')\bigr),
  \]
  where the codomain is \(\Gamma(X',\mathcal{F}')\) viewed as a \(B\)-module by restriction
  of scalars along \(B \to \Gamma(X',\OO_{X'})\) (Definition~\ref{def:fbcb_pullbackGroundRingAlg}).
  \begin{proof}
    \uses{lem:fbcb_gammaTopEquivEqLocus, lem:pullback_spec_tilde_iso,
          def:fbcb_pullbackGroundRingAlg, lem:pullback_tilde_restriction_natural,
          lem:base_change_chart_tensor_iso, lem:tensor_finite_product_comm}
    The base-changed cover \(\{U_i'\}\) is a finite affine cover of \(X'\) with
    \(\bigvee_i U_i' = \top\) (base change preserves affineness and open immersions and
    carries \(\bigvee_i U_i = \top\) across). Three pieces compose:
    \emph{(a) per-chart and per-overlap identification.} Each \(U_i\) and (separated case)
    each \(U_i\cap U_j\) is affine, so the concrete pullback dictionary
    Lemma~\ref{lem:pullback_spec_tilde_iso} (Tag 01I9) reads
    \(\Gamma(U_i',\mathcal{F}') \cong \Gamma(U_i,\mathcal{F})\otimes_A B\) (and likewise on
    overlaps) naturally in restriction, intertwining the base-changed legs
    \(\mathrm{id}_B\otimes\mathrm{leftRes}\), \(\mathrm{id}_B\otimes\mathrm{rightRes}\) with
    the restriction legs \(\mathrm{leftRes}'\), \(\mathrm{rightRes}'\) of \(\{U_i'\}\); this
    gives an isomorphism of equalizer loci
    \(\operatorname{eqLocus}(\mathrm{id}_B\otimes\mathrm{leftRes},\mathrm{id}_B\otimes\mathrm{rightRes})
    \cong \operatorname{eqLocus}(\mathrm{leftRes}',\mathrm{rightRes}')\).
    Finiteness of the index set (so that \(B\otimes_A\prod_i \cong \prod_i (B\otimes_A -)\),
    Mathlib's tensor-preserves-finite-products) is what makes \(\mathrm{id}_B\otimes(-)\) of
    the legs land in the finite products of the base-changed cover.
    \emph{(b) recognising the right-hand locus.} The equalizer-locus presentation
    Lemma~\ref{lem:fbcb_gammaTopEquivEqLocus} for \(\mathcal{F}'\) and \(\{U_i'\}\) identifies
    \(\operatorname{eqLocus}(\mathrm{leftRes}',\mathrm{rightRes}') \cong \Gamma(X',\mathcal{F}')\)
    as \(\Gamma(X',\OO_{X'})\)-modules.
    \emph{(c) restriction of scalars.} Transporting (b) along \(B \to \Gamma(X',\OO_{X'})\)
    (Definition~\ref{def:fbcb_pullbackGroundRingAlg}) yields the \(B\)-linear codomain as
    stated; chasing \(\Gamma(X',\OO_{X'}) = B\) is avoided. Composing (a)–(c) gives the
    isomorphism. No abstract base-change map and no mate identification enter, every piece
    being affine.
  \end{proof}
\end{lemma}

\begin{theorem}\leanok
[Global \(H^0\) flat base change via the finite-cover equalizer, direct route]
  \label{thm:fbcb_global_direct}
  \lean{AlgebraicGeometry.Modules.baseChangeGammaPullbackEquiv}
  \uses{lem:finite_affine_cover_qcqs, lem:fbcb_gammaTopEquivEqLocus,
        lem:fbcb_baseChangeGammaEquiv, lem:pullback_spec_tilde_iso,
        lem:fbcb_baseChangeEqLocusToPullbackGamma, lem:flat_base_change_reduce_global_sections,
        def:fbcb_pullbackGroundRingAlg, lem:pullback_tilde_restriction_natural,
        lem:base_change_chart_tensor_iso, lem:tensor_finite_product_comm}
  % SOURCE: [Stacks Project], Tag 02KH (``Flat base change''), part (2), $i = 0$ case
  % (read from references/stacks-coherent.tex, L966--968).
  % SOURCE QUOTE: "\item if $S = \Spec(A)$ and $S' = \Spec(B)$, then
  % $H^i(X, \mathcal{F}) \otimes_A B = H^i(X', \mathcal{F}')$."
  \textit{Source: Stacks Project, Tag 02KH (``Flat base change''), part (2), \(i = 0\) case.}
  Let \(X\) be a quasi-compact, quasi-separated scheme with ground ring
  \(A = \Gamma(X,\mathcal{O}_X)\), let \(\mathcal{F}\) be a quasi-coherent \(\mathcal{O}_X\)-module, and
  let \(B\) be a flat \(A\)-algebra. Write \(X' = X_B = X \times_{\operatorname{Spec}(A)}
  \operatorname{Spec}(B)\) with projection \(g' : X' \to X\) and \(\mathcal{F}' = (g')^*\mathcal{F}\) the
  pullback. Then the comparison map
  \[
    \Gamma(X,\mathcal{F}) \otimes_A B \;\xrightarrow{\;\sim\;}\; \Gamma(X', \mathcal{F}')
  \]
  is a \(B\)-linear isomorphism, assembled directly from the finite affine-cover equalizer presentation;
  consequently the sheaf-level base-change map \(g^*(f_*\mathcal{F}) \to f'_*\mathcal{F}'\) is an
  isomorphism, recovering the named global target Theorem~\ref{thm:flat_base_change_pushforward}. The
  route uses neither the adjoint-mate keystone of the conjugate route nor the Mayer--Vietoris induction
  on the number of cover members.

  % SOURCE QUOTE PROOF: "\item We have $\psi^* \widetilde M = \widetilde{S \otimes_R M}$
  % functorially in the $R$-module $M$." (references/stacks-schemes.tex, Tag 01I9, part (1) —
  % the per-chart pullback-is-base-change identification feeding the assembly.)
  \begin{proof}
    \uses{lem:finite_affine_cover_qcqs, lem:fbcb_gammaTopEquivEqLocus,
          lem:fbcb_baseChangeGammaEquiv, lem:pullback_spec_tilde_iso,
          lem:fbcb_baseChangeEqLocusToPullbackGamma, lem:flat_base_change_reduce_global_sections,
          lem:pullback_tilde_restriction_natural, lem:base_change_chart_tensor_iso,
          lem:tensor_finite_product_comm}
    Choose a finite affine cover \(\{U_i\}_{i\in\iota}\) of \(X\) with quasi-compact pairwise overlaps
    (Lemma~\ref{lem:finite_affine_cover_qcqs}), so that \(\bigvee_i U_i = \top\).

    \emph{Flat base change past the finite equalizer.} By
    Lemma~\ref{lem:fbcb_baseChangeGammaEquiv}, flatness of \(B\) over \(A\) gives a \(B\)-linear
    isomorphism
    \[
      B \otimes_A \Gamma(X,\mathcal{F}) \;\cong\;
        \operatorname{eqLocus}\bigl(\mathrm{id}_B\otimes\mathrm{leftRes},\ \mathrm{id}_B\otimes\mathrm{rightRes}\bigr),
    \]
    where \(\mathrm{leftRes}, \mathrm{rightRes} : \prod_i \Gamma(U_i,\mathcal{F}) \rightrightarrows
    \prod_{i,j}\Gamma(U_i\cap U_j,\mathcal{F})\) are the two restriction legs of the cover.

    \emph{Per-chart identification of the base-changed legs.} The base-changed cover
    \(\{(U_i)_B\}_{i\in\iota}\) is again a finite affine cover, now of \(X'\): base change preserves
    affineness and open immersions and carries \(\bigvee_i U_i = \top\) to \(\bigvee_i (U_i)_B = \top\).
    On each chart the affine pullback dictionary (Lemma~\ref{lem:pullback_spec_tilde_iso}, Stacks
    Tag 01I9: \((\operatorname{Spec}\varphi)^*\widetilde M \cong \widetilde{B\otimes_A M}\)) reads the
    sections of \(\mathcal{F}'\) over \((U_i)_B\) as the base change of the sections of \(\mathcal{F}\),
    \[
      \Gamma\bigl((U_i)_B, \mathcal{F}'\bigr) \;\cong\; \Gamma(U_i,\mathcal{F}) \otimes_A B,
    \]
    and likewise on each overlap: in the separated case the overlaps \(U_i\cap U_j\) are
    again affine, so the \emph{same} concrete dictionary
    Lemma~\ref{lem:pullback_spec_tilde_iso} supplies
    \(\Gamma((U_i\cap U_j)_B, \mathcal{F}') \cong \Gamma(U_i\cap U_j,\mathcal{F})\otimes_A B\)
    — no abstract base-change map and no mate identification is needed. These
    identifications are natural in restriction, so they intertwine the base-changed legs
    \(\mathrm{id}_B\otimes\mathrm{leftRes}\), \(\mathrm{id}_B\otimes\mathrm{rightRes}\) with the
    restriction legs \(\mathrm{leftRes}'\), \(\mathrm{rightRes}'\) of the cover \(\{(U_i)_B\}\) of
    \(X'\); hence they restrict to an isomorphism of equalizer loci
    \[
      \operatorname{eqLocus}\bigl(\mathrm{id}_B\otimes\mathrm{leftRes},\ \mathrm{id}_B\otimes\mathrm{rightRes}\bigr)
      \;\cong\;
      \operatorname{eqLocus}\bigl(\mathrm{leftRes}',\ \mathrm{rightRes}'\bigr).
    \]

    \emph{Recognising the right-hand equalizer.} Applying the equalizer-locus presentation
    (Lemma~\ref{lem:fbcb_gammaTopEquivEqLocus}) to \(\mathcal{F}'\) and the cover \(\{(U_i)_B\}\) of
    \(X'\) identifies the latter equalizer locus with the global sections,
    \(\operatorname{eqLocus}(\mathrm{leftRes}',\mathrm{rightRes}') \cong \Gamma(X',\mathcal{F}')\).
    Composing the three isomorphisms yields the \(B\)-linear isomorphism
    \(\Gamma(X,\mathcal{F})\otimes_A B \cong \Gamma(X',\mathcal{F}')\), which is the comparison map.

    \emph{Back to the sheaf-level statement.} Finally, the reduction
    Lemma~\ref{lem:flat_base_change_reduce_global_sections} — each pushforward is the tilde of its module
    of global sections, and pullback along \(g\) is \(-\otimes_A B\) on modules — turns this module
    isomorphism into the isomorphism of the sheaf-level base-change map
    \(g^*(f_*\mathcal{F}) \to f'_*\mathcal{F}'\), the named global target
    Theorem~\ref{thm:flat_base_change_pushforward}. No conjugate/mate identification and no {\v C}ech
    cohomology beyond the degree \(0/1\) sheaf condition enters.
  \end{proof}
\end{theorem}
